% ===========================================================================
%  KTH -- Team Reference Document  (NWERC)
%  GENERATED FILE: edit src/*.tex, then run tools/build.sh
%  Paste notebook.tex into Overleaf and hit Recompile -- plain pdfLaTeX.
%
%  Layout matches KACTL exactly: A4 LANDSCAPE, 3 columns, 8pt code.
%  (kactl uses \documentclass[9pt,a4paper]{extreport} + geometry
%   a4paper,landscape,hmargin={0.5cm,0.5cm},vmargin={1.0cm,0.4cm}
%   + multicols*{3} + \footnotesize listings, i.e. 8pt code.)
%
%  NWERC rules encoded here: <= 25 pages, single-sided, A4, university name
%  in the UPPER LEFT, page number in the UPPER RIGHT.
% ===========================================================================
\documentclass[9pt,a4paper]{extarticle}

\newcommand{\uni}{KTH Royal Institute of Technology}
\newcommand{\unishort}{KTH}
\newcommand{\team}{Wery Accepted}
\newcommand{\members}{Ruben \quad Ruiming \quad Teo}
\newcommand{\contest}{NWERC 2026}

% Code size. 8pt is exactly KACTL's, which is printed for NWERC every year.
% DO NOT go below it -- raise both numbers if the print looks tight.
\newcommand{\codesize}{\fontsize{8}{8.8}\selectfont}

% includehead makes geometry reserve the header INSIDE the top margin. Without it
% LaTeX puts the header box top at 1in+\topmargin, which lands ABOVE the paper edge
% here and slices the header in half. (KACTL papers over the same problem by setting
% \headheight to 25pt; reserving the space properly is the honest fix.)
\usepackage[a4paper,landscape,hmargin={0.5cm,0.5cm},vmargin={0.4cm,0.4cm},
            includehead,headsep=4pt]{geometry}
\usepackage{multicol}
\usepackage{lmodern}
\usepackage[T1]{fontenc}
\usepackage{listings}
\usepackage{xcolor}
\usepackage{graphicx}
\usepackage{titlesec}
\usepackage{fancyhdr}
\usepackage{amsmath,amssymb}
\usepackage{tocloft}
\usepackage[expansion=false]{microtype}

\definecolor{secbg}{HTML}{1F3B57}
\definecolor{hdbg}{HTML}{E8ECF1}
\definecolor{kw}{HTML}{16507B}
\definecolor{cmt}{HTML}{6B7C6B}
\definecolor{str}{HTML}{8B3A2E}
\definecolor{rule}{HTML}{9AA7B4}

\setlength{\parindent}{0pt}
\setlength{\parskip}{0pt}
\setlength{\columnsep}{4mm}
\setlength{\columnseprule}{0.2pt}

\newcommand{\secbar}[1]{%
  \colorbox{secbg}{\parbox{\dimexpr\columnwidth-2\fboxsep}{%
    \strut\color{white}\bfseries #1}}}
\titleformat{\section}{\normalfont}{}{0pt}{\secbar}
\titlespacing{\section}{0pt}{1.1ex}{0.6ex}

\pagestyle{fancy}
\fancyhf{}
\renewcommand{\headrulewidth}{0.2pt}
\lhead{\footnotesize\uni}                 % NWERC: university, upper left
\chead{\footnotesize\team}
\rhead{\footnotesize\thepage}             % NWERC: page number, upper right
% headsep is set by geometry above; do not override it here or the header
% drifts back off the top of the page.

\lstset{
  language=C++,
  basicstyle=\ttfamily\codesize,
  keywordstyle=\color{kw}\bfseries,
  commentstyle=\color{cmt},
  stringstyle=\color{str},
  showstringspaces=false,
  keepspaces=true,
  columns=fullflexible,
  breaklines=false,
  tabsize=4,
  aboveskip=0.8ex,
  belowskip=0.3ex,
  xleftmargin=0pt,
  morekeywords={ll,ld,vi,pii,matrix,rep,down,all,i128,ull}
}

% Compact index: every \section and every \entry gets a table-of-contents line.
% secnumdepth 0 = unnumbered sections. The section bars show no number either,
% and with numbering on, \numberline would overprint the title in the index.
\setcounter{secnumdepth}{0}
\setcounter{tocdepth}{2}
\setlength{\cftbeforesecskip}{0.35ex}
\setlength{\cftbeforesubsecskip}{0pt}
\setlength{\cftsecindent}{0pt}
\setlength{\cftsubsecindent}{0.9em}
\setlength{\cftsecnumwidth}{0em}
\setlength{\cftsubsecnumwidth}{0em}
\renewcommand{\cftsecfont}{\bfseries\footnotesize}
\renewcommand{\cftsecpagefont}{\bfseries\footnotesize}
\renewcommand{\cftsecleader}{\normalfont\hfill}
\renewcommand{\cftsubsecfont}{\footnotesize}
\renewcommand{\cftsubsecpagefont}{\footnotesize}
\renewcommand{\cftsubsecleader}{\normalfont\hfill}
\renewcommand{\cftdot}{}

% #1 name, #2 complexity / note, #3 description
\newcommand{\entry}[3]{%
  \addcontentsline{toc}{subsection}{#1}%
  \vspace{0.8ex}%
  \noindent\colorbox{hdbg}{\parbox{\dimexpr\columnwidth-2\fboxsep}{%
    \strut\textbf{#1}\hfill{\footnotesize\textsf{#2}}}}\par
  \vspace{0.4ex}%
  {\footnotesize #3\par}%
}
% same, but no index line -- for a second listing belonging to the entry above
\newcommand{\entryn}[3]{%
  \vspace{0.8ex}%
  \noindent\colorbox{hdbg}{\parbox{\dimexpr\columnwidth-2\fboxsep}{%
    \strut\textbf{#1}\hfill{\footnotesize\textsf{#2}}}}\par
  \vspace{0.4ex}%
  {\footnotesize #3\par}%
}

\newcommand{\code}[1]{{\ttfamily\small #1}}
% compact two-column description table, sized for a 93mm column
\newenvironment{ctab}[2]{%
  \begingroup\setlength{\tabcolsep}{3pt}%
  \begin{tabular}{@{}p{#1\columnwidth}p{#2\columnwidth}@{}}}%
  {\end{tabular}\endgroup}

\begin{document}

% ------------------------------- front page --------------------------------
% Same layout as KACTL's team page: logo, university, team, members, and the
% contest at the bottom. The logo is kth.pdf next to this file (upload it to
% Overleaf too); if it is missing the cover simply has no logo.
\pagestyle{empty}
\begin{center}
\strut\vspace{2.6cm}\\
\IfFileExists{kth.pdf}{\includegraphics[width=3cm]{kth.pdf}\\[0.4cm]}{}
{\fontsize{25}{30}\selectfont \uni\\[0.5cm]}
{\fontsize{60}{60}\selectfont \team\\[0.9cm]}
{\huge \members\\}
\vfill
{\fontsize{21}{25}\selectfont \contest}\\[0.5cm]
{\huge \today}\\[1cm]
\end{center}
\newpage
\setcounter{page}{1}
% The cover uses \pagestyle{empty}; switch back here. \thispagestyle does NOT
% work for this -- its effect lands on the index page instead of the cover and
% page 1 of the index loses its header.
\pagestyle{fancy}

% ------------------------------ index + body -------------------------------
{\large\bfseries Index}\hfill{\footnotesize\members}
\vspace{0.5ex}
\textcolor{rule}{\hrule height 0.6pt}
\vspace{1.0ex}

\begin{multicols*}{3}
\renewcommand{\contentsname}{\vspace{-3.5ex}}
\tableofcontents
% tocloft forces \thispagestyle{plain} on the table-of-contents page, which
% wipes the header off the index page. Undo it right here -- verified by
% bisecting the preamble against a minimal file.
\thispagestyle{fancy}
\vspace{1.5ex}

\section{Contest setup}

\entry{Header}{paste first}{%
Every entry in this document assumes this block. Ranges are \textbf{half-open}
\code{[l, r)} and graphs are \textbf{0-indexed} unless the entry says otherwise
--- both are stated wherever it matters.}

%HEADER
\begin{lstlisting}
#include <bits/stdc++.h>
#include <cassert>   // bits/stdc++.h no longer pulls this in
using namespace std;
typedef long long ll;
typedef long double ld;
typedef unsigned long long ull;
typedef __int128 i128;
typedef pair<ll, ll> pii;
typedef vector<ll> vi;
typedef vector<vi> matrix;
#define rep(i, a, b) for(ll i = a; i < b; i++)
#define down(i, b, a) for(ll i = b - 1; i >= a; i--)
#define all(v) (v).begin(), (v).end()
const ll INF = 1e18;

int main(){
    ios_base::sync_with_stdio(false);
    cin.tie(NULL);
}
\end{lstlisting}

\entry{Compile and stress test}{}{%
Drop the sanitizers for the real submission --- they cost $\sim$3x.
\code{-Wconversion} additionally catches silent
\code{ll}$\rightarrow$\code{int} narrowing.\\[0.3ex]
Stress test: \code{gen.py} prints a small random case, \code{mine} is the fast
solution, \code{ans} the slow one you trust.}

%NOCOMPILE
\begin{lstlisting}
g++ -std=gnu++17 -O2 -Wall -Wextra -Wshadow \
    -fsanitize=address,undefined -g a.cpp

for i in $(seq 1000); do
  python3 gen.py $i > t.txt
  ./mine < t.txt > o1.txt; ./ans < t.txt > o2.txt
  cmp -s o1.txt o2.txt || { echo "WA $i"; cat t.txt; break; }
done
\end{lstlisting}

\entry{int128 I/O}{}{%
\code{cin} / \code{cout} cannot handle \code{\_\_int128}. Range is
$\pm1.7\cdot10^{38}$, so it is the fix for \code{a*b \% m} when $m$ is up to
$9\cdot10^{18}$: \code{(i128)a * b \% m}.}

\begin{lstlisting}
i128 read128() {
    i128 x = 0; int s = 1; char c = getchar();
    while (c == ' ' || c == '\n') c = getchar();
    if (c == '-') { s = -1; c = getchar(); }
    for (; c >= '0' && c <= '9'; c = getchar())
        x = x * 10 + (c - '0');
    return s * x;
}
void print128(i128 x) {
    if (x < 0) { putchar('-'); x = -x; }
    if (x > 9) print128(x / 10);
    putchar('0' + (int)(x % 10));
}
\end{lstlisting}

\section{Range structures}

All ranges are half-open: \code{[l, r)}, \code{r} not included.

\entry{Fenwick (BIT)}{add, sum, kth $O(\log n)$ \quad mem $n$}{%
Needs an operation with an inverse (sum, xor); min/max/gcd need a segment tree.
$\sim$3x faster than \code{Seg}, $\sim$15x faster than \code{LazySeg}.\\[0.3ex]
Range add + point get: \code{add(l, x)}, \code{add(r, -x)}; then \code{sum(i+1)}
\emph{is} \code{a[i]}. Range add + range sum: two BITs, or \code{LazySeg}.
\code{kth} needs every \code{a[i] >= 0} and \code{k >= 1}. Xor instead of sum:
every \code{+=} becomes \code{\textasciicircum=}, and drop the \code{-} in
\code{sum(l, r)}.}

%LABEL BIT
\begin{lstlisting}
struct BIT {
    int n; vector<ll> t;
    BIT(int n = 0) : n(n), t(n + 1, 0) {}
    BIT(const vector<ll>& a) : BIT((int)a.size()) {
        for (int i = 1; i <= n; i++) {
            t[i] += a[i-1];
            int j = i + (i & -i);
            if (j <= n) t[j] += t[i];
        }
    }
    void add(int i, ll x) {                     // a[i] += x
        for (++i; i <= n; i += i & -i) t[i] += x;
    }
    ll sum(int r) {                             // a[0..r)
        ll s = 0;
        for (; r > 0; r -= r & -r) s += t[r];
        return s;
    }
    ll sum(int l, int r) {                      // a[l..r)
        return l < r ? sum(r) - sum(l) : 0;
    }
    int kth(ll k) {     // min i with sum(i+1) >= k, else n
        int p = 0;
        for (int w = 1 << __lg(max(n, 1)); w; w >>= 1)
            if (p + w <= n && t[p + w] < k) p += w, k -= t[p];
        return p;
    }
};
\end{lstlisting}

\entry{Segment tree}{build $O(n)$, set / query $O(\log n)$ \quad mem $2n$}{%
Point assign + range query.\\[0.3ex]
\begin{ctab}{0.15}{0.80}
\code{T} & what one element is, and what a query gives back.\\
\code{op(a,b)} & answers for two neighbouring pieces, \code{a} the left one;
return the answer for both together. Associative, need not commute.\\
\code{ID} & a \code{T} that \code{op} ignores: \code{op(x,ID) = op(ID,x) = x}.
An empty query returns it.\\
\end{ctab}\\[0.3ex]
{\ttfamily\footnotesize
\begin{ctab}{0.16}{0.79}
query= & op \hfill ID\\
sum & \code{a + b} \hfill \code{0}\\
min & \code{min(a,b)} \hfill \code{LLONG\_MAX}\\
max & \code{max(a,b)} \hfill \code{LLONG\_MIN}\\
gcd & \code{\_\_gcd(a,b)} \hfill \code{0}\\
xor & \code{a \textasciicircum{} b} \hfill \code{0}\\
and & \code{a \& b} \hfill \code{\textasciitilde 0LL}\\
or & \code{a | b} \hfill \code{0}\\
\end{ctab}}\\[0.3ex]
Count of min: \code{T = pair<ll,int>}, keep the smaller \code{.first} and sum
\code{.second} on ties, \code{ID = \{LLONG\_MAX, 0\}}. Struct \code{T}: write
\code{inline static const T ID = T\{...\};}\\[0.3ex]
Read-modify-write is \code{s.set(i, s[i] + v)}. Range updates need
\code{LazySeg}. Two different trees in one problem: copy, rename, edit.}

\begin{lstlisting}
struct Seg {
    using T = ll;
    static T op(T a, T b) { return a + b; }
    static constexpr T ID = 0;

    int n; vector<T> t;
    Seg(int n = 0) : n(n), t(2 * n, ID) {}
    Seg(const vector<T>& a)
            : n((int)a.size()), t(2*a.size()) {
        copy(a.begin(), a.end(), t.begin() + n);
        for (int i = n - 1; i > 0; i--)
            t[i] = op(t[2*i], t[2*i+1]);
    }
    T operator[](int i) const { return t[i + n]; }
    void set(int i, T v) {
        for (t[i += n] = v; i /= 2;)
            t[i] = op(t[2*i], t[2*i+1]);
    }
    T query(int l, int r) {               // combine a[l..r)
        T x = ID, y = ID;              // x = left, y = right
        for (l += n, r += n; l < r; l /= 2, r /= 2) {
            if (l & 1) x = op(x, t[l++]);
            if (r & 1) y = op(t[--r], y);
        }
        return op(x, y);
    }
};
\end{lstlisting}

\entry{Sparse table}{build $O(n \log n)$, query $O(1)$ \quad mem $n \log n$}{%
\textbf{Idempotent} \code{op} only --- min, max, gcd, and, or. \code{op(x,x)} must
be \code{x}, because the two halves of a query overlap. No \code{ID} needed, and
no updates ever. For sum use \code{BIT}, for updates use \code{Seg}.}

\begin{lstlisting}
struct SpTab {
    using T = ll;
    static T op(T a, T b) { return min(a, b); }

    int n; vector<vector<T>> t;
    SpTab(const vector<T>& a) : n((int)a.size()), t(1, a) {
        for (int k = 1; (1 << k) <= n; k++) {
            t.push_back(vector<T>(n - (1 << k) + 1));
            for (int i = 0; i < (int)t[k].size(); i++)
                t[k][i] = op(t[k-1][i],
                             t[k-1][i + (1 << (k-1))]);
        }
    }
    T query(int l, int r) {          // a[l..r), needs l < r
        int k = __lg(r - l);
        return op(t[k][l], t[k][r - (1 << k)]);
    }
};
\end{lstlisting}

\entry{Sparse segment tree}{set / add / query $O(\log C)$ \quad mem $V \log_2 (C/V)$ nodes}{%
Same \code{T / op / ID} as \code{Seg}, but \code{[lo, hi)} is given at
construction, may be $10^{18}$ wide or start negative, and nodes appear only where
you write. \code{ID} is \emph{also} what every untouched position reads as, so pick
the neutral element you want spread over the whole range.\\[0.3ex]
\begin{ctab}{0.19}{0.76}
\code{add(i,v)} & \code{a[i] = op(a[i], v)} in one descent; prefer it over
\code{set(i, s[i]+v)}, which pays two.\\
\code{hint} & 3rd ctor arg, how many \code{set}/\code{add} you expect.
\code{reserve} only; 0 or wrong is safe.\\
\code{kth(k)} & sum config, all \code{a[i] >= 0}, \code{k >= 1}: smallest
\code{i} with \code{query(lo, i+1) >= k}, else \code{hi}.\\
\end{ctab}\\[0.3ex]
Multiset on a value axis: count config, \code{add(v, 1)} / \code{add(v, -1)},
\code{query(a, b+1)} counts $[a,b]$, \code{kth} gives order statistics.\\[0.3ex]
Use it \emph{only} when the positions are not all known up front (interactive, or an
update that depends on an earlier answer). Otherwise coordinate-compress and use
\code{BIT}: 3.5x faster, 4.5x less memory. Nodes $\approx 2V + V\log_2(C/V)$, 16 B
each ($V$ = positions touched), so $V = 3\cdot10^5$ over $C = 10^9$ is $\sim$62 MB.}

\begin{lstlisting}
struct SparseSeg {
    using T = ll;
    static T op(T a, T b) { return a + b; }
    static constexpr T ID = 0;

    struct Nd { int l, r; T v; };
    ll lo, hi;
    vector<Nd> nd;               // node 0 = null, 1 = root

    SparseSeg(ll lo = 0, ll hi = 1, int hint = 0)
            : lo(lo), hi(hi) {
        ll C = max(hi - lo, 1LL), h = max(hint, 1);
        size_t cap = 2 + h * (3 + (C > h ? __lg(C / h) : 0));
        nd.reserve(cap);
        node(); node();
    }
    int node() {
        nd.push_back({0, 0, ID});
        return (int)nd.size() - 1;
    }
    void set(ll i, T v) { pt(i, v, true);  }
    void add(ll i, T v) { pt(i, v, false); }
    T operator[](ll i) { return query(i, i + 1); }
    T query(ll l, ll r) {                      // fold a[l..r)
        return l < r ? qry(1, lo, hi, l, r) : ID;
    }
    ll kth(T k) {
        if (nd[1].v < k) return hi;
        int x = 1; ll l = lo, r = hi;
        while (r - l > 1) {
            ll m = l + (r - l) / 2;
            if (nd[nd[x].l].v >= k) { x = nd[x].l; r = m; }
            else { k -= nd[nd[x].l].v; x = nd[x].r; l = m; }
        }
        return l;
    }
    void pt(ll i, T v, bool assign) {
        int st[64], d = 0, x = 1; ll l = lo, r = hi;
        while (r - l > 1) {
            st[d++] = x;
            ll m = l + (r - l) / 2;
            if (i < m) {
                if (!nd[x].l) { int c = node(); nd[x].l = c; }
                x = nd[x].l; r = m;
            } else {
                if (!nd[x].r) { int c = node(); nd[x].r = c; }
                x = nd[x].r; l = m;
            }
        }
        nd[x].v = assign ? v : op(nd[x].v, v);
        while (d--) {
            int p = st[d];
            nd[p].v = op(nd[nd[p].l].v, nd[nd[p].r].v);
        }
    }
    T qry(int x, ll l, ll r, ll ql, ll qr) {
        if (!x || qr <= l || r <= ql) return ID;
        if (ql <= l && r <= qr) return nd[x].v;
        ll m = l + (r - l) / 2;
        return op(qry(nd[x].l, l, m, ql, qr),
                  qry(nd[x].r, m, r, ql, qr));
    }
};
\end{lstlisting}

\entry{Lazy segment tree}{build $O(n)$, update / query $O(\log n)$ \quad mem $4n$}{%
A node holds \code{t}, the answer for its own piece, and \code{d}, an update already
counted into \code{t} but not yet given to its children.\\[0.3ex]
\begin{ctab}{0.15}{0.80}
\code{T} & what one element is, and what a query gives back.\\
\code{F} & one update: exactly what you hand to \code{update(l, r, f)}.\\
\code{op(a,b)} & merge two neighbouring answers, \code{a} the left one.
Associative, need not commute.\\
\code{ID} & neutral for \code{op}; an empty query returns it.\\
\code{ap(x,f,k)} & a node of \code{k} elements has answer \code{x} and \code{f} now
hits all \code{k}: return the node's new answer. You never see the elements, only
\code{x} and \code{k} --- if the new answer is not decided by those two, lazy
cannot do this update at all.\\
\code{cm(a,f)} & \code{a} was already waiting and \code{f} lands on top: return the
single update meaning ``\code{a} first, then \code{f}''.\\
\code{sh(f,k)} & the same update shifted \code{k} places right. An \code{f} starts at
the first element it covers, so the left child gets \code{f} and the right child
\code{sh(f, size of left child)}. Return \code{f} unchanged unless the update varies
along the range: ``add $m, m{+}1, \dots$'' becomes ``add $m{+}k, m{+}k{+}1, \dots$''.\\
\code{NOP} & an \code{F} meaning nothing waiting. For assigning updates use a value
you will never assign (\code{LLONG\_MIN}).\\
\end{ctab}\\[0.3ex]
{\ttfamily\footnotesize
\begin{ctab}{0.30}{0.65}
f does \hfill query= & ap(x,f,k) \hfill cm(a,f) \hfill ID \hfill NOP\\
x += f \hfill sum & x + f*k \hfill a+f \hfill 0 \hfill 0\\
x += f \hfill min & x + f \hfill a+f \hfill INF \hfill 0\\
x += f \hfill max & x + f \hfill a+f \hfill -INF \hfill 0\\
x = f \hfill sum & f * k \hfill f \hfill 0 \hfill NONE\\
x = f \hfill min & f \hfill f \hfill INF \hfill NONE\\
x = f \hfill max & f \hfill f \hfill -INF \hfill NONE\\
x \textasciicircum= f \hfill xor & k\&1 ? x\textasciicircum f : x \hfill a\textasciicircum f \hfill 0 \hfill 0\\
\end{ctab}}\\[0.3ex]
\code{NONE = LLONG\_MIN}. Affine (one problem mixing ``multiply a range'' and ``add a
range''): \code{F = pair<ll,ll>}, \code{f = \{p,q\}} means $x \mapsto px+q$,
\code{ap = f.first*x + f.second*k}, \code{cm(a,f) = \{a.first*f.first,
f.first*a.second + f.second\}}, \code{NOP = \{1,0\}}.\\[0.3ex]
Check a new row on paper with
$\mathrm{op}(\mathrm{ap}(a,f,k_a), \mathrm{ap}(b,\mathrm{sh}(f,k_a),k_b)) = \mathrm{ap}(\mathrm{op}(a,b), f, k_a{+}k_b)$
and $\mathrm{ap}(\mathrm{ap}(x,a,k),f,k) = \mathrm{ap}(x,\mathrm{cm}(a,f),k)$.
$\sim$4x slower than \code{Seg} --- use \code{Seg} for point updates.}

%LABEL LazySeg
\begin{lstlisting}
struct LazySeg {
    using T = ll;                            // node value
    using F = ll;                            // pending update
    static T op(T a, T b)        { return a + b; }
    static T ap(T x, F f, int k) { return x + f * k; }
    static F cm(F a, F f)        { return a + f; }
    static F sh(F f, int k)      { return f; }
    static constexpr T ID  = 0;
    static constexpr F NOP = 0;

    int n; vector<T> t; vector<F> d;
    LazySeg(int n = 0) : n(n), t(4*n, ID), d(4*n, NOP) {}
    LazySeg(const vector<T>& a) : LazySeg((int)a.size()) {
        if (n) build(1, 0, n, a);
    }
    void update(int l, int r, F f) {
        if (l < r) upd(1, 0, n, l, r, f);
    }
    T query(int l, int r) {
        return l < r ? qry(1, 0, n, l, r) : ID;
    }
    T operator[](int i) { return query(i, i + 1); }

    // node v covers [l, r)
    void build(int v, int l, int r, const vector<T>& a) {
        if (r - l == 1) { t[v] = a[l]; return; }
        int m = (l + r) / 2;
        build(2*v, l, m, a); build(2*v+1, m, r, a);
        t[v] = op(t[2*v], t[2*v+1]);
    }
    void put(int v, int k, F f) {
        t[v] = ap(t[v], f, k); d[v] = cm(d[v], f);
    }
    void push(int v, int l, int m, int r) {
        if (d[v] == NOP) return;
        put(2*v, m - l, d[v]);
        put(2*v+1, r - m, sh(d[v], m - l));
        d[v] = NOP;
    }
    void upd(int v, int l, int r, int ql, int qr, F f) {
        if (ql <= l && r <= qr) { put(v, r - l, f); return; }
        int m = (l + r) / 2;
        push(v, l, m, r);
        if (ql < m) upd(2*v, l, m, ql, qr, f);
        int c = max(0, m - max(ql, l));   // update left of m
        if (m < qr) upd(2*v+1, m, r, ql, qr, sh(f, c));
        t[v] = op(t[2*v], t[2*v+1]);
    }
    T qry(int v, int l, int r, int ql, int qr) {
        if (ql <= l && r <= qr) return t[v];
        int m = (l + r) / 2;
        push(v, l, m, r);
        if (qr <= m) return qry(2*v, l, m, ql, qr);
        if (m <= ql) return qry(2*v+1, m, r, ql, qr);
        return op(qry(2*v, l, m, ql, qr),
                  qry(2*v+1, m, r, ql, qr));
    }
};
\end{lstlisting}

\entry{Implicit treap}{insert / erase / at / query $O(\log n)$ \quad mem $n$}{%
A \emph{sequence} you can insert into and erase from the middle of, by index --- what
\code{vector} cannot do and \code{Seg} cannot do. Indices are positions in the
current sequence, so every insert renumbers everything after it.\\[0.3ex]
\code{agg} is the subtree fold; change the \code{+} in \code{pull} and the \code{T\{\}}
in \code{ag} together to get min/max instead of sum. \code{query(l, r)} folds
\code{[l, r)}.\\[0.3ex]
\code{Treap<ll> t; t.insert(0, 5); t.insert(1, 7); t.at(0); t.erase(0);}}

\begin{lstlisting}
mt19937 rng(chrono::steady_clock::now()
                .time_since_epoch().count());

template<typename T>
struct Treap {
    struct Node {
        T val, agg; ll prio, sz; Node *l, *r;
        Node(T v) : val(v), agg(v), prio(rng()), sz(1),
                    l(nullptr), r(nullptr) {}
    };
    Node* root = nullptr;

    static ll sz(Node* t) { return t ? t->sz : 0; }
    static T  ag(Node* t) { return t ? t->agg : T{}; }
    static void pull(Node* t) {
        if (!t) return;
        t->sz  = 1 + sz(t->l) + sz(t->r);
        t->agg = ag(t->l) + t->val + ag(t->r);   // op
    }
    // first k elements -> a, rest -> b
    static void split(Node* t, ll k, Node*& a, Node*& b) {
        if (!t) { a = b = nullptr; return; }
        if (sz(t->l) >= k) { split(t->l, k, a, t->l); b = t; }
        else { split(t->r, k-sz(t->l)-1, t->r, b); a = t; }
        pull(t);
    }
    static Node* merge(Node* a, Node* b) {
        if (!a || !b) return a ? a : b;
        if (a->prio > b->prio) {
            a->r = merge(a->r, b); pull(a); return a;
        }
        b->l = merge(a, b->l); pull(b); return b;
    }
    void insert(ll k, T val) {       // ends up at index k
        Node *a, *b;
        split(root, k, a, b);
        root = merge(merge(a, new Node(val)), b);
    }
    void erase(ll k) {
        Node *a, *b, *c;
        split(root, k, a, b); split(b, 1, b, c);
        delete b; root = merge(a, c);
    }
    T at(ll k) {
        Node* t = root;
        while (t) {
            ll ls = sz(t->l);
            if (k == ls) return t->val;
            if (k <  ls) t = t->l;
            else { k -= ls + 1; t = t->r; }
        }
        assert(false);
    }
    T query(ll l, ll r) {                  // fold [l, r)
        Node *a, *b, *c;
        split(root, l, a, b); split(b, r - l, b, c);
        T res = ag(b);
        root = merge(a, merge(b, c));
        return res;
    }
    ll size() const { return sz(root); }
};
\end{lstlisting}

\entry{Order-statistic tree}{insert / erase / rank / kth $O(\log n)$}{%
A \code{set} that also answers ``how many are below $x$'' and ``what is the $k$-th''.
\code{find\_by\_order(k)} is an iterator to the \code{k}-th smallest (0-indexed);
\code{order\_of\_key(x)} counts elements \emph{strictly} below \code{x}.\\[0.3ex]
It is a \emph{set}, so duplicates collapse. For a multiset store
\code{pair<val, unique\_id>} and query with \code{\{x, 0\}} / \code{\{x, INF\}}.}

\begin{lstlisting}
#include <ext/pb_ds/assoc_container.hpp>
using namespace __gnu_pbds;
template<class T> using Tree = tree<T, null_type,
    less<T>, rb_tree_tag,
    tree_order_statistics_node_update>;
// Tree<ll> s; s.insert(4);
// *s.find_by_order(0);   s.order_of_key(4);
\end{lstlisting}

\entry{Line container (CHT)}{add / query $O(\log n)$ \quad mem $n$}{%
Upper hull of lines: \code{add(k, m)} inserts $y = kx+m$, \code{query(x)} returns the
\emph{maximum} over all inserted lines at $x$. Lines may arrive in any order.
For a minimum, insert \code{(-k, -m)} and negate the answer.\\[0.3ex]
This is the tool for DP transitions of the form
$dp_i = \max_j (dp_j + a_j \cdot b_i)$ --- insert line
\code{(a[j], dp[j])} once $dp_j$ is known, then \code{query(b[i])}.
Integer arithmetic only; \code{div} floors towards $-\infty$.}

\begin{lstlisting}
struct Line {
    mutable ll k, m, p;
    bool operator<(const Line& o) const { return k < o.k; }
    bool operator<(ll x) const { return p < x; }
};
struct LineContainer : multiset<Line, less<>> {
    static const ll LINF = LLONG_MAX;
    ll fdiv(ll a, ll b) {      // floored division
        return a / b - ((a ^ b) < 0 && a % b);
    }
    bool isect(iterator x, iterator y) {
        if (y == end()) return x->p = LINF, 0;
        if (x->k == y->k) x->p = x->m > y->m ? LINF : -LINF;
        else x->p = fdiv(y->m - x->m, x->k - y->k);
        return x->p >= y->p;
    }
    void add(ll k, ll m) {
        auto z = insert({k, m, 0}), y = z++, x = y;
        while (isect(y, z)) z = erase(z);
        if (x != begin() && isect(--x, y))
            isect(x, y = erase(y));
        while ((y = x) != begin() && (--x)->p >= y->p)
            isect(x, erase(y));
    }
    ll query(ll x) {                   // max over lines
        assert(!empty());
        auto l = *lower_bound(x);
        return l.k * x + l.m;
    }
};
\end{lstlisting}

\entry{Mo's algorithm}{$O((n + q)\sqrt n)$ moves}{%
Offline range queries where no segment tree works, but you \emph{can} cheaply add or
remove one element at an end of the current range: number of distinct values, most
frequent value count, number of pairs of equal values in \code{[l, r)}, ...\\[0.3ex]
You write three functions: \code{add(i, end)} puts \code{a[i]} into the current range,
\code{del(i, end)} takes it out (\code{end} is 0 for the left end, 1 for the right, in
case it matters), \code{calc()} gives the answer for the current range. \code{mo}
sorts the queries so that the two ends move only $O(n\sqrt q)$ steps in total, and
returns the answers in the original order.\\[0.3ex]
Queries are \textbf{half-open} \code{[l, r)}. Set \code{blk} $\approx n / \sqrt q$.
Removal hard but adding easy? sort the same way and, inside each block, restart the
left end every time (``Mo with rollback''). Tree paths: KACTL's \code{moTree} (not
here) or flatten with an Euler tour. Shown with ``count distinct''. From KACTL.}

%LABEL mo
\begin{lstlisting}
vector<int> a, cnt(1 << 20);   // example: distinct values
ll distinct = 0;
void add(int i, int end) { if (cnt[a[i]]++ == 0) distinct++; }
void del(int i, int end) { if (--cnt[a[i]] == 0) distinct--; }
ll calc() { return distinct; }

vi mo(const vector<pii>& Q) {           // Q[i] = [l, r)
    int L = 0, R = 0, blk = 350;          // ~ n / sqrt(q)
    vector<int> s(Q.size()); vi res(Q.size());
    iota(all(s), 0);
    auto K = [&](pii x) {
        return pii(x.first / blk,
                   x.second ^ -(x.first / blk & 1));
    };
    sort(all(s), [&](int x, int y) {
        return K(Q[x]) < K(Q[y]); });
    for (int qi : s) {
        pii q = Q[qi];
        while (L > q.first) add(--L, 0);
        while (R < q.second) add(R++, 1);
        while (L < q.first) del(L++, 0);
        while (R > q.second) del(--R, 1);
        res[qi] = calc();
    }
    return res;
}
\end{lstlisting}

\entry{Set of disjoint intervals}{$O(\log n)$ amortised}{%
Keeps a \code{set} of non-overlapping \textbf{half-open} intervals \code{[L, R)}.
\code{addInterval} merges the new interval with everything it overlaps or touches and
returns an iterator to the merged piece; \code{removeInterval} cuts a range out,
splitting an interval in two if needed. Total covered length, ``is \code{x} covered''
(\code{upper\_bound(\{x, LLONG\_MAX\})} then step back), painting a line, merging
bookings --- all one line each on top of this. From KACTL.}

%LABEL intervals
\begin{lstlisting}
set<pii>::iterator addInterval(set<pii>& is, ll L, ll R) {
    if (L == R) return is.end();
    auto it = is.lower_bound({L, R}), before = it;
    while (it != is.end() && it->first <= R) {
        R = max(R, it->second);
        before = it = is.erase(it);
    }
    if (it != is.begin() && (--it)->second >= L) {
        L = min(L, it->first);
        R = max(R, it->second);
        is.erase(it);
    }
    return is.insert(before, {L, R});
}
void removeInterval(set<pii>& is, ll L, ll R) {
    if (L == R) return;
    auto it = addInterval(is, L, R);
    ll r2 = it->second;
    if (it->first == L) is.erase(it);
    else const_cast<ll&>(it->second) = L;
    if (R != r2) is.emplace(R, r2);
}
\end{lstlisting}

\entry{Fast hash map}{$\sim$3x faster than \code{unordered\_map}}{%
\code{gp\_hash\_table} has the same \code{h[key]}, \code{h.find}, \code{h.count}
interface as \code{unordered\_map}. The custom hash spreads the bits, so keys that are
all multiples of 1024 (or a test built to break \code{std::hash}) do not pile into one
bucket. The last constructor argument is the starting capacity, a power of two. Reach
for it when a \code{map} of $10^6$+ operations is the bottleneck; for keys up to
$\sim 10^7$ a plain \code{vector} is faster still. From KACTL.}

%LABEL hashmap
\begin{lstlisting}
#include <ext/pb_ds/assoc_container.hpp>
struct chash {                 // large odd constant
    const uint64_t C = ll(4e18 * acos(0)) | 71;
    ll operator()(ll x) const {
        return __builtin_bswap64(x * C); }
};
__gnu_pbds::gp_hash_table<ll, int, chash>
    h({}, {}, {}, {}, {1 << 16});
\end{lstlisting}

\section{Union-find and MST}

\entry{DSU}{$O(\alpha(n))$ amortised \quad mem $n$}{%
\code{p[a] < 0} marks a root and then \code{-p[a]} is its component size, otherwise
\code{p[a]} is the parent. Union by size plus path compression.\\[0.3ex]
\textbf{Moving} one element \code{x} to another set (b5xu's \code{UF\_MV}): never move
the node itself, it may be some other element's parent. Keep \code{id[x]}, the node that
currently stands for \code{x}; to move it, make a fresh node (\code{DSU(n + moves)}),
set \code{id[x]} to it and \code{join} it into the target. The old node stays behind as
a ghost, so keep your own per-set counts (or subtract ghosts from \code{size}).}

%LABEL DSU
\begin{lstlisting}
struct DSU {
    vector<int> p;
    DSU(int n = 0) : p(n, -1) {}
    int find(int a) { return p[a]<0 ? a : p[a]=find(p[a]); }
    int size(int a) { return -p[find(a)]; }
    bool same(int a, int b) { return find(a) == find(b); }
    bool join(int a, int b) {   // false if already joined
        a = find(a); b = find(b);
        if (a == b) return false;
        if (p[a] > p[b]) swap(a, b);         // a = bigger
        p[a] += p[b]; p[b] = a;
        return true;
    }
};
\end{lstlisting}

\entry{DSU with rollback}{join / find $O(\log n)$, undo $O(1)$ \quad mem $n+q$}{%
Union-find where every join can be taken back. Path compression is dropped because it
cannot be undone, so \code{find} costs a log. \code{snap()} records the current state,
\code{undo(t)} restores it exactly.\\[0.3ex]
Use it when merges happen on the way \emph{down} a recursion and must be un-merged on
the way back up: offline dynamic connectivity (segment tree over time, see
\code{Trees}), or a batch of MST/connectivity queries that each add a different set of
edges on top of a common prefix.}

%LABEL RDSU
\begin{lstlisting}
struct RDSU {
    vector<int> p;
    vector<pair<int,int>> hist;  // (joined root, its old p)
    RDSU(int n = 0) : p(n, -1) {}
    int find(int a) { while (p[a] >= 0) a = p[a]; return a; }
    int size(int a) { return -p[find(a)]; }
    bool same(int a, int b) { return find(a) == find(b); }
    bool join(int a, int b) {
        a = find(a); b = find(b);
        if (a == b) return false;
        if (p[a] > p[b]) swap(a, b);
        hist.push_back({b, p[b]});
        p[a] += p[b]; p[b] = a;
        return true;
    }
    int snap() { return (int)hist.size(); }
    void undo(int t) {
        while ((int)hist.size() > t) {
            auto [b, old] = hist.back(); hist.pop_back();
            p[p[b]] -= old; p[b] = old;
        }
    }
};
\end{lstlisting}

\entry{Weighted DSU (potentials)}{$O(\log n)$ \quad mem $n$}{%
Union-find that also remembers \emph{differences}. \code{join(a, b, w)} asserts
$w(a) - w(b) = w$ and returns \code{false} if that contradicts what is already known
--- so a run of \code{join} calls decides whether a system of
``$x_i - x_j = c$'' equations is consistent. \code{diff(a, b)} gives
$w(a) - w(b)$, meaningful only when \code{same(a, b)}.\\[0.3ex]
For $\bmod\ k$ versions (``these two are on opposite sides'') use $k = 2$ and keep
every weight in $\{0,1\}$: that is bipartiteness under edge insertions.}

%LABEL WDSU
\begin{lstlisting}
struct WDSU {
    vector<int> p, sz; vector<ll> d;  // d[x]=w(x)-w(root)
    WDSU(int n = 0) : p(n), sz(n, 1), d(n, 0) {
        iota(all(p), 0);
    }
    int find(int a) {
        if (p[a] == a) return a;
        int r = find(p[a]);
        d[a] += d[p[a]]; p[a] = r;     // parent first!
        return r;
    }
    ll pot(int a) { find(a); return d[a]; }
    ll diff(int a, int b) { return pot(a) - pot(b); }
    bool same(int a, int b) { return find(a) == find(b); }
    bool join(int a, int b, ll w) {    // w(a) - w(b) == w
        ll da = pot(a), db = pot(b);
        int ra = find(a), rb = find(b);
        if (ra == rb) return da - db == w;
        if (sz[ra] < sz[rb]) {
            swap(ra, rb); swap(da, db); w = -w;
        }
        p[rb] = ra; sz[ra] += sz[rb];
        d[rb] = da - db - w;
        return true;
    }
};
\end{lstlisting}

\entry{Kruskal}{$O(m \log m)$ \quad mem $n + m$}{%
Needs \code{DSU}. Sort by weight, keep an edge iff it joins two components.
\code{mst.size() < n-1} means the graph was disconnected.\\[0.3ex]
What you usually want is not the tree but the \emph{join order}: after processing all
edges of weight $\le w$ the components are exactly the connected components at that
threshold.\\[0.3ex]
The max edge on the MST path between $a$ and $b$ equals the minimum, over all
$a$--$b$ paths, of the max edge on it (minimax path) --- answer those with \code{LCA}
carrying a max table, not with a new search. For ``is this edge in some MST'' and for
the second-best MST, process all edges of equal weight as one block.}

%DEP DSU
%LABEL kruskal
\begin{lstlisting}
// e = {w, a, b}. returns {total cost, chosen edges}
pair<ll, vector<array<ll,3>>>
kruskal(int n, vector<array<ll,3>> e) {
    sort(all(e));
    DSU d(n); ll cost = 0; vector<array<ll,3>> mst;
    for (auto [w, a, b] : e)
        if (d.join(a, b)) {
            cost += w; mst.push_back({w, a, b});
        }
    return {cost, mst};
}
\end{lstlisting}
\entry{Directed MST (min arborescence)}{$O(m \log n)$}{%
Needs \code{RDSU}. The directed version of an MST: choose the cheapest set of edges
such that every vertex can be reached from the root \code{r} along chosen edges. Every
vertex except \code{r} gets exactly one incoming edge. Kruskal is \textbf{wrong} here
(direction matters). Returns \code{\{cost, par\}} with \code{par[v]} the start of
\code{v}'s chosen incoming edge (\code{par[r] = -1}), or \code{\{-1, \{\}\}} if some
vertex cannot be reached at all.\\[0.3ex]
Idea (Edmonds / Chu-Liu): every vertex grabs its cheapest incoming edge; if that makes a
cycle, shrink the cycle to one vertex, lower the other incoming edges by what the cycle
already paid, and repeat. Heaps that can be merged and an undoable DSU make it fast.
No root given (``any root'')? add a new vertex with edges of a huge cost $M$ to every
vertex, root there, and subtract $M$; if the cost is $\ge 2M$ there is no answer.
Weights may be negative. From KACTL.}

%DEP RDSU
%LABEL dmst
\begin{lstlisting}
struct DEdge { int a, b; ll w; };
struct SkewNode {                      // lazy skew heap
    DEdge key; SkewNode *l, *r; ll delta;
    void prop() {
        key.w += delta;
        if (l) l->delta += delta;
        if (r) r->delta += delta;
        delta = 0;
    }
    DEdge top() { prop(); return key; }
};
SkewNode* merge(SkewNode* a, SkewNode* b) {
    if (!a || !b) return a ? a : b;
    a->prop(), b->prop();
    if (a->key.w > b->key.w) swap(a, b);
    swap(a->l, (a->r = merge(b, a->r)));
    return a;
}
void pop(SkewNode*& a) { a->prop(); a = merge(a->l, a->r); }

pair<ll, vector<int>> dmst(int n, int r, vector<DEdge>& g) {
    RDSU uf(n);
    vector<SkewNode*> heap(n);
    for (DEdge e : g)
        heap[e.b] = merge(heap[e.b], new SkewNode{e});
    ll res = 0;
    vector<int> seen(n, -1), path(n), par(n);
    seen[r] = r;
    vector<DEdge> Q(n), in(n, {-1, -1});
    deque<tuple<int, int, vector<DEdge>>> cycs;
    rep(s, 0, n) {
        int u = s, qi = 0, w;
        while (seen[u] < 0) {
            if (!heap[u]) return {-1, {}};
            DEdge e = heap[u]->top();
            heap[u]->delta -= e.w, pop(heap[u]);
            Q[qi] = e, path[qi++] = u, seen[u] = s;
            res += e.w, u = uf.find(e.a);
            if (seen[u] == s) {          // cycle: contract
                SkewNode* cyc = 0;
                int end = qi, t = uf.snap();
                do cyc = merge(cyc, heap[w = path[--qi]]);
                while (uf.join(u, w));
                u = uf.find(u), heap[u] = cyc, seen[u] = -1;
                cycs.push_front({u, t, {&Q[qi], &Q[end]}});
            }
        }
        rep(i, 0, qi) in[uf.find(Q[i].b)] = Q[i];
    }
    for (auto& [u, t, comp] : cycs) {     // rebuild the tree
        uf.undo(t);
        DEdge inEdge = in[u];
        for (auto& e : comp) in[uf.find(e.b)] = e;
        in[uf.find(inEdge.b)] = inEdge;
    }
    rep(i, 0, n) par[i] = in[i].a;
    return {res, par};
}
\end{lstlisting}


\section{Trees}

\entry{Euler tour}{$O(n)$ \quad mem $n$}{%
The subtree of \code{u} is the range \code{[tin[u], tout[u])} in any array indexed by
\code{tin}, and \code{u} is an ancestor of \code{v} iff
\code{tin[u] <= tin[v] < tout[u]}. Turns ``update / query a whole subtree'' into a
range operation, so pair it with \code{BIT} / \code{Seg} / \code{LazySeg}. Edge values
go on the \emph{lower} endpoint of the edge.\\[0.3ex]
Iterative, so depth never matters.}

%LABEL euler
\begin{lstlisting}
void euler(const vector<vector<int>>& g, int root,
           vector<int>& tin, vector<int>& tout) {
    int n = (int)g.size(), t = 0;
    tin.assign(n, 0); tout.assign(n, 0);
    vector<int> par(n, -1), it(n, 0), st{root};
    tin[root] = t++;
    while (!st.empty()) {
        int u = st.back();
        if (it[u] < (int)g[u].size()) {
            int v = g[u][it[u]++];
            if (v == par[u]) continue;
            par[v] = u; tin[v] = t++; st.push_back(v);
        } else { tout[u] = t; st.pop_back(); }
    }
}
\end{lstlisting}

\entry{LCA (binary lifting)}{build $O(n \log n)$, query $O(\log n)$ \quad mem $n \log n$}{%
\code{up[k][v]} is the $2^k$-th ancestor of \code{v}; the root is its own ancestor so
jumps past it clamp there. Lift the deeper vertex to the other's depth, then jump both
up together while the ancestors differ. Built by BFS, so no recursion depth
trouble.\\[0.3ex]
Max / min / sum edge on a path: add a second table \code{mx} with \code{mx[0][v]} =
weight of the edge to the parent, \code{mx[0][root] = ID},
\code{mx[k][v] = op(mx[k-1][v], mx[k-1][up[k-1][v]])}, and merge \code{mx[i][v]} into
the answer at every jump taken in \code{kth} and \code{lca}.}

%LABEL LCA
\begin{lstlisting}
struct LCA {
    int LG; vector<vector<int>> up; vector<int> dep;
    LCA(const vector<vector<int>>& g, int root = 0) {
        int n = g.size(); LG = 1 + __lg(max(n, 2));
        up.assign(LG, vector<int>(n, root)); dep.assign(n, 0);
        vector<int> order{root};
        vector<bool> seen(n); seen[root] = true;
        for (int i = 0; i < (int)order.size(); i++) {
            int u = order[i];
            for (int v : g[u]) if (!seen[v]) {
                seen[v] = true; up[0][v] = u;
                dep[v] = dep[u] + 1; order.push_back(v);
            }
        }
        for (int k = 1; k < LG; k++)
            for (int v = 0; v < n; v++)
                up[k][v] = up[k-1][up[k-1][v]];
    }
    int kth(int v, int k) {        // root if k is too big
        k = min(k, dep[v]);
        for (int i = 0; k; k >>= 1, i++)
            if (k & 1) v = up[i][v];
        return v;
    }
    int lca(int a, int b) {
        if (dep[a] < dep[b]) swap(a, b);
        a = kth(a, dep[a] - dep[b]);
        if (a == b) return a;
        down(k, LG, 0) if (up[k][a] != up[k][b])
            a = up[k][a], b = up[k][b];
        return up[0][a];
    }
    int dist(int a, int b) {
        return dep[a] + dep[b] - 2 * dep[lca(a, b)];
    }
};
\end{lstlisting}

\entry{Heavy-light decomposition}{build $O(n)$, path / subtree $O(\log^2 n)$ \quad mem $n$}{%
Every vertex points at its heaviest child, so a root-to-leaf walk changes chain at most
$\log n$ times. Positions are handed out so that each chain and each subtree is a
contiguous range; feed those to \code{Seg} / \code{LazySeg} / \code{BIT}, which is the
second log. \code{lca} comes for free, no extra table.\\[0.3ex]
\code{path} calls \code{f} on $O(\log n)$ ranges in \emph{arbitrary} order, so
\code{op} must commute. For a non-commutative op keep two trees, one on \code{pos} and
one on reversed \code{pos}, and collect the two sides of the LCA separately.
\code{edges = true} stores each edge at its lower endpoint and skips the LCA itself;
then \code{sub} should start at \code{pos[u] + 1} too. Recursion depth is $O(n)$ ---
on a path of $2\cdot10^5$ vertices raise the stack or rewrite the two \code{dfs}
calls iteratively.\\[0.3ex]
\code{HLD h(g); Seg s(n); s.set(h.pos[u], val[u]);}\\
\code{h.path(a, b, [\&](int l, int r)\{ ... s.query(l, r) ... \});}}

\begin{lstlisting}
struct HLD {
    int n, t = 0;
    vector<vector<int>> g;      // heavy child kept first
    vector<int> par, dep, sz, head, pos;

    HLD(vector<vector<int>> adj, int root = 0)
        : n(adj.size()), g(move(adj)), par(n, -1), dep(n),
          sz(n), head(n, root), pos(n) {
        dfs(root); build(root, root);
    }
    void dfs(int u) {
        if (par[u] != -1) g[u].erase(find(all(g[u]), par[u]));
        sz[u] = 1;
        for (int& v : g[u]) {
            par[v] = u; dep[v] = dep[u] + 1;
            dfs(v); sz[u] += sz[v];
            if (sz[v] > sz[g[u][0]]) swap(v, g[u][0]);
        }
    }
    void build(int u, int h) {      // heavy child keeps head
        head[u] = h; pos[u] = t++;
        for (int i = 0; i < (int)g[u].size(); i++)
            build(g[u][i], i ? g[u][i] : h);
    }
    int lca(int a, int b) {
        for (; head[a] != head[b]; b = par[head[b]])
            if (dep[head[a]] > dep[head[b]]) swap(a, b);
        return dep[a] < dep[b] ? a : b;
    }
    pair<int,int> sub(int u) {              // subtree range
        return {pos[u], pos[u] + sz[u]};
    }
    template<class F>
    void path(int a, int b, F f, bool edges = false) {
        for (; head[a] != head[b]; b = par[head[b]]) {
            if (dep[head[a]] > dep[head[b]]) swap(a, b);
            f(pos[head[b]], pos[b] + 1);
        }
        if (dep[a] > dep[b]) swap(a, b);
        f(pos[a] + edges, pos[b] + 1);
    }
};
\end{lstlisting}

\entry{Tree diameter}{$O(n)$ \quad mem $n$}{%
Two BFS runs: the farthest vertex from anywhere is an endpoint of some diameter, and
the farthest vertex from \emph{that} is the other endpoint. Unweighted only --- with
weights use the same shape with Dijkstra, or a single DFS keeping the two deepest
children per vertex. Iterative.}

\begin{lstlisting}
// returns {length in edges, endpoint a, endpoint b}
array<int,3> diameter(const vector<vector<int>>& g) {
    int n = (int)g.size();
    auto bfs = [&](int s) {
        vector<int> d(n, -1); d[s] = 0;
        vector<int> q{s};
        for (int i = 0; i < (int)q.size(); i++)
            for (int v : g[q[i]]) if (d[v] < 0) {
                d[v] = d[q[i]] + 1; q.push_back(v);
            }
        int b = s;
        for (int i = 0; i < n; i++)
            if (d[i] > d[b]) b = i;
        return pair<int,int>{b, d[b]};
    };
    int a = bfs(0).first;
    auto [b, len] = bfs(a);
    return {len, a, b};
}
\end{lstlisting}

\entry{Centroid}{$O(n)$ \quad mem $n$}{%
The centroid minimises the largest component left after deleting it; that component is
always $\le n/2$. Every tree has one or two, and two only when they are
adjacent.\\[0.3ex]
Recurse on the pieces (removing the centroid each time) and you get centroid
decomposition: $O(\log n)$ levels, every path in the tree passes through the centroid
of exactly one level, which is how ``count paths with property X'' problems are
solved in $O(n \log n)$.}

\begin{lstlisting}
vector<int> centroids(const vector<vector<int>>& g) {
    int n = (int)g.size();
    vector<int> sub(n, 1), par(n, -1), it(n, 0), st{0}, ord;
    while (!st.empty()) {              // iterative dfs order
        int u = st.back();
        if (it[u] < (int)g[u].size()) {
            int v = g[u][it[u]++];
            if (v == par[u]) continue;
            par[v] = u; st.push_back(v);
        } else { ord.push_back(u); st.pop_back(); }
    }
    for (int u : ord)                        // children first
        if (par[u] >= 0) sub[par[u]] += sub[u];
    vector<int> res;
    for (int u = 0; u < n; u++) {
        int mx = n - sub[u];
        for (int v : g[u]) if (v != par[u])
            mx = max(mx, sub[v]);
        if (mx <= n / 2) res.push_back(u);
    }
    return res;
}
\end{lstlisting}

\entry{Centroid decomposition}{build $O(n \log n)$ \quad mem $n$}{%
Take the centroid \code{c} of the tree, remove it, and recurse on every piece; the
centroids of the pieces become children of \code{c}. The result is a new tree
(\code{cpar[v]} = parent, \code{-1} at the top) of depth $\le \log_2 n$.\\[0.3ex]
The one fact to remember: \textbf{the path between any $u$ and $v$ goes through their
lowest common ancestor in the centroid tree}. So ``something about all paths'' becomes
``for each centroid, something about paths through it'', and every vertex lies in only
$\log n$ pieces.\\[0.3ex]
Classic use, ``nearest red vertex'' with updates: keep \code{best[c]} = smallest
\code{dist(c, red)} over red vertices in \code{c}'s piece. Paint \code{v}: for each
ancestor \code{c} of \code{v} (walk \code{cpar}), \code{best[c] = min(best[c],
dist(c, v))}. Query \code{v}: min over ancestors \code{c} of
\code{best[c] + dist(c, v)}. \code{dist} from \code{LCA} on the \emph{original} tree;
$O(\log^2 n)$ per operation. Counting paths of length $k$: at each centroid, count
depths of each piece with an array, combine, then subtract pairs from the same piece.
Iterative; b5xu's \code{CD\_tree}, rewritten.}

%LABEL centroidTree
\begin{lstlisting}
vector<int> centroidTree(const vector<vector<int>>& g) {
    int n = (int)g.size();
    vector<int> cpar(n, -1), sz(n), pr(n);
    vector<char> del(n, 0);
    vector<pair<int,int>> todo{{0, -1}};  // (vertex, parent)
    while (!todo.empty()) {
        auto [s, up] = todo.back(); todo.pop_back();
        vector<int> ord{s}; pr[s] = -1;     // bfs the piece
        for (int i = 0; i < (int)ord.size(); i++)
            for (int v : g[ord[i]])
                if (!del[v] && v != pr[ord[i]])
                    pr[v] = ord[i], ord.push_back(v);
        for (int i = (int)ord.size() - 1; i >= 0; i--) {
            int u = ord[i]; sz[u] = 1;
            for (int v : g[u])
                if (!del[v] && v != pr[u]) sz[u] += sz[v];
        }
        int tot = (int)ord.size(), c = s;
        for (bool moved = true; moved;) {   // go heavy
            moved = false;
            for (int v : g[c])
                if (!del[v] && v != pr[c] && 2*sz[v] > tot) {
                    c = v; moved = true; break;
                }
        }
        cpar[c] = up; del[c] = 1;
        for (int v : g[c])
            if (!del[v]) todo.push_back({v, c});
    }
    return cpar;
}
\end{lstlisting}

\entry{Virtual tree (compress tree)}{$O(k \log k)$ per query}{%
Needs \code{LCA} and \code{euler} (same root). A query names $k$ important vertices of a
huge tree, and the sum of all $k$ is small. Build a small tree containing just those
vertices plus the LCAs needed to keep the shape (at most $k-1$ extra), then run the
normal tree DP on it. Edge lengths are \code{dep[b] - dep[a]} from the \code{LCA} struct
(or weighted depth).\\[0.3ex]
Returns pairs \code{(parent index, vertex)}; index 0 is the root and points to itself,
and parents always come earlier in the list, so a loop from the back is a bottom-up DP.
Sorting by \code{tin} is the whole trick: neighbours in that order have all the LCAs you
need. From KACTL.}

%DEP LCA euler
%LABEL vtree
\begin{lstlisting}
vector<pair<int,int>> virtualTree(LCA& L,
        const vector<int>& tin, vector<int> li) {
    static vector<int> rev; rev.resize(tin.size());
    auto cmp = [&](int a, int b) { return tin[a] < tin[b]; };
    sort(all(li), cmp);
    int m = (int)li.size() - 1;
    rep(i, 0, m) li.push_back(L.lca(li[i], li[i+1]));
    sort(all(li), cmp);
    li.erase(unique(all(li)), li.end());
    rep(i, 0, (ll)li.size()) rev[li[i]] = (int)i;
    vector<pair<int,int>> ret = {{0, li[0]}};
    rep(i, 0, (ll)li.size() - 1) {
        int a = li[i], b = li[i+1];
        ret.push_back({rev[L.lca(a, b)], b});
    }
    return ret;
}
\end{lstlisting}

\entry{Rerooting DP}{$O(n)$ \quad mem $n$}{%
Answers a subtree DP for \emph{every} root at once, in two passes: up to collect each
subtree, then down to hand each child what the rest of the tree contributes. Shown for
$\mathrm{ans}[u] = \sum_v \mathrm{dist}(u,v)$.\\[0.3ex]
The pattern generalises: whatever \code{dp} accumulates from children, in the second
pass give child \code{v} the value ``everything except \code{v}'s subtree'', which for
a sum is \code{dp[u] - contribution(v)} and for a max needs the top two children.}

\begin{lstlisting}
// ans[u] = sum over v of dist(u, v)
vi sumOfDists(const vector<vector<int>>& g) {
    int n = (int)g.size();
    vi sub(n, 1), dp(n, 0), ans(n, 0);
    vector<int> par(n, -1), it(n, 0), st{0}, ord;
    while (!st.empty()) {
        int u = st.back();
        if (it[u] < (int)g[u].size()) {
            int v = g[u][it[u]++];
            if (v == par[u]) continue;
            par[v] = u; st.push_back(v);
        } else { ord.push_back(u); st.pop_back(); }
    }
    for (int u : ord) if (par[u] >= 0) {     // up: children
        sub[par[u]] += sub[u];
        dp[par[u]]  += dp[u] + sub[u];
    }
    ans[ord.back()] = dp[ord.back()];        // root
    down(i, (ll)ord.size(), 0) {             // down: parents
        int u = ord[i];
        if (par[u] >= 0)
            ans[u] = ans[par[u]] + n - 2 * sub[u];
    }
    return ans;
}
\end{lstlisting}

\entry{Offline DFS + rollback}{$O((n + q) \cdot \mathrm{apply})$ \quad mem $n + q$}{%
Not a structure, a shape. Bucket every query at the vertex it belongs to, then run
\emph{one} DFS from the root holding a global structure that always describes the
current root-to-\code{u} path. Apply \code{u}'s updates on entry, answer \code{u}'s
queries, recurse, and on exit undo exactly what was applied --- so the structure is
correct at every vertex without ever being rebuilt. Needs invertible updates and
offline queries.\\[0.3ex]
\begin{ctab}{0.33}{0.62}
array / counter & subtract what you added\\
diff array by depth & \code{d[dep] += x}, \code{d[dep+len+1] -= x}, prefix-summed
along the path\\
\code{BIT} & \code{add(i, x)} then \code{add(i, -x)}\\
\code{RDSU} & \code{snap()} then \code{undo(t)}\\
\end{ctab}\\[0.3ex]
Replace the recursion over the tree with a recursion over a segment tree of
\emph{time} and the same shape becomes offline dynamic connectivity: put each edge on
the $O(\log q)$ nodes covering the interval it is alive, DFS that tree, \code{join} on
entry, \code{undo} on exit. That is the only place \code{RDSU} is really needed ---
with an array or a BIT the undo is just the negated update.}

%NOCOMPILE
\begin{lstlisting}
void dfs(int u, int par) {
    int t = snapshot();                // e.g. rdsu.snap()
    for (auto& upd : at[u]) apply(upd);          // enter
    for (auto& q : qs[u]) ans[q.id] = get(q);
    for (int v : g[u]) if (v != par) dfs(v, u);
    undo(t);                     // exit: exact inverse
}
\end{lstlisting}

\section{Connectivity}

\entry{SCC (Tarjan, iterative)}{$O(n + m)$ \quad mem $n + m$}{%
Directed. Takes any adjacency of an integer type. No recursion, so depth never matters;
self-loops and multi-edges are fine.\\[0.3ex]
\code{comp} is in \textbf{topological} order --- every edge \code{u -> v} has
\code{comp[u] <= comp[v]} --- so a DP over the condensation is just a loop over
\code{0 .. nc-1}, no toposort. \code{members()} gives the vertices of each comp,
\code{dag()} the condensation with duplicates removed.\\[0.3ex]
A comp holds a cycle iff it has $\ge 2$ vertices, or 1 with a self-loop; the graph is a
DAG iff \code{nc == n} and no self-loop exists. Sources of the condensation have
in-degree 0 in \code{dag()}, sinks out-degree 0 --- count degrees, \code{comp[v] == 0}
is not enough.\\[0.3ex]
\textbf{2-SAT.} Node \code{2i} is $x_i$ false, \code{2i+1} is true, so \code{n = 2*vars}.
A clause $(a \lor b)$ adds \emph{both} \code{!a -> b} and \code{!b -> a}; forcing $a$
true is the single edge \code{!a -> a}. Satisfiable iff
\code{comp[2i] != comp[2i+1]} for every $i$, and then
$x_i$ = \code{comp[2i+1] > comp[2i]} (needs \code{comp} topological, as here).}

%LABEL SCC
\begin{lstlisting}
struct SCC {
    int n, nc = 0;              // nc = number of comps
    vector<vector<int>> g;
    vector<int> comp;           // topological order

    template<class G>          // any vector<vector<intlike>>
    SCC(const G& adj) : n(adj.size()), g(n), comp(n, -1) {
        for (int u = 0; u < n; u++)
            for (auto v : adj[u]) g[u].push_back((int)v);
        vector<int> tin(n, -1), low(n), ptr(n, 0), stk, cs;
        vector<char> on(n, 0);
        int timer = 0;
        for (int r = 0; r < n; r++) {
            if (tin[r] != -1) continue;
            tin[r] = low[r] = timer++;
            stk.push_back(r); on[r] = 1; cs.push_back(r);
            while (!cs.empty()) {
                int u = cs.back();
                if (ptr[u] < (int)g[u].size()) {
                    int v = g[u][ptr[u]++];
                    if (tin[v] == -1) {
                        tin[v] = low[v] = timer++;
                        stk.push_back(v); on[v] = 1;
                        cs.push_back(v);
                    } else if (on[v])
                        low[u] = min(low[u], tin[v]);
                } else {
                    cs.pop_back();
                    if (!cs.empty()) {
                        int p = cs.back();
                        low[p] = min(low[p], low[u]);
                    }
                    if (low[u] == tin[u]) {  // u roots comp
                        int w;
                        do {
                            w = stk.back(); stk.pop_back();
                            on[w] = 0; comp[w] = nc;
                        } while (w != u);
                        nc++;
                    }
                }
            }
        }
        for (int v = 0; v < n; v++)         // -> topological
            comp[v] = nc - 1 - comp[v];
    }
    vector<vector<int>> members() {
        vector<vector<int>> c(nc);
        for (int v = 0; v < n; v++) c[comp[v]].push_back(v);
        return c;
    }
    vector<vector<int>> dag() {      // condensation, dedup
        vector<vector<int>> d(nc);
        for (int u = 0; u < n; u++)
            for (int v : g[u])
                if (comp[u] != comp[v])
                    d[comp[u]].push_back(comp[v]);
        for (auto& a : d) {
            sort(all(a));
            a.erase(unique(all(a)), a.end());
        }
        return d;
    }
};
\end{lstlisting}

\entry{2-SAT}{$O(n + \text{clauses})$ \quad mem same}{%
Needs \code{SCC}. You have $n$ yes/no choices and rules that each talk about at most
\textbf{two} choices: ``$a$ or $b$'', ``not both $a$ and $b$'', ``$a$ forces $b$'',
``$a$ equals $b$''. Find any assignment that obeys every rule, or say none exists.
Three choices in one rule (``$a$ or $b$ or $c$'') is NP-hard --- not this.\\[0.3ex]
A literal is \code{x} for ``choice $x$ is true'' and \code{\~{}x} for ``false''.
\code{solve()} returns \code{false} if impossible, else \code{val[x]} is the answer.\\[0.3ex]
\begin{ctab}{0.35}{0.60}
\code{either(a, b)} & $a \lor b$ (at least one)\\
\code{either(\~{}a, \~{}b)} & not both\\
\code{implies(a, b)} & if $a$ then $b$\\
\code{implies} both ways & $a = b$; with \code{\~{}b}: $a \ne b$\\
\code{setValue(a)} & force $a$ true (\code{\~{}a}: false)\\
\code{atMostOne(\{a,b,..\})} & at most one true, $O(k)$ clauses via helper
variables; ``exactly one'' adds \code{either} over all of them\\
\end{ctab}\\[0.3ex]
Typical disguises: every item goes left or right / up or down / gets one of two
colours, and some pairs of items cannot both take a certain option. Iterative (no deep
recursion), because \code{SCC} is.}

%DEP SCC
%LABEL twosat
\begin{lstlisting}
struct TwoSat {             // literal x = true, ~x = false
    int n; vector<vector<int>> g; vector<char> val;
    TwoSat(int n = 0) : n(n), g(2 * n) {}
    int node(int x) { return x >= 0 ? 2*x + 1 : 2*~x; }
    int addVar() {
        g.emplace_back(); g.emplace_back(); return n++;
    }
    void either(int a, int b) {                // a or b
        g[node(a) ^ 1].push_back(node(b));
        g[node(b) ^ 1].push_back(node(a));
    }
    void implies(int a, int b) { either(~a, b); }
    void setValue(int a) { either(a, a); }
    void atMostOne(const vector<int>& li) {
        if (li.size() <= 1) return;
        int cur = ~li[0];
        for (int i = 2; i < (int)li.size(); i++) {
            int nx = addVar();
            either(cur, ~li[i]); either(cur, nx);
            either(~li[i], nx); cur = ~nx;
        }
        either(cur, ~li[1]);
    }
    bool solve() {
        SCC s(g); val.assign(n, 0);
        for (int i = 0; i < n; i++) {
            if (s.comp[2*i] == s.comp[2*i+1]) return false;
            val[i] = s.comp[2*i+1] > s.comp[2*i];
        }
        return true;
    }
};
\end{lstlisting}

\entry{Bridges, cut vertices, 2ECC}{$O(n + m)$ \quad mem $n + m$}{%
Undirected, given as an \textbf{edge list} (pairs, 2-element arrays or tuples of any
integer type), so \code{bri[i]} indexes your own input. Multi-edges and self-loops are
safe, disconnected input is fine, no recursion.\\[0.3ex]
\begin{ctab}{0.15}{0.80}
\code{bri[i]} & input edge \code{i} is a bridge.\\
\code{art[v]} & \code{v} is a cut vertex.\\
\code{comp[v]} & 2-edge-connected comp: all you reach without crossing a bridge.\\
\code{tree()} & the bridge forest on those comps.\\
\end{ctab}\\[0.3ex]
Bridges are exactly the edges between comps, so \code{tree()} is a forest, one tree per
connected component --- feed it to \code{LCA} / \code{HLD} for ``how many bridges on
the \code{u}--\code{v} path''. Adding one edge kills every bridge on the path joining
its two comps.\\[0.3ex]
The parent edge is skipped by \textbf{edge id}, never by vertex --- by vertex reports
parallel edges as bridges, which is the usual wrong answer. A bridge endpoint need not
be a cut vertex (degree 1) and a cut vertex need not carry one (two cycles sharing it);
the DFS root is one iff it has $\ge 2$ children.}

\begin{lstlisting}
struct Bridges {
    int n, m, nc = 0;                  // nc = number of 2ecc
    vector<vector<pair<int,int>>> g;   // (to, edge id)
    vector<int> comp;                  // 2ecc id per vertex
    vector<char> bri, art;             // per edge / vertex
    vector<pair<int,int>> es;

    template<class E>      // list of pair / array / tuple
    Bridges(int n, const E& edges)
            : n(n), m(edges.size()), g(n), comp(n, -1),
              bri(m, 0), art(n, 0) {
        es.reserve(m);
        for (auto& e : edges)
            es.push_back({(int)get<0>(e), (int)get<1>(e)});
        vector<int> deg(n, 0);
        for (int i = 0; i < m; i++) {
            deg[es[i].first]++; deg[es[i].second]++;
        }
        for (int v = 0; v < n; v++) g[v].reserve(deg[v]);
        for (int i = 0; i < m; i++) {
            g[es[i].first].push_back({es[i].second, i});
            g[es[i].second].push_back({es[i].first, i});
        }
        vector<int> tin(n, -1), low(n), ptr(n, 0), cs;
        vector<int> pe(n, -1);
        int timer = 0;
        for (int r = 0; r < n; r++) {
            if (tin[r] != -1) continue;
            int kids = 0;
            tin[r] = low[r] = timer++; cs.push_back(r);
            while (!cs.empty()) {
                int u = cs.back();
                if (ptr[u] < (int)g[u].size()) {
                    auto [v, id] = g[u][ptr[u]++];
                    if (id == pe[u]) continue;   // skip by ID
                    if (tin[v] == -1) {
                        pe[v] = id; tin[v] = low[v] = timer++;
                        cs.push_back(v); kids += (u == r);
                    } else low[u] = min(low[u], tin[v]);
                } else {
                    cs.pop_back();
                    if (cs.empty()) break;
                    int p = cs.back();
                    low[p] = min(low[p], low[u]);
                    if (low[u] > tin[p]) bri[pe[u]] = 1;
                    if (low[u] >= tin[p]) art[p] = 1;
                }
            }
            art[r] = kids >= 2;                  // root rule
        }
        vector<int> q;                // 2ecc: drop bridges
        for (int r = 0; r < n; r++) {
            if (comp[r] != -1) continue;
            comp[r] = nc; q.clear(); q.push_back(r);
            while (!q.empty()) {
                int u = q.back(); q.pop_back();
                for (auto [v, id] : g[u])
                    if (!bri[id] && comp[v] == -1) {
                        comp[v] = nc; q.push_back(v);
                    }
            }
            nc++;
        }
    }
    Bridges(int n, initializer_list<pair<int,int>> e)
            : Bridges(n, vector<pair<int,int>>(e)) {}

    vector<vector<int>> tree() {     // bridge forest
        vector<vector<int>> t(nc);
        for (int i = 0; i < m; i++) if (bri[i]) {
            int a = comp[es[i].first], b = comp[es[i].second];
            t[a].push_back(b); t[b].push_back(a);
        }
        return t;
    }
};
\end{lstlisting}

\entry{Topological sort (Kahn)}{$O(n + m)$ \quad mem $n$}{%
Directed. Returns an empty vector iff the graph has a cycle, so it doubles as a cycle
test. If you already built \code{SCC} you have a topological order for free in
\code{comp} --- use this only when you do not need the components.\\[0.3ex]
Pop from the back instead of the front and you get a DFS-flavoured order; use a
\code{priority\_queue} for the lexicographically smallest one.}

\begin{lstlisting}
vector<int> topo(const vector<vector<int>>& g) {
    int n = (int)g.size();
    vector<int> deg(n, 0), q;
    for (int u = 0; u < n; u++)
        for (int v : g[u]) deg[v]++;
    for (int u = 0; u < n; u++) if (!deg[u]) q.push_back(u);
    for (int i = 0; i < (int)q.size(); i++)
        for (int v : g[q[i]])
            if (--deg[v] == 0) q.push_back(v);
    return (int)q.size() == n ? q : vector<int>{};
}
\end{lstlisting}

\section{Shortest paths}

\entry{Dijkstra}{$O((n + m) \log n)$ \quad mem $n + m$}{%
\textbf{Non-negative} weights only --- one negative edge and the answer is silently
wrong, use \code{bellman}. \code{d[v] == INF} means unreachable. Pass \code{\&par} to
recover a path: walk \code{par} back from the target and reverse.\\[0.3ex]
Variants that are the same loop with a different relax step: maximise the bottleneck
(\code{max(d[u], w)} with a max-heap), count shortest paths (a second array,
\code{cnt[v] += cnt[u]} on tie), fewest edges among shortest paths (compare
\code{pair<dist, edges>}). For 0/1 weights use a \code{deque} instead of the heap
(push\_front on 0) and drop the log.}

%LABEL dijkstra
\begin{lstlisting}
// g[u] = {(v, w)}
vi dijkstra(const vector<vector<pii>>& g, int s,
            vector<int>* par = nullptr) {
    int n = (int)g.size();
    vi d(n, INF);
    if (par) par->assign(n, -1);
    priority_queue<pii, vector<pii>, greater<pii>> q;
    d[s] = 0; q.push({0, s});
    while (!q.empty()) {
        auto [dv, u] = q.top(); q.pop();
        if (dv > d[u]) continue;            // stale entry
        for (auto [v, w] : g[u])
            if (dv + w < d[v]) {
                d[v] = dv + w;
                if (par) (*par)[v] = (int)u;
                q.push({d[v], v});
            }
    }
    return d;
}
\end{lstlisting}

\entry{Bellman-Ford}{$O(nm)$ \quad mem $n + m$}{%
Handles negative weights and reports a negative cycle reachable from \code{s}. Stops
early when a round changes nothing, so it is fast on easy inputs.\\[0.3ex]
\textbf{Difference constraints.} A system of $x_v - x_u \le w$ is exactly a shortest
path problem: add edge \code{u -> v} of weight \code{w}, add a super-source with a
0-edge to every vertex, run this from it. The resulting \code{d} is a feasible
assignment, and the system is infeasible iff there is a negative cycle. $x_v - x_u \ge
w$ becomes $x_u - x_v \le -w$; $x_v - x_u = w$ is both.\\[0.3ex]
To mark every vertex whose distance is truly $-\infty$, run $n$ more rounds and set
\code{d[v] = -INF} whenever it still relaxes.}

\begin{lstlisting}
// e = {u, v, w}; neg = a negative cycle is reachable from s
vi bellman(int n, const vector<array<ll,3>>& e, int s,
           bool& neg) {
    vi d(n, INF); d[s] = 0; neg = false;
    rep(i, 0, n) {
        bool ch = false;
        for (auto [u, v, w] : e)
            if (d[u] < INF && d[u] + w < d[v]) {
                d[v] = d[u] + w; ch = true;
                if (i == n - 1) neg = true;
            }
        if (!ch) break;
    }
    return d;
}
\end{lstlisting}

\entry{Floyd-Warshall}{$O(n^3)$ \quad mem $n^2$}{%
All pairs, negative weights allowed. Fill \code{d[i][i] = 0}, \code{d[i][j]} = the
\emph{cheapest} edge \code{i -> j} (careful with multi-edges: take a \code{min}), and
\code{INF} elsewhere. Rewrites \code{d} in place.\\[0.3ex]
A negative cycle exists iff \code{d[i][i] < 0} for some \code{i} afterwards; then
\code{d[a][b]} is $-\infty$ whenever \code{d[a][i] < INF} and \code{d[i][b] < INF} for
such an \code{i}. The \code{k} loop must be outermost. The \code{< INF} guards are what
stop \code{INF + INF} from wrapping.\\[0.3ex]
Same triple loop, other semirings: transitive closure (\code{|=} with \code{\&\&}),
maximum bottleneck (\code{max}/\code{min}), minimax path.}

\begin{lstlisting}
void floyd(matrix& d) {
    int n = (int)d.size();
    rep(k, 0, n) rep(i, 0, n) if (d[i][k] < INF)
        rep(j, 0, n) if (d[k][j] < INF)
            d[i][j] = min(d[i][j], d[i][k] + d[k][j]);
}
\end{lstlisting}

\entry{K shortest walks}{$O((m + k n)\log)$ \quad mem $n + m$}{%
The \code{k} smallest \code{s}$\to$\code{t} lengths in ascending order, \textbf{walks}
--- vertices may repeat, so on a graph with a cycle you get the same path plus loops.
That is what most ``k-th shortest route'' problems want. Non-negative weights.
Fewer than \code{k} results means fewer exist.\\[0.3ex]
For genuinely \emph{simple} k shortest paths you need Yen's algorithm --- much longer,
and rarely what is asked.}

\begin{lstlisting}
vi kShortest(const vector<vector<pii>>& g, int s, int t,
             int k) {
    int n = (int)g.size();
    vector<int> cnt(n, 0); vi res;
    priority_queue<pii, vector<pii>, greater<pii>> q;
    q.push({0, s});
    while (!q.empty() && (int)res.size() < k) {
        auto [d, u] = q.top(); q.pop();
        if (cnt[u] >= k) continue;
        cnt[u]++;
        if (u == t) res.push_back(d);
        for (auto [v, w] : g[u])
            if (cnt[v] < k) q.push({d + w, v});
    }
    return res;
}
\end{lstlisting}

\section{Flow and matching}

\entry{Flow: how to spot it}{read this first}{%
Think of every edge as a pipe that carries at most \code{c} units. \textbf{Max flow} is
how much can go from \code{s} to \code{t} at once. A \textbf{cut} is a set of edges
whose removal makes \code{t} unreachable from \code{s}; the cheapest cut always weighs
exactly the max flow, so ``cheapest way to separate'' is a flow problem too.\\[0.3ex]
\textbf{Bipartite} means the things split into two groups (people / jobs, rows /
columns, black / white squares of a chessboard, odd / even numbers) and every allowed
pair has one thing from each group. Test: colour with BFS in two colours; it works iff
no edge joins two same-coloured things (no odd cycle).\\[0.3ex]
\textbf{Build it.} One source \code{S}, one sink \code{T}. \code{S -> x} with how many
times \code{x} may be used, \code{x -> y} for every allowed pair, \code{y -> T} with how
many times \code{y} may be used. Several sources or sinks: one super source / super sink
in front of them.\\[0.3ex]
{\footnotesize
\begin{ctab}{0.47}{0.48}
\textbf{The problem says...} & \textbf{Use}\\
pair as many as possible, two groups, each used once & \code{Matching} (Hopcroft-Karp)\\
same, but things may be used $k$ times & \code{Dinic}, capacity $k$ on \code{S->x}\\
everyone must be paired and each pair has a cost; min (or max) total & \code{Hungarian},
dense $n \le 1000$\\
pairs have costs but not everyone must be paired, or the graph is sparse & \code{MCMF};
\code{run(s,t,m.INFF,1)} stops when pairing more only costs\\
send exactly $k$ units as cheaply as possible & \code{MCMF} \code{run(s, t, k)}\\
fewest things to remove so that no allowed pair is left / every edge is touched & min
vertex cover = matching size, \code{cover()} lists them\\
most things you can pick with no two of them forming an allowed pair & independent set
$= n - $ matching (complement of \code{cover()})\\
pairs inside \emph{one} group (people pair with people) & not bipartite:
\code{generalMatching}\\
cheapest way to disconnect \code{s} from \code{t} by deleting edges & \code{Dinic}, then
\code{cut()}\\
delete \emph{vertices} instead of edges / a vertex may be used $k$ times & split
\code{v} into \code{v\_in -> v\_out} with capacity = its cost or $k$\\
$k$ paths from \code{s} to \code{t} that share no edge (no vertex) & capacity 1 on edges
(and split vertices); max flow is the number of paths\\
fewest paths covering every vertex of a DAG & $n -$ matching, left copy \code{u} to
right copy \code{v} per edge; if paths may share vertices, add all reachable pairs
first\\
choose projects with profit; each needs tools that cost money & project selection: see
\code{Dinic}\\
cut the graph in two, no \code{s}/\code{t} given & \code{globalMinCut}\\
min cut between \emph{every} pair of vertices & \code{gomoryHu}: $n-1$ flows\\
``smallest time / budget so that everything fits'' & binary search the answer, run a
flow inside the check\\
\end{ctab}}\\[0.3ex]
\textbf{Edges must carry at least $l$} (lower bounds): give the edge capacity $c - l$,
add $l$ to \code{need[v]} and subtract $l$ from \code{need[u]}. Add \code{T -> S} with
capacity \code{INF}, a new source \code{S'} to every \code{v} with \code{need[v] > 0}
(capacity \code{need[v]}), and every \code{u} with \code{need[u] < 0} to a new sink
\code{T'}. A valid flow exists iff max flow \code{S' -> T'} fills all of those
edges.\\[0.3ex]
\textbf{Do not fear the size}: Dinic is far faster than its bound; $10^5$ edges is
routine. If it feels too slow, the graph is usually too big --- look for a smaller
graph (merge equal vertices, only keep useful edges), not a faster flow.}

\entry{Dinic (max flow, iterative)}{$O(n^2 m)$, unit caps $O(m\sqrt n)$ \quad mem $n + m$}{%
Directed, 0-indexed. No recursion, so path length never matters; multi-edges,
self-loops and 0-capacity edges are fine. \code{add(u, v, c)} returns an \textbf{even}
edge id, \code{id \^{} 1} is its reverse; an \textbf{undirected} edge is
\code{add(u, v, c, c)}.\\[0.3ex]
\begin{ctab}{0.22}{0.73}
\code{flow(id)} & flow on edge \code{id}; negative = it runs \code{v -> u}.\\
\code{inS(v)} & after \code{maxflow}: \code{v} on the \code{s} side of the min cut
closest to \code{s}.\\
\code{cut()} & edge ids crossing $S \to T$, all saturated; their \code{c} sums to the
flow. An odd id = undirected edge crossing reversed, \code{id \& \~{}1} is what
\code{add} returned.\\
\code{reset()} & all flow back to 0, edges kept.\\
\end{ctab}\\[0.3ex]
\textbf{Config.} \code{F} = capacity type (\code{int} saves memory, \code{double} works
with \code{EPS = 1e-9}). \code{INF} is max/4, so up to 3 can add up; a result
\code{>= INF} means no finite cut.\\[0.3ex]
\code{maxflow} is \textbf{incremental}: add edges or raise \code{d.e[id].c}, call it
again, and it returns only the \emph{extra} flow. A different \code{(s, t)} also
continues from the old flow --- \code{reset()} first.\\[0.3ex]
\textbf{Recipes.} Bipartite matching: use \code{Matching} below, about 5x faster than
the same graph here. Min vertex cover in a bipartite graph = max matching (K\"onig), and
max independent set = $n$ minus that.
Vertex capacities: split \code{v} into \code{v\_in -> v\_out} with that capacity.
Minimum path cover of a DAG = $n$ minus max matching on the split graph.
Project selection / max closure: source$\to$profit nodes, cost nodes$\to$sink,
prerequisites as \code{INF} edges, answer = total profit $-$ max flow.\\[0.3ex]
Measured at -O2: random $10^5$ vertices $10^6$ edges 0.4\,s; unit-cap bipartite
$2\cdot10^5$ / $1.2\cdot10^6$ 0.8\,s; a $300\times300$ grid with random two-way caps
4.5\,s --- a bad case for every Dinic.}

%LABEL Dinic
\begin{lstlisting}
struct Dinic {
    using F = ll;
    static constexpr F EPS = 0;            // double: 1e-9
    static constexpr F INF = numeric_limits<F>::max() / 4;
    struct E { int to; F c, f; };          // cap, flow
    int n; vector<E> e; vector<vector<int>> g;
    vector<int> lv, it;                    // level, edge ptr

    Dinic(int n = 0) : n(n), g(n), lv(n), it(n) {}
    int add(int u, int v, F c, F rc = 0) {     // id of u->v
        g[u].push_back(e.size()); e.push_back({v, c, 0});
        g[v].push_back(e.size()); e.push_back({u, rc, 0});
        return (int)e.size() - 2;
    }
    F res(int id) { return e[id].c - e[id].f; }
    F flow(int id) { return e[id].f; }
    bool inS(int v) { return lv[v] >= 0; }
    void reset() { for (auto& x : e) x.f = 0; }

    bool bfs(int s, int t) {
        fill(lv.begin(), lv.end(), -1); lv[s] = 0;
        vector<int> q{s};
        for (int i = 0; i < (int)q.size() && lv[t] < 0; i++)
            for (int id : g[q[i]]) {
                int v = e[id].to;
                if (lv[v] < 0 && res(id) > EPS)
                    lv[v] = lv[q[i]] + 1, q.push_back(v);
            }
        return lv[t] >= 0;
    }
    F maxflow(int s, int t) {
        F tot = 0; vector<int> st;         // path, edge ids
        while (bfs(s, t)) {
            fill(it.begin(), it.end(), 0);
            int u = s;
            while (true) {
                if (u == t) {                  // augment
                    F b = INF; int k = 0;  // k = 1st sat.
                    for (int id : st) b = min(b, res(id));
                    for (int i = st.size(); i--;) {
                        e[st[i]].f += b; e[st[i] ^ 1].f -= b;
                        if (res(st[i]) <= EPS) k = i;
                    }
                    tot += b; st.resize(k);
                    u = k ? e[st.back()].to : s;
                } else if (it[u] < (int)g[u].size()) {
                    int id = g[u][it[u]], v = e[id].to;
                    if (lv[v] == lv[u] + 1 && res(id) > EPS)
                        st.push_back(id), u = v;
                    else it[u]++;
                } else {                       // dead end
                    if (u == s) break;
                    lv[u] = -1;
                    u = e[st.back() ^ 1].to;
                    st.pop_back(); it[u]++;
                }
            }
        }
        return tot;
    }
    vector<int> cut() {                    // after maxflow
        vector<int> r;
        for (int id = 0; id < (int)e.size(); id++) {
            int u = e[id ^ 1].to, v = e[id].to;
            if (e[id].c > 0 && inS(u) && !inS(v))
                r.push_back(id);
        }
        return r;
    }
};
\end{lstlisting}

\entry{Min cost max flow}{$O(\mathrm{aug} \cdot m \log n)$ \quad mem $n + m$}{%
Dijkstra with potentials. \textbf{Negative costs are fine} (one SPFA pass first);
negative \emph{cycles} are not. No recursion. \code{add(u, v, cap, cost)} returns the
edge id, cost is per unit; \code{flow(id)} reads it back.\\[0.3ex]
\begin{ctab}{0.36}{0.59}
\code{run(s, t)} & \code{\{flow, cost\}}: max flow, min cost among those.\\
\code{run(s, t, k)} & at most \code{k} units, min cost.\\
\code{run(s, t, m.INFF, 1)} & \emph{any} amount, min cost: stops once one more unit
costs \code{>= 0}.\\
\end{ctab}\\[0.3ex]
\textbf{Config.} \code{F} flow type, \code{C} cost type (\code{double} works --- compare
with a tolerance). \code{INFF} / \code{INFC} are max/4.\\[0.3ex]
\textbf{Max} profit: add every edge with \code{-cost}, negate the answer. Undirected edge
with cost: two separate \code{add} calls. Paths go cheapest first, so after $k$ units
the cost is the minimum for exactly $k$. \code{run} is incremental like \code{Dinic}.\\[0.3ex]
One Dijkstra per augmenting path ($\le$ total flow), fine up to $\sim10^4$--$10^5$ of
them. Measured: $10^4$ vertices $10^5$ edges 0.4\,s; a $300\times300$ assignment
0.42\,s, while \code{Hungarian} does $1000\times1000$ in 0.08\,s --- use that for a
dense square cost matrix.}

%LABEL MCMF
\begin{lstlisting}
struct MCMF {
    using F = ll; using C = ll;
    static constexpr F INFF = numeric_limits<F>::max() / 4;
    static constexpr C INFC = numeric_limits<C>::max() / 4;
    struct E { int to; F c, f; C w; };     // cap, flow, cost
    int n; vector<E> e; vector<vector<int>> g;
    vector<C> pi, dist; vector<int> pe;    // pe = parent edge

    MCMF(int n = 0) : n(n), g(n), pi(n, 0), dist(n), pe(n) {}
    int add(int u, int v, F c, C w) {
        g[u].push_back(e.size()); e.push_back({v, c, 0, w});
        g[v].push_back(e.size()); e.push_back({u, 0, 0, -w});
        return (int)e.size() - 2;
    }
    F flow(int id) { return e[id].f; }

    void spfa(int s) {                     // exact potentials
        fill(pi.begin(), pi.end(), INFC); pi[s] = 0;
        vector<char> inq(n, 0); deque<int> q{s};
        while (!q.empty()) {
            int u = q.front(); q.pop_front(); inq[u] = 0;
            for (int id : g[u]) {
                auto& x = e[id];
                if (x.c - x.f > 0 && pi[u] + x.w < pi[x.to]) {
                    pi[x.to] = pi[u] + x.w;
                    if (!inq[x.to])
                        inq[x.to] = 1, q.push_back(x.to);
                }
            }
        }
        for (auto& p : pi) if (p == INFC) p = 0;  // unreached
    }
    bool dijkstra(int s, int t) {
        fill(dist.begin(), dist.end(), INFC); dist[s] = 0;
        priority_queue<pair<C,int>, vector<pair<C,int>>,
                       greater<pair<C,int>>> pq;
        pq.push({0, s});
        while (!pq.empty()) {
            auto [d, u] = pq.top(); pq.pop();
            if (d > dist[u]) continue;
            for (int id : g[u]) {
                auto& x = e[id];
                C nd = d + x.w + pi[u] - pi[x.to];   // >= 0
                if (x.c - x.f > 0 && nd < dist[x.to]) {
                    dist[x.to] = nd; pe[x.to] = id;
                    pq.push({nd, x.to});
                }
            }
        }
        for (int v = 0; v < n; v++)
            if (dist[v] < INFC) pi[v] += dist[v];
        return dist[t] < INFC;
    }
    pair<F,C> run(int s, int t, F lim = INFF, bool neg = 0) {
        F fl = 0; C cost = 0;
        for (auto& x : e)
            if (x.c - x.f > 0 && x.w < 0) { spfa(s); break; }
        while (fl < lim && dijkstra(s, t)) {
            C pc = pi[t] - pi[s];          // cost of one unit
            if (neg && pc >= 0) break;
            F b = lim - fl;
            for (int v = t; v != s; v = e[pe[v] ^ 1].to)
                b = min(b, e[pe[v]].c - e[pe[v]].f);
            for (int v = t; v != s; v = e[pe[v] ^ 1].to)
                e[pe[v]].f += b, e[pe[v] ^ 1].f -= b;
            fl += b; cost += b * pc;
        }
        return {fl, cost};
    }
};
\end{lstlisting}

\entry{Bipartite matching (Hopcroft-Karp)}{$O(m\sqrt n)$ \quad mem $n + m$}{%
Unweighted maximum matching, iterative, multi-edges fine. Left \code{[0, nl)} and right
\code{[0, nr)} are numbered \textbf{independently}; \code{add(l, r)}, then
\code{run()} returns the size $k$. \code{ml[l]} / \code{mr[r]} give the partner or
\code{-1}.\\[0.3ex]
\textbf{K\"onig.} After \code{run}, \code{auto [L, R] = m.cover()} is a minimum vertex
cover, \code{|L| + |R| == k}. Its complement --- left \emph{not} in \code{L} plus right
\emph{not} in \code{R} --- is a maximum independent set, size \code{nl + nr - k}.\\[0.3ex]
\code{run} is incremental: add edges and call it again, old pairs are kept. You may also
pre-fill \code{ml}/\code{mr} with a valid greedy matching.\\[0.3ex]
Weighted: \code{Hungarian} or \code{MCMF}. General (non-bipartite) graphs need
\code{generalMatching}, not this. Measured: $10^5 + 10^5$ vertices, $10^6$ edges 0.15\,s, 15\,MB.}

%LABEL Matching
\begin{lstlisting}
struct Matching {
    int nl, nr; vector<vector<int>> g;
    vector<int> ml, mr, d, it;             // d = bfs layer

    Matching(int nl, int nr)
        : nl(nl), nr(nr), g(nl), ml(nl, -1), mr(nr, -1) {}
    void add(int l, int r) { g[l].push_back(r); }

    bool bfs() {                           // layers
        d.assign(nl, -1); vector<int> q; bool found = false;
        for (int l = 0; l < nl; l++)
            if (ml[l] < 0) d[l] = 0, q.push_back(l);
        for (int i = 0; i < (int)q.size(); i++)
            for (int r : g[q[i]]) {
                int w = mr[r];
                if (w < 0) found = true;
                else if (d[w] < 0)
                    d[w] = d[q[i]] + 1, q.push_back(w);
            }
        return found;
    }
    bool aug(int s) {                      // s = free left
        vector<int> st{s};
        while (!st.empty()) {
            int u = st.back();
            if (it[u] == (int)g[u].size()) {
                d[u] = -1; st.pop_back(); continue;
            }
            int r = g[u][it[u]], w = mr[r];
            if (w < 0) {                   // free: flip path
                for (int x : st)
                    ml[x] = g[x][it[x]], mr[ml[x]] = x;
                return true;
            }
            if (d[w] == d[u] + 1) st.push_back(w);
            else it[u]++;
        }
        return false;
    }
    int run() {
        while (bfs()) {
            it.assign(nl, 0);
            for (int l = 0; l < nl; l++)
                if (ml[l] < 0) aug(l);
        }
        int k = 0;
        for (int l = 0; l < nl; l++) k += ml[l] >= 0;
        return k;
    }
    pair<vector<int>, vector<int>> cover() { // after run
        vector<char> vl(nl, 0), vr(nr, 0);
        vector<int> q, L, R;
        for (int l = 0; l < nl; l++)
            if (ml[l] < 0) vl[l] = 1, q.push_back(l);
        for (int i = 0; i < (int)q.size(); i++)
            for (int r : g[q[i]]) if (!vr[r]) {
                vr[r] = 1; int w = mr[r];
                if (w >= 0 && !vl[w])
                    vl[w] = 1, q.push_back(w);
            }
        for (int l = 0; l < nl; l++)
            if (!vl[l]) L.push_back(l);
        for (int r = 0; r < nr; r++)
            if (vr[r]) R.push_back(r);
        return {L, R};
    }
};
\end{lstlisting}

\entry{Hungarian (assignment)}{$O(n^2 m)$ \quad mem $nm$}{%
\code{n} rows (workers), \code{m >= n} columns (jobs), \code{a[i][j]} = cost of giving
job \code{j} to worker \code{i}. Every worker gets exactly one job, no job is used
twice, and the total is as small as possible. The constructor does all the work:
\code{Hungarian h(a);} then \code{h.cost} and \code{h.job[i]} (the column row \code{i}
got). Costs may be negative; \code{C = double} works.\\[0.3ex]
\begin{ctab}{0.30}{0.65}
maximise & pass \code{-a}, negate \code{cost}.\\
more rows than columns & transpose: jobs choose workers.\\
forbidden pair & cost \code{BIG}, with \code{n*BIG} far below \code{INF} and
\code{BIG} above any real total; \code{h.cost >= BIG} means one had to be used.\\
a row may stay idle & add one 0-cost dummy column per such row.\\
a worker takes up to $k$ jobs & repeat its row $k$ times.\\
\end{ctab}\\[0.3ex]
Not the same thing as \code{Matching}: that one maximises the \emph{number} of pairs,
this one fixes the number and minimises the \emph{cost}. Measured at -O2:
$n = m = 500$ 0.02\,s, $1000$ 0.08\,s. Internally 1-indexed, column 0 is a virtual
helper --- that is why the matrix is read as \code{a[i0-1][j-1]}.}

%LABEL hungarian
\begin{lstlisting}
struct Hungarian {
    using C = ll;
    static constexpr C INF = numeric_limits<C>::max() / 4;
    C cost; vector<int> job;

    template<class M>                  // vector<vector<C>>
    Hungarian(const M& a) {
        int n = a.size(), m = n ? a[0].size() : 0;
        vector<C> u(n + 1), v(m + 1);  // potentials, 1-idx
        vector<int> p(m + 1), way(m + 1); // p[j]: row of j
        for (int i = 1; i <= n; i++) {
            p[0] = i; int j0 = 0;
            vector<C> mn(m + 1, INF);
            vector<char> used(m + 1, 0);
            do {
                used[j0] = 1;
                int i0 = p[j0], j1 = 0; C dl = INF;
                for (int j = 1; j <= m; j++) if (!used[j]) {
                    C cur = a[i0-1][j-1] - u[i0] - v[j];
                    if (cur < mn[j]) mn[j] = cur, way[j] = j0;
                    if (mn[j] < dl) dl = mn[j], j1 = j;
                }
                for (int j = 0; j <= m; j++)
                    if (used[j]) u[p[j]] += dl, v[j] -= dl;
                    else mn[j] -= dl;
                j0 = j1;
            } while (p[j0]);
            do {                        // flip the path back
                int j1 = way[j0]; p[j0] = p[j1]; j0 = j1;
            } while (j0);
        }
        job.assign(n, -1); cost = 0;
        for (int j = 1; j <= m; j++)
            if (p[j]) job[p[j] - 1] = j - 1;
        for (int i = 0; i < n; i++) cost += a[i][job[i]];
    }
};
\end{lstlisting}

\entry{Global min cut (Stoer-Wagner)}{$O(n^3)$ \quad mem $n^2$}{%
Cheapest set of edges whose removal splits an \textbf{undirected} graph into two
non-empty parts, when nobody tells you which two vertices must end up apart. With
\code{Dinic} you would need $n-1$ flows; this does it in one go from an adjacency
\emph{matrix} of edge weights (add parallel edges together, keep the diagonal 0).
Returns \code{\{weight, the vertices of one side\}}.\\[0.3ex]
How it works, roughly: repeatedly grow a set by always adding the vertex most strongly
connected to it. The last vertex added, cut away from everything else, is a candidate;
then it is merged with the second-to-last and the process repeats. Weights must be
$\ge 0$. ``Is the graph $k$-edge-connected'' is \code{globalMinCut(..).first >= k}.
From KACTL.}

%LABEL globalMinCut
\begin{lstlisting}
pair<ll, vector<int>> globalMinCut(vector<vi> mat) {
    pair<ll, vector<int>> best = {LLONG_MAX, {}};
    int n = (int)mat.size();
    vector<vector<int>> co(n);
    rep(i, 0, n) co[i] = {(int)i};
    rep(ph, 1, n) {
        vi w = mat[0];
        size_t s = 0, t = 0;
        rep(it, 0, n - ph) {     // O(n^2); heap -> O(m log n)
            w[t] = LLONG_MIN;
            s = t, t = max_element(all(w)) - w.begin();
            rep(i, 0, n) w[i] += mat[t][i];
        }
        best = min(best, {w[t] - mat[t][t], co[t]});
        co[s].insert(co[s].end(), all(co[t]));
        rep(i, 0, n) mat[s][i] += mat[t][i];
        rep(i, 0, n) mat[i][s] = mat[s][i];
        mat[0][t] = LLONG_MIN;
    }
    return best;
}
\end{lstlisting}

\entry{Gomory-Hu tree (all-pairs min cut)}{$n - 1$ calls to \code{Dinic}}{%
Needs \code{Dinic}. For an \textbf{undirected} graph, builds a tree on the same vertices
such that the min cut (= max flow) between \emph{any} two vertices \code{u, v} equals
the \textbf{smallest edge weight on the tree path} from \code{u} to \code{v}. So $n^2$
flow questions become $n-1$ flows plus tree queries.\\[0.3ex]
Input edges are \code{\{u, v, cap\}}; output tree edges are \code{\{i, parent, flow\}}.
For many queries, build \code{LCA} on the tree with a min table (see \code{LCA}), or for
$n \le 3000$ just run a DFS from every vertex. The smallest tree edge overall is also
the global min cut. Gusfield's version, from KACTL.}

%DEP Dinic
%LABEL gomoryHu
\begin{lstlisting}
vector<array<ll,3>> gomoryHu(int n, vector<array<ll,3>> ed) {
    vector<array<ll,3>> tree;
    vector<int> par(n, 0);
    for (int i = 1; i < n; i++) {
        Dinic d(n);
        for (auto [u, v, c] : ed) d.add(u, v, c, c);
        tree.push_back({i, par[i], d.maxflow(i, par[i])});
        for (int j = i + 1; j < n; j++)
            if (par[j] == par[i] && d.inS(j)) par[j] = i;
    }
    return tree;
}
\end{lstlisting}

\entry{Matrix inverse mod p}{$O(n^3)$ \quad mem $n^2$}{%
Needs \code{power} (Modular basics). Inverts \code{A} modulo a prime \code{md} in place
and returns the rank; if the rank is below $n$ the matrix has no inverse and \code{A} is
garbage. Used by \code{generalMatching} below, but on its own it also solves
$Ax = b \pmod p$ for many different \code{b} at once ($x = A^{-1}b$).\\[0.3ex]
\code{md} must be below $2^{31}$ (because of \code{power}). For a real-valued inverse,
use \code{gauss} on each unit vector instead.}

%DEP modbasics
%LABEL matInv
\begin{lstlisting}
int matInv(vector<vi>& A, ll md) {
    int n = (int)A.size(); vector<int> col(n);
    vector<vi> tmp(n, vi(n));
    rep(i, 0, n) tmp[i][i] = 1, col[i] = i;
    rep(i, 0, n) {
        int r = i, c = i;
        rep(j, i, n) rep(k, i, n) if (A[j][k]) {
            r = j; c = k; goto found;
        }
        return i;
found:
        A[i].swap(A[r]); tmp[i].swap(tmp[r]);
        rep(j, 0, n) swap(A[j][i], A[j][c]),
                     swap(tmp[j][i], tmp[j][c]);
        swap(col[i], col[c]);
        ll v = power(A[i][i], md - 2, md);
        rep(j, i + 1, n) {
            ll f = A[j][i] * v % md;
            A[j][i] = 0;
            rep(k, i + 1, n)
                A[j][k] = (A[j][k] - f * A[i][k]) % md;
            rep(k, 0, n)
                tmp[j][k] = (tmp[j][k] - f * tmp[i][k]) % md;
        }
        rep(j, i + 1, n) A[i][j] = A[i][j] * v % md;
        rep(j, 0, n) tmp[i][j] = tmp[i][j] * v % md;
        A[i][i] = 1;
    }
    for (int i = n - 1; i > 0; --i) rep(j, 0, i) {
        ll v = A[j][i];
        rep(k, 0, n) tmp[j][k] = (tmp[j][k]-v*tmp[i][k]) % md;
    }
    rep(i, 0, n) rep(j, 0, n)
        A[col[i]][col[j]] = (tmp[i][j] % md + md) % md;
    return n;
}
\end{lstlisting}

\entry{General matching}{$O(n^3)$ \quad mem $n^2$}{%
Needs \code{matInv}. Maximum matching in a graph that is \textbf{not} bipartite ---
pairing people with people, where triangles and other odd cycles exist and
\code{Matching} gives wrong answers. Returns the matched pairs. Edges are undirected,
0-indexed.\\[0.3ex]
It is randomised (puts random numbers in a skew-symmetric matrix and uses its rank, the
Tutte matrix trick): wrong with probability about $n / 10^9$, i.e. never in practice.
$n \le 500$ is comfortable. For a \emph{weighted} general matching there is nothing
here; for small $n$ ($\le 20$) use bitmask DP over the set of unmatched vertices.
From KACTL.}

%DEP modbasics matInv
%LABEL generalMatching
\begin{lstlisting}
vector<pii> generalMatching(int N, const vector<pii>& ed) {
    const ll md = 1e9 + 7;
    mt19937 rng(2024);
    vector<vi> mat(N, vi(N)), A;
    for (auto [a, b] : ed) {
        ll r = rng() % md;
        mat[a][b] = r, mat[b][a] = (md - r) % md;
    }
    int r = matInv(A = mat, md), M = 2 * N - r, fi, fj;
    assert(r % 2 == 0);
    if (M != N) do {
        mat.resize(M, vi(M));
        rep(i, 0, N) {
            mat[i].resize(M);
            rep(j, N, M) {
                ll x = rng() % md;
                mat[i][j] = x, mat[j][i] = (md - x) % md;
            }
        }
    } while (matInv(A = mat, md) != M);
    vector<int> has(M, 1); vector<pii> ret;
    rep(it, 0, M / 2) {
        rep(i, 0, M) if (has[i])
            rep(j, i + 1, M) if (A[i][j] && mat[i][j]) {
                fi = i; fj = j; goto done;
            }
        assert(0); done:
        if (fj < N) ret.emplace_back(fi, fj);
        has[fi] = has[fj] = 0;
        rep(sw, 0, 2) {
            ll a = power(A[fi][fj], md - 2, md);
            rep(i, 0, M) if (has[i] && A[i][fj]) {
                ll b = A[i][fj] * a % md;
                rep(j, 0, M)
                    A[i][j] = (A[i][j] - A[fi][j] * b) % md;
            }
            swap(fi, fj);
        }
    }
    return ret;
}
\end{lstlisting}

\entry{Eulerian path}{$O(n + m)$ \quad mem $n + m$}{%
Hierholzer, iterative, handles both directions with the \code{dir} flag. Returns the
\textbf{vertex} sequence, length \code{m+1}, using every edge exactly once --- or an
empty vector when none exists (wrong degrees, or the edges are disconnected). A closed
tour comes back with \code{path.front() == path.back()}.\\[0.3ex]
Existence: \emph{directed} needs \code{out-in} equal to 0 everywhere except possibly one
$+1$ (the start) and one $-1$ (the end); \emph{undirected} needs zero or exactly two odd
vertices. Both also need all edges in one connected component --- which the final
\code{size() != m+1} test catches, so isolated vertices do not matter.\\[0.3ex]
For the edge sequence instead, record \code{id} alongside each pop.}

%LABEL eulerPath
\begin{lstlisting}
vector<int> eulerPath(int n, const vector<pii>& es,
                      bool dir) {
    int m = (int)es.size();
    if (!m) return {};
    vector<vector<pii>> g(n);          // (to, edge id)
    vi outd(n, 0), ind(n, 0);
    for (int i = 0; i < m; i++) {
        int a = (int)es[i].first, b = (int)es[i].second;
        g[a].push_back({b, i}); outd[a]++; ind[b]++;
        if (!dir) {
            g[b].push_back({a, i}); outd[b]++; ind[a]++;
        }
    }
    int s = -1;
    if (dir) {
        int plus = 0, minus = 0;
        for (int v = 0; v < n; v++) {
            ll df = outd[v] - ind[v];
            if (df == 1) { plus++; s = v; }
            else if (df == -1) minus++;
            else if (df) return {};
        }
        if (plus > 1 || minus > 1) return {};
    } else {
        int odd = 0;
        for (int v = 0; v < n; v++)
            if (outd[v] & 1) { odd++; s = v; }
        if (odd != 0 && odd != 2) return {};
    }
    if (s < 0)
        for (int v = 0; v < n; v++)
            if (outd[v]) { s = v; break; }
    vector<int> it(n, 0), st{s}, path;
    vector<char> used(m, 0);
    while (!st.empty()) {
        int u = st.back();
        while (it[u] < (int)g[u].size()
               && used[g[u][it[u]].second]) it[u]++;
        if (it[u] == (int)g[u].size()) {
            path.push_back(u); st.pop_back();
        } else {
            auto [v, id] = g[u][it[u]++];
            used[id] = 1; st.push_back((int)v);
        }
    }
    if ((int)path.size() != m + 1) return {};   // disjoint
    reverse(all(path));
    return path;
}
\end{lstlisting}

\section{Bitmask and DP}

\entry{PathDP (bitmask)}{$O(2^n n^2)$ \quad mem $2^n n$}{%
Folds every simple path covering \emph{exactly} a mask into one value, $n \le 20$.
Directed --- fill both directions for undirected. Pick a row of the table, edit the 5
config lines, done. The constructor does all the work; every call below is a
lookup.\\[0.3ex]
{\ttfamily\scriptsize
\begingroup\setlength{\tabcolsep}{2.5pt}
\begin{tabular}{@{}llllll@{}}
want & T & UNIT & NONE & ext(x,wt) & merge(a,b)\\
shortest path & ll & 0 & INF & x + wt & min(a, b)\\
longest path & ll & 0 & -INF & x + wt & max(a, b)\\
count paths & ll & 1 & 0 & x & a + b\\
count mod p & ll & 1 & 0 & x & (a + b) \% p\\
min bottleneck & ll & 0 & INF & max(x, wt) & min(a, b)\\
max probability & double & 1 & 0 & x * wt & max(a, b)\\
reachable lengths & bitset & 1<<0 & 0 & x << wt & a | b\\
\end{tabular}\endgroup}\\[0.3ex]
\begin{ctab}{0.25}{0.70}
\code{PathDP p(w,st,cn)} & \code{w[u][v]} = weight, \code{INF} = no edge.\\
\code{val(m,v)} & paths covering exactly \code{m}, ending at \code{v}.\\
\code{hamPath()} & over \code{v} of \code{val(full, v)}.\\
\code{hamPathTo(v)} & \code{val(full, v)}.\\
\code{hamCycle(s)} & closes \code{w[v][s]}; build with \code{st = 1LL<<s}.\\
\code{endingAt(v)} & over every mask, ending at \code{v}.\\
\code{bySize(k)} & over the masks of exactly \code{k} vertices.\\
\code{anyPath()} & over every mask and every \code{v}.\\
\code{cycles()} & over every simple cycle; build with \code{cn = true}.\\
\code{path(m,v)} & vertex order; \code{MIN} / \code{MAX} only.\\
\code{induced(g,vs)} & \code{w} for the vertices \code{vs}, weights all 1, local
\code{i} is \code{vs[i]}; translate back with \code{vs[x]}.\\
\end{ctab}\\[0.3ex]
\code{st} is a \textbf{mask}, not a vertex: \code{-1} = any start, \code{1LL<<s} pins
it. \code{cn} makes every path begin at the lowest vertex of its mask, for
\code{cycles()} only.\\[0.3ex]
\code{val(m,v)} is exactly \code{m}, not a subset, and \code{v} must be inside
\code{m}. Undirected counts come out $2\times$. \code{ext} must distribute over
\code{merge}, and \code{merge} must commute. Struct or bitset \code{T}:
\code{inline static const}, not \code{constexpr}. Unmodded counts overflow \code{ll} by
dense $n = 18$. At \code{-O2}, 30\% density: $n{=}16$ 0.01 s 13 MB, $n{=}18$ 0.05 s
44 MB, $n{=}20$ 0.22 s 184 MB.}

\begin{lstlisting}
struct PathDP {
    using T = ll;
    static T ext(T x, ll wt) { return x + wt; }
    static T merge(T a, T b) { return min(a, b); }
    static constexpr T UNIT = 0;
    static constexpr T NONE = 1e18;
    static constexpr ll INF = 1e18;      // missing edge in w

    int n; ll starts, full;
    matrix w; vector<T> dp; vector<char> par;

    PathDP(const matrix& wt, ll ss = -1, bool canon = false)
            : n(wt.size()), starts(ss), full((1LL << n) - 1),
              w(wt) {
        dp.assign((full + 1) * n, NONE);
        par.assign((full + 1) * n, -1);
        rep(v, 0, n) if (starts >> v & 1)
            dp[(1LL << v) * n + v] = UNIT;
        vector<ll> nb(n, 0);        // out-edges of u as bits
        rep(u, 0, n) rep(v, 0, n)
            if (w[u][v] < INF) nb[u] |= 1LL << v;
        rep(m, 1, full + 1) {
            ll base = m * n;
            for (ll b = m; b; b &= b - 1) {   // set bits of m
                int u = __builtin_ffsll(b) - 1;
                T cur = dp[base + u];
                if (cur == NONE) continue;
                for (ll c = nb[u] & ~m; c; c &= c - 1) {
                    int v = __builtin_ffsll(c) - 1;
                    if (canon && (1LL << v) < (m & -m))
                        continue;
                    ll t = (m | 1LL << v) * n + v;
                    T nv = merge(dp[t], ext(cur, w[u][v]));
                    if (nv != dp[t]) {
                        dp[t] = nv; par[t] = (char)u;
                    }
                }
            }
        }
    }
    T val(ll m, int v) { return dp[m * n + v]; }
    T hamPathTo(int v) { return val(full, v); }
    T hamPath() {
        T r = NONE;
        rep(v, 0, n) r = merge(r, val(full, v));
        return r;
    }
    T hamCycle(int s) {          // build with 1LL << s
        T r = NONE;
        rep(v, 0, n) {
            T d = val(full, v);
            if (d != NONE && w[v][s] < INF)
                r = merge(r, ext(d, w[v][s]));
        }
        return r;
    }
    T endingAt(int v) {
        T r = NONE;
        rep(m, 1, full + 1) r = merge(r, val(m, v));
        return r;
    }
    T bySize(int k) {            // exactly k vertices
        T r = NONE;
        rep(m, 1, full + 1) if (__builtin_popcountll(m) == k)
            rep(v, 0, n) r = merge(r, val(m, v));
        return r;
    }
    T anyPath() {
        T r = NONE;
        rep(m, 1, full + 1) rep(v, 0, n)
            r = merge(r, val(m, v));
        return r;
    }
    T cycles() {                 // build with canon
        T r = NONE;
        rep(m, 1, full + 1) {
            if (__builtin_popcountll(m) < 3) continue;
            int s = 0;
            while (!(m >> s & 1)) s++;   // lowest of m
            rep(v, 0, n) {
                T d = val(m, v);
                if (v != s && d != NONE && w[v][s] < INF)
                    r = merge(r, ext(d, w[v][s]));
            }
        }
        return r;
    }
    vi path(ll m, int v) {               // MIN / MAX only
        vi p;
        while (v >= 0) {
            p.push_back(v);
            int u = par[m * n + v]; m ^= 1LL << v; v = u;
        }
        reverse(p.begin(), p.end());
        return p;
    }
    template<class G, class V>    // g, vs: any integer type
    static matrix induced(const G& g, const V& vs) {
        static vi idx;        // reused, all -1 on entry
        if (idx.size() < g.size()) idx.assign(g.size(), -1);
        int k = (int)vs.size();
        rep(i, 0, k) idx[vs[i]] = i;
        matrix w(k, vi(k, INF));
        rep(i, 0, k) for (ll u : g[vs[i]])
            if (idx[u] >= 0) w[i][idx[u]] = 1;
        rep(i, 0, k) idx[vs[i]] = -1;   // wipe what we set
        return w;
    }
};
\end{lstlisting}

\entry{Subset sums (SOS / zeta)}{$O(2^n n)$ \quad mem $2^n$}{%
\code{sos(n, a)[m]} is the sum of \code{a[s]} over \emph{all} submasks \code{s} of
\code{m}. Replace \code{+=} with \code{|=} / \code{max} for ``or over submasks'' /
``best submask''. Supersets instead: flip the condition to
\code{if (!(m >> b \& 1)) f[m] += f[m | 1LL<<b]}.\\[0.3ex]
Inverse (M\"obius, recovers \code{a} from \code{f}): same loops with \code{-=}.
Subset convolution and ``count pairs with \code{a \& b == 0}'' both reduce to
this.\\[0.3ex]
To walk the submasks of one mask only: \code{for (ll s = m; s; s = (s-1) \& m)}, which
visits every nonempty submask once --- $O(3^n)$ over all masks.}

\begin{lstlisting}
// f has size 1<<n and holds a on entry
vi sos(int n, vi f) {
    rep(b, 0, n) rep(m, 0, 1LL << n)
        if (m >> b & 1) f[m] += f[m ^ (1LL << b)];
    return f;
}
\end{lstlisting}

\entry{LIS}{$O(n \log n)$ \quad mem $n$}{%
\code{c[k]} ends up as the smallest possible tail of an increasing subsequence of
length \code{k+1}, so \code{c.size()} is the answer. \code{c} itself is \emph{not} a
valid subsequence.\\[0.3ex]
\code{lower\_bound} gives \textbf{strictly} increasing; \code{upper\_bound} gives
non-decreasing. Longest decreasing: reverse the input, or negate it. Minimum number of
non-increasing chains covering the sequence = length of the LIS (Dilworth).\\[0.3ex]
To reconstruct, store for each element the index it was placed at plus the predecessor
\code{c}-index, then walk back from the last placement.}

\begin{lstlisting}
ll lis(const vi& a) {
    vi c;
    for (ll e : a) {
        auto it = lower_bound(all(c), e);
        if (it == c.end()) c.push_back(e);
        else *it = e;
    }
    return (ll)c.size();
}
\end{lstlisting}

\entry{Interval DP}{$O(n^3)$, optimised $O(n^2)$ \quad mem $n^2$}{%
Shortest interval first, so both halves are ready when you need them. Circular
versions: double the array to length $2n$ and take the best window of length
$n$.\\[0.3ex]
\textbf{Knuth} cuts it to $O(n^2)$ when \code{opt[i][j-1] <= opt[i][j] <= opt[i+1][j]}
--- restrict \code{k} to that window and keep an \code{opt} table. It holds when the
cost satisfies the quadrangle inequality, which is true for the classic ``merge
adjacent piles, cost = sum of the pile'' problem. \textbf{Divide and conquer
optimisation} is the same idea one layer at a time, for
\code{dp[i][j] = min over k of dp[i-1][k] + C(k,j)} with monotone \code{opt}:
recurse on \code{(jlo,jhi,klo,khi)}, evaluate the midpoint, split. Both are worth it
only when $n^3$ is too slow --- check the bound first.}

%NOCOMPILE
\begin{lstlisting}
rep(len, 2, n + 1)
  rep(i, 0, n - len + 1) {
    ll j = i + len - 1;                  // inclusive
    dp[i][j] = INF;
    rep(k, i, j)                         // split point
      dp[i][j] = min(dp[i][j],
                     dp[i][k] + dp[k+1][j] + cost(i, j));
  }
\end{lstlisting}

\entry{Divide and conquer DP}{$O(k \cdot n \log n)$}{%
For ``split the array into $k$ groups, the cost of a group is \code{C(l, r)}, minimise
the total'': $dp_j[i] = \min_{k < i} dp_{j-1}[k] + C(k, i)$. Normally $O(k n^2)$. If the
best split point never moves left as $i$ grows (true when \code{C} satisfies
$C(a,c) + C(b,d) \le C(a,d) + C(b,c)$ for $a \le b \le c \le d$, e.g. ``sum squared'' or
``number of equal pairs'' in a group), solve the middle $i$ first, and its best $k$
splits the search range for both halves.\\[0.3ex]
Run one layer at a time: \code{prv} is layer $j-1$ (layer 0 is \code{0} at index 0 and
\code{INF} elsewhere), call \code{dnc(1, n+1, 0, n, prv, cur, C)}, then
\code{prv = cur}. Answer is \code{prv[n]} after $k$ layers. \code{C(k, i)} is the cost
of elements \code{[k, i)}, usually $O(1)$ from prefix sums. Not sure the condition
holds? test against the $O(n^2)$ version on random input.}

%LABEL dnc
\begin{lstlisting}
template<class F>   // cur[i] = min_{k<i} prv[k] + C(k, i)
void dnc(int l, int r, int optl, int optr,
         const vi& prv, vi& cur, F& C) {
    if (l >= r) return;
    int m = (l + r) / 2, best = optl;
    ll bv = INF;
    for (int k = optl; k <= min(m - 1, optr); k++) {
        if (prv[k] >= INF) continue;
        ll v = prv[k] + C(k, m);
        if (v < bv) bv = v, best = k;
    }
    cur[m] = bv;
    dnc(l, m, optl, best, prv, cur, C);
    dnc(m + 1, r, best, optr, prv, cur, C);
}
\end{lstlisting}

\entry{Knuth optimisation}{$O(n^2)$ instead of $O(n^3)$}{%
Interval DP of the form $dp[i][j] = \min_{i<k<j} dp[i][k] + dp[k][j] + w(i, j)$, like
``cut a stick at given points, each cut costs the current length'' or ``merge adjacent
piles, a merge costs the total size''. If the best $k$ satisfies
$opt[i][j{-}1] \le opt[i][j] \le opt[i{+}1][j]$ (true when $w$ satisfies the condition
from D\&C DP above and $w(b,c) \le w(a,d)$ for $a \le b \le c \le d$), searching only
between those makes the whole table $O(n^2)$.\\[0.3ex]
Points are $0..n$, \code{dp[i][i+1] = 0}, answer \code{dp[0][n]}. \code{w(i, j)} is the
cost of combining the block from point $i$ to point $j$.}

%LABEL knuth
\begin{lstlisting}
template<class F> vector<vi> knuth(int n, F w) {
    vector<vi> dp(n + 1, vi(n + 1, 0));
    vector<vector<int>> opt(n + 1, vector<int>(n + 1));
    rep(i, 0, n) opt[i][i + 1] = (int)i + 1;
    for (int len = 2; len <= n; len++)
        for (int i = 0; i + len <= n; i++) {
            int j = i + len; dp[i][j] = INF;
            int lo = max(i + 1, opt[i][j-1]);
            int hi = min(j - 1, opt[i+1][j]);
            for (int k = lo; k <= hi; k++) {
                ll v = dp[i][k] + dp[k][j] + w(i, j);
                if (v < dp[i][j]) dp[i][j] = v, opt[i][j] = k;
            }
        }
    return dp;
}
\end{lstlisting}

\section{Linear algebra and numerics}

\entry{Gauss mod p}{$O(n^2 m)$ \quad mem $nm$}{%
Needs \code{power}. Solves $Ax \equiv b \pmod p$, $p$ prime, $A$ is $n \times m$ ($n$
equations, $m$ unknowns). Returns \code{-1} if there is no solution, else the rank, and
\code{x} is one solution with free variables 0. Rank $< m$ means $p^{m - \text{rank}}$
solutions in total --- that is how ``count the solutions'' questions are answered.
Exact, so no \code{eps}. Ruiming's \code{mod\_elimination}, rewritten 0-indexed with
the inverse taken by \code{power}.}

%DEP modbasics
%LABEL gaussMod
\begin{lstlisting}
int gaussMod(vector<vi> a, vi b, vi& x, ll p) {
    int n = (int)a.size(), m = (int)a[0].size(), rank = 0;
    for (auto& r : a) for (ll& v : r) v = (v % p + p) % p;
    for (ll& v : b) v = (v % p + p) % p;
    vector<int> where(m, -1);
    for (int col = 0; col < m && rank < n; col++) {
        int sel = -1;
        for (int i = rank; i < n; i++)
            if (a[i][col]) { sel = i; break; }
        if (sel < 0) continue;
        swap(a[sel], a[rank]); swap(b[sel], b[rank]);
        where[col] = rank;
        ll iv = power(a[rank][col], p - 2, p);
        for (int j = col; j < m; j++)
            a[rank][j] = a[rank][j] * iv % p;
        b[rank] = b[rank] * iv % p;
        for (int i = 0; i < n; i++)
            if (i != rank && a[i][col]) {
                ll c = a[i][col];
                for (int j = col; j < m; j++)
                    a[i][j] = ((a[i][j] - c * a[rank][j]) % p
                               + p) % p;
                b[i] = ((b[i] - c * b[rank]) % p + p) % p;
            }
        rank++;
    }
    for (int i = rank; i < n; i++) if (b[i]) return -1;
    x.assign(m, 0);
    for (int j = 0; j < m; j++)
        if (where[j] >= 0) x[j] = b[where[j]];
    return rank;
}
\end{lstlisting}

\entry{Gauss mod 2 (XOR equations)}{$O(n^2 m / 64)$}{%
The same thing over bits: each equation is ``the XOR of some variables equals 0 or 1''.
Rows are \code{bitset}s, so $n = m = 2000$ is instant. Lights-out puzzles, ``flip
switches so every lamp is off'', parity constraints. Fill \code{A[i][j] = 1} when
variable \code{j} appears in equation \code{i}, and \code{b[i]} is the right side.
Returns \code{-1} if impossible, else the rank; $2^{m - \text{rank}}$ solutions exist.
Destroys \code{A} and \code{b}. Bits at columns $\ge m$ must be 0. KACTL's, replacing
Ruiming's \code{xor\_gauss} (same algorithm, general interface).}

%LABEL gaussXor
\begin{lstlisting}
typedef bitset<2048> bs;
int gaussXor(vector<bs>& A, vector<int>& b, bs& x, int m) {
    int n = (int)A.size(), rank = 0, br;
    vector<int> col(m); iota(all(col), 0);
    rep(i, 0, n) {
        for (br = i; br < n; ++br) if (A[br].any()) break;
        if (br == n) {
            rep(j, i, n) if (b[j]) return -1;
            break;
        }
        int bc = (int)A[br]._Find_next(i - 1);
        swap(A[i], A[br]); swap(b[i], b[br]);
        swap(col[i], col[bc]);
        rep(j, 0, n) if (A[j][i] != A[j][bc]) {
            A[j].flip(i); A[j].flip(bc);
        }
        rep(j, i + 1, n) if (A[j][i]) {
            b[j] ^= b[i]; A[j] ^= A[i];
        }
        rank++;
    }
    x = bs();
    for (int i = rank; i--;) {
        if (!b[i]) continue;
        x[col[i]] = 1;
        rep(j, 0, i) b[j] ^= A[j][i];
    }
    return rank;
}
\end{lstlisting}

\entry{XOR basis}{add $O(\log C)$ \quad mem 63 words}{%
Keeps a small set of numbers (at most 63) whose XOR combinations are exactly the XORs
of all subsets of everything added so far. Answers: ``can $x$ be made as a XOR of some
of the numbers?'' (\code{can}), ``largest XOR of a subset'' (\code{maxXor()}), ``largest
\code{x \^{} subset}'' (\code{maxXor(x)}), ``$k$-th smallest reachable value''
(\code{kth}, counting 0 as the 0-th).\\[0.3ex]
\code{add} returns \code{false} when the number was already reachable: then the empty
subset is not the only way to get 0, which matters if the problem forbids empty subsets
(if \emph{every} add returned \code{true}, 0 is not reachable by a non-empty subset, so
shift \code{k} by one). Count of distinct reachable values: $2^{\text{rank}}$. Works
for numbers $< 2^{63}$; combine with a DFS for ``max XOR path in a graph'' (add every
cycle's XOR). Ruiming's \code{xor\_basis}, as a struct.}

%LABEL xorbasis
\begin{lstlisting}
struct XorBasis {
    static const int B = 63;
    ll b[B] = {}; int rank = 0;
    bool add(ll x) {
        for (int i = B - 1; i >= 0; i--) if (x >> i & 1) {
            if (!b[i]) { b[i] = x; rank++; return true; }
            x ^= b[i];
        }
        return false;                    // already reachable
    }
    bool can(ll x) {
        for (int i = B - 1; i >= 0; i--) if (x >> i & 1) {
            if (!b[i]) return false;
            x ^= b[i];
        }
        return true;
    }
    ll maxXor(ll x = 0) {
        for (int i = B - 1; i >= 0; i--) x = max(x, x ^ b[i]);
        return x;
    }
    ll kth(ll k) {                // 0-th = 0; -1 if too big
        vector<ll> v;
        for (int i = 0; i < B; i++) if (b[i]) {
            for (int j = i - 1; j >= 0; j--)  // reduce
                if (b[i] >> j & 1) b[i] ^= b[j];
            v.push_back(b[i]);
        }
        if (v.size() < 63 && (k >> v.size())) return -1;
        ll r = 0;
        rep(i, 0, (ll)v.size()) if (k >> i & 1) r ^= v[i];
        return r;
    }
};
\end{lstlisting}

\entry{Determinant}{$O(n^3)$}{%
\code{det} for real matrices (partial pivoting); \code{detMod} is exact modulo
\emph{any} \code{md}, prime or not, because it only uses Euclid-style row subtraction
(no division). Reduce the input mod \code{md} first.\\[0.3ex]
\textbf{Count spanning trees} (Kirchhoff): build the Laplacian $L$ ($L_{ii}$ = degree,
$L_{ij} = -$(number of edges $i$--$j$)), delete row 0 and column 0, take the
determinant. Directed, trees pointing into root $r$: $L_{ii}$ = out-degree, delete row
and column $r$. Also: $\det = 0$ iff the rows are dependent. From KACTL.}

%LABEL det
\begin{lstlisting}
ld det(vector<vector<ld>> a) {
    int n = (int)a.size(); ld res = 1;
    rep(i, 0, n) {
        int b = i;
        rep(j, i + 1, n)
            if (fabsl(a[j][i]) > fabsl(a[b][i])) b = j;
        if (i != b) swap(a[i], a[b]), res *= -1;
        res *= a[i][i];
        if (res == 0) return 0;
        rep(j, i + 1, n) {
            ld v = a[j][i] / a[i][i];
            if (v != 0) rep(k, i + 1, n) a[j][k] -= v*a[i][k];
        }
    }
    return res;
}
ll detMod(vector<vi> a, ll md) {
    int n = (int)a.size(); ll ans = 1;
    rep(i, 0, n) {
        rep(j, i + 1, n) while (a[j][i] != 0) {   // gcd step
            ll t = a[i][i] / a[j][i];
            if (t) rep(k, i, n)
                a[i][k] = (a[i][k] - a[j][k] * t) % md;
            swap(a[i], a[j]);
            ans *= -1;
        }
        ans = ans * a[i][i] % md;
        if (!ans) return 0;
    }
    return (ans + md) % md;
}
\end{lstlisting}

\entry{Berlekamp-Massey + linear recurrence}{$O(n^2)$, then $O(n^2 \log k)$}{%
Needs \code{power}. \textbf{The ``guess the formula'' tool.} If the answer for $n$
follows \emph{some} linear recurrence $a_i = c_1 a_{i-1} + \dots + c_d a_{i-d}$ (mod a
prime) --- which is true for every ``count paths / tilings / strings with a fixed
small automaton'' problem and anything a matrix power could solve --- then:\\[0.3ex]
1. Brute force the first $2d + 10$ terms (more is safer), mod the prime.\\
2. \code{tr = berlekampMassey(terms, md)} returns $c_1, \dots, c_d$.\\
3. \code{linearRec(terms, tr, k, md)} gives term $k$ (0-indexed), $k$ up to $10^{18}$.\\[0.3ex]
Much faster and simpler than building the matrix by hand: you never have to find the
recurrence yourself. If \code{tr} comes out as long as half your input, you did not
give enough terms or there is no linear recurrence. Example:
\code{berlekampMassey(\{0,1,1,3,5,11\}, md)} is \code{\{1, 2\}}. From KACTL.}

%DEP modbasics
%LABEL bm
\begin{lstlisting}
vi berlekampMassey(vi s, ll md) {
    int n = (int)s.size(), L = 0, m = 0;
    vi C(n), B(n), T;
    C[0] = B[0] = 1;
    ll b = 1;
    rep(i, 0, n) { ++m;
        ll d = s[i] % md;
        rep(j, 1, L + 1) d = (d + C[j] * s[i - j]) % md;
        if (!d) continue;
        T = C; ll coef = d * power(b, md - 2, md) % md;
        rep(j, m, n) C[j] = (C[j] - coef * B[j - m]) % md;
        if (2 * L > i) continue;
        L = i + 1 - L; B = T; b = d; m = 0;
    }
    C.resize(L + 1); C.erase(C.begin());
    for (ll& x : C) x = (md - x) % md;
    return C;
}
ll linearRec(vi S, vi tr, ll k, ll md) {   // S[k]
    int n = (int)tr.size();
    auto combine = [&](vi a, vi b) {
        vi res(n * 2 + 1);
        rep(i, 0, n + 1) rep(j, 0, n + 1)
            res[i + j] = (res[i + j] + a[i] * b[j]) % md;
        for (int i = 2 * n; i > n; --i) rep(j, 0, n)
            res[i-1-j] = (res[i-1-j] + res[i] * tr[j]) % md;
        res.resize(n + 1);
        return res;
    };
    vi pol(n + 1), e(pol);
    pol[0] = e[1] = 1;
    for (++k; k; k /= 2) {
        if (k % 2) pol = combine(pol, e);
        e = combine(e, e);
    }
    ll res = 0;
    rep(i, 0, n) res = (res + pol[i + 1] * S[i]) % md;
    return res;
}
\end{lstlisting}

\entry{Linear programming (simplex)}{fast in practice}{%
Maximise $c^T x$ subject to $Ax \le b$ and $x \ge 0$. Returns the best value, or
\code{-inf} if the rules contradict each other, or \code{inf} if it can grow forever;
\code{x} gets an optimal point.\\[0.3ex]
Turning a problem into this form: $\ge$ constraints are $\le$ with both sides negated;
an equality is one $\le$ and one $\ge$; minimise $f$ = maximise $-f$; a variable that
may be negative is $x = x^+ - x^-$ with two non-negative variables. Use it when the
problem is ``choose real amounts, every rule is a linear inequality, optimise a linear
total'' --- e.g. diet / mixing / scheduling with fractional amounts. Answers are
\code{double}: print with enough digits. Integer answers are a different (hard)
problem. From KACTL.\\[0.3ex]
\code{LPSolver::vd b = \{1, 1, -4\}, c = \{-1, -1\}, x;}\\
\code{LPSolver(A, b, c).solve(x);} with \code{A = \{\{1,-1\},\{-1,1\},\{-1,-2\}\}}.}

%LABEL simplex
\begin{lstlisting}
#define MP make_pair
#define ltj(X) if (s==-1 || MP(X[j],N[j]) < MP(X[s],N[s])) s=j
struct LPSolver {
    using T = double; using vd = vector<T>;
    static constexpr T eps = 1e-8;
    static constexpr T inf = numeric_limits<T>::infinity();
    int m, n; vector<int> N, B; vector<vd> D;
    LPSolver(const vector<vd>& A, const vd& b, const vd& c)
        : m((int)b.size()), n((int)c.size()), N(n + 1),
          B(m), D(m + 2, vd(n + 2)) {
        rep(i, 0, m) rep(j, 0, n) D[i][j] = A[i][j];
        rep(i, 0, m) {
            B[i] = n + i; D[i][n] = -1; D[i][n+1] = b[i];
        }
        rep(j, 0, n) { N[j] = j; D[m][j] = -c[j]; }
        N[n] = -1; D[m+1][n] = 1;
    }
    void pivot(int r, int s) {
        T *a = D[r].data(), inv = 1 / a[s];
        rep(i, 0, m + 2) if (i != r && abs(D[i][s]) > eps) {
            T *b = D[i].data(), inv2 = b[s] * inv;
            rep(j, 0, n + 2) b[j] -= a[j] * inv2;
            b[s] = a[s] * inv2;
        }
        rep(j, 0, n + 2) if (j != s) D[r][j] *= inv;
        rep(i, 0, m + 2) if (i != r) D[i][s] *= -inv;
        D[r][s] = inv;
        swap(B[r], N[s]);
    }
    bool simplex(int phase) {
        int x = m + phase - 1;
        for (;;) {
            int s = -1;
            rep(j, 0, n + 1) if (N[j] != -phase) ltj(D[x]);
            if (D[x][s] >= -eps) return true;
            int r = -1;
            rep(i, 0, m) {
                if (D[i][s] <= eps) continue;
                if (r == -1 || MP(D[i][n+1] / D[i][s], B[i])
                             < MP(D[r][n+1] / D[r][s], B[r]))
                    r = i;
            }
            if (r == -1) return false;
            pivot(r, s);
        }
    }
    T solve(vd& x) {
        int r = 0;
        rep(i, 1, m) if (D[i][n+1] < D[r][n+1]) r = i;
        if (D[r][n+1] < -eps) {
            pivot(r, n);
            if (!simplex(2) || D[m+1][n+1] < -eps)
                return -inf;
            rep(i, 0, m) if (B[i] == -1) {
                int s = 0;
                rep(j, 1, n + 1) ltj(D[i]);
                pivot(i, s);
            }
        }
        bool ok = simplex(1); x = vd(n);
        rep(i, 0, m) if (B[i] < n) x[B[i]] = D[i][n+1];
        return ok ? D[m][n+1] : inf;
    }
};
\end{lstlisting}

\entry{AND / OR / XOR convolution}{$O(n \log n)$, $n = 2^k$}{%
\code{c[z]} = sum of \code{a[x]*b[y]} over all pairs with \code{x OP y == z}, where OP
is AND (\code{mode 0}), OR (\code{1}) or XOR (\code{2}). The same job as \code{conv} for
$+$, but on bit masks. Pad both arrays with zeros to the same power of two.\\[0.3ex]
Uses: ``number of pairs whose XOR is $z$'' for all $z$ at once; combining two
distributions over subsets; Nim-like games (XOR of piles) summed over choices.
\code{FST} on its own: OR forward = sum over subsets (same as \code{sos}), AND forward
= sum over supersets. Values grow like $n \cdot \max^2$; to work mod a prime, add
\code{\% md} to every line and replace the final \code{/ n} by multiplying with
\code{power(n, md-2, md)}. b5xu's \code{bitwise-xor-convolution}, fixed (the original
does not compile) and extended with KACTL's AND/OR.}

%LABEL fst
\begin{lstlisting}
void FST(vi& a, bool inv, int mode) {    // 0 AND 1 OR 2 XOR
    int n = (int)a.size();
    for (int s = 1; s < n; s *= 2)
        for (int i = 0; i < n; i += 2 * s) rep(j, i, i + s) {
            ll &u = a[j], &v = a[j + s];
            if (mode == 0) u += inv ? -v : v;
            else if (mode == 1) v += inv ? -u : u;
            else tie(u, v) = pii(u + v, u - v);
        }
    if (inv && mode == 2) for (ll& x : a) x /= n;
}
vi bitConv(vi a, vi b, int mode) {     // same size, 2^k
    FST(a, 0, mode); FST(b, 0, mode);
    rep(i, 0, (ll)a.size()) a[i] *= b[i];
    FST(a, 1, mode);
    return a;
}
\end{lstlisting}

\entry{Polynomial interpolation}{$O(n^2)$ / $O(n \log p)$}{%
Any $n$ points with different $x$ lie on exactly one polynomial of degree $< n$.
\code{interpolate} returns its coefficients (\code{double}; inaccurate for large $n$).
\code{lagrange} is the exact, more useful one: given $y_i = p(i)$ for
$i = 0..n-1$ mod a prime $\ge n$, it returns $p(X)$ for a huge $X$.\\[0.3ex]
Typical use: ``the answer is a polynomial in $n$ of degree $d$'' --- e.g.
$\sum_{i \le n} i^k$ has degree $k+1$, and many counting answers over a growing box are
polynomials. Brute force $d+1$ (or a few more) small values, then \code{lagrange} at
the real $n$. Linear recurrence instead of polynomial? use Berlekamp-Massey.}

%DEP modbasics
%LABEL interp
\begin{lstlisting}
vector<ld> interpolate(vector<ld> x, vector<ld> y) {
    int n = (int)x.size();
    vector<ld> res(n), temp(n);
    rep(k, 0, n - 1) rep(i, k + 1, n)
        y[i] = (y[i] - y[k]) / (x[i] - x[k]);
    ld last = 0; temp[0] = 1;
    rep(k, 0, n) rep(i, 0, n) {
        res[i] += y[k] * temp[i];
        swap(last, temp[i]);
        temp[i] -= last * x[k];
    }
    return res;                      // res[i] * x^i
}
ll lagrange(const vi& y, ll X, ll md) {  // y[i] = p(i)
    int n = (int)y.size(); X %= md;
    vi pre(n + 1, 1), suf(n + 1, 1), f(n, 1);
    rep(i, 0, n) pre[i+1] = pre[i] * ((X - i + md) % md) % md;
    down(i, n, 0) suf[i] = suf[i+1] * ((X - i + md) % md)%md;
    rep(i, 1, n) f[i] = f[i-1] * i % md;
    ll res = 0;
    rep(i, 0, n) {
        ll t = y[i] % md * (pre[i] * suf[i+1] % md) % md
             * power(f[i] * f[n-1-i] % md, md - 2, md) % md;
        res = ((n - 1 - i) & 1) ? (res - t + md) % md
                                : (res + t) % md;
    }
    return res;
}
\end{lstlisting}

\entry{Numerical integration}{adaptive Simpson}{%
\code{quad(a, b, f)} integrates \code{f} from \code{a} to \code{b}. It splits the
interval in halves only where the function is hard, so it is fast for smooth functions
and still fine for ones with a kink. Area of a weird shape: integrate its height
(e.g. the length of a vertical cut through a union of circles) over $x$. Volume:
nest it, \code{quad(.., [\&](ld x)\{ return quad(.., [\&](ld y)\{ ... \}); \})}.\\[0.3ex]
Trouble with jumps (a function that is 0 almost everywhere): split at the known jump
points yourself and call it per piece. From KACTL.}

%LABEL quad
\begin{lstlisting}
template<class F> ld simp(F& f, ld a, ld b) {
    return (f(a) + 4 * f((a + b) / 2) + f(b)) * (b - a) / 6;
}
template<class F>
ld simpRec(F& f, ld a, ld b, ld eps, ld S) {
    ld c = (a + b) / 2;
    ld S1 = simp(f, a, c), S2 = simp(f, c, b), T = S1 + S2;
    if (fabsl(T - S) <= 15 * eps || b - a < 1e-10)
        return T + (T - S) / 15;
    return simpRec(f, a, c, eps / 2, S1)
         + simpRec(f, c, b, eps / 2, S2);
}
template<class F> ld quad(ld a, ld b, F f, ld eps = 1e-8) {
    return simpRec(f, a, b, eps, simp(f, a, b));
}
\end{lstlisting}

\section{Number theory}

\entry{Modular basics}{$O(\log)$ each}{%
\code{power(a, b, m)} is $a^b \bmod m$; it needs $m < 2^{31}$ or the \code{a*a} inside
overflows --- for bigger moduli cast to \code{i128}.\\[0.3ex]
\code{ext\_gcd(a, b, x, y)} returns $g=\gcd$ and fills $ax+by=g$. \code{inv(a, m)}
gives $a^{-1} \bmod m$ for \emph{any} \code{m} with $\gcd(a,m)=1$, and \code{-1} when
none exists. When \code{m} is prime, \code{power(a, m-2, m)} is the same thing and
shorter.\\[0.3ex]
$a^b \bmod m$ with a \emph{gigantic} $b$: read $b$ as a string, reduce it mod
$\varphi(m)$ --- but only add $\varphi(m)$ back (i.e. use $b \bmod \varphi(m) +
\varphi(m)$) when the true $b \ge \varphi(m)$, otherwise you get the wrong answer for
small $b$. That correction is needed exactly when $\gcd(a,m) \ne 1$.}

%LABEL modbasics
\begin{lstlisting}
ll power(ll a, ll b, ll mod) {         // a^b % mod
    a %= mod; ll res = 1;
    for (; b > 0; b >>= 1, a = a * a % mod)
        if (b & 1) res = res * a % mod;
    return res;
}
ll ext_gcd(ll a, ll b, ll &x, ll &y) {     // a*x + b*y = g
    if (b == 0) { x = 1; y = 0; return a; }
    ll x1, y1;
    ll g = ext_gcd(b, a % b, x1, y1);
    x = y1;
    y = x1 - (a / b) * y1;
    return g;
}
ll inv(ll a, ll m) {                 // -1 if none exists
    ll x, y;
    if (ext_gcd(a, m, x, y) != 1) return -1;
    return ((x % m) + m) % m;
}
\end{lstlisting}

\entry{CRT}{$O(\log)$ \quad }{%
Merges $x \equiv r_1 \pmod{m_1}$ and $x \equiv r_2 \pmod{m_2}$ into one congruence
modulo $\mathrm{lcm}(m_1,m_2)$. Moduli need \emph{not} be coprime; returns
\code{\{-1,-1\}} when the two are incompatible. Fold a list by merging one at a
time.\\[0.3ex]
Watch the size: the intermediate \code{m1 / g * m2} must fit, so with several moduli
near $10^9$ switch to \code{i128}.}

%LABEL crt
%DEP modbasics
\begin{lstlisting}
// x = r1 (mod m1), x = r2 (mod m2)  ->  {r, lcm} or {-1,-1}
pair<ll,ll> crt(ll r1, ll m1, ll r2, ll m2) {
    ll x, y, g = ext_gcd(m1, m2, x, y);
    if ((r2 - r1) % g) return {-1, -1};
    ll l = m1 / g * m2, t = m2 / g;
    ll k = ((r2 - r1) / g % t) * (x % t) % t;
    ll r = (r1 + (i128)m1 * k) % l;
    return {(r % l + l) % l, l};
}
\end{lstlisting}

\entry{Sieves}{$O(n \log\log n)$ / $O(n)$ \quad mem $n$}{%
\code{sieve(n)} lists the primes up to \code{n} inclusive. \code{spf(n)} gives the
smallest prime factor of every number, which factorises any \code{x <= n} in
$O(\log x)$ by dividing repeatedly --- that is the one to build when you must factorise
many numbers. \code{phis(n)} computes Euler's $\varphi$ for everything at once by the
linear sieve.\\[0.3ex]
$\varphi(n) = n \prod_{p \mid n} (1 - 1/p)$; $\sum_{d \mid n} \varphi(d) = n$. Number of
divisors and their sum are multiplicative, so the same linear sieve shape computes
them.}

%LABEL sieves
\begin{lstlisting}
vector<int> sieve(int n) {
    vector<char> comp(n + 1, 0); vector<int> ps;
    for (int i = 2; i <= n; i++) if (!comp[i]) {
        ps.push_back(i);
        for (ll j = (ll)i * i; j <= n; j += i) comp[j] = 1;
    }
    return ps;
}
vector<int> spf(int n) {         // smallest prime factor
    vector<int> s(n + 1);
    iota(all(s), 0);
    for (int i = 2; (ll)i * i <= n; i++) if (s[i] == i)
        for (int j = i * i; j <= n; j += i)
            if (s[j] == j) s[j] = i;
    return s;
}
vi phis(int n) {                      // linear sieve
    vi phi(n + 1); vector<int> ps;
    vector<char> comp(n + 1, 0);
    phi[1] = 1;
    for (int i = 2; i <= n; i++) {
        if (!comp[i]) { ps.push_back(i); phi[i] = i - 1; }
        for (int p : ps) {
            if ((ll)p * i > n) break;
            comp[p * i] = 1;
            if (i % p == 0) { phi[p*i] = p * phi[i]; break; }
            phi[p * i] = (p - 1) * phi[i];
        }
    }
    return phi;
}
\end{lstlisting}

\entry{Miller-Rabin + Pollard rho}{$O(\log^2 n)$ / $O(n^{1/4})$}{%
\code{isPrime} is \textbf{deterministic} for every \code{n} that fits in
\code{ll}: those 12 bases are proven sufficient below $3.3\cdot10^{24}$.
\code{factor(n)} returns the prime factors with multiplicity, unsorted --- sort if you
want them grouped. Both use \code{i128} for the multiplications, so \code{n} up to
$9\cdot10^{18}$ is safe.\\[0.3ex]
Use trial division up to $\sqrt n$ instead when $n \le 10^{12}$ or so; it is shorter and
you will not mistype it under pressure.}

%LABEL pollard
\begin{lstlisting}
ll mulm(ll a, ll b, ll m) { return (i128)a * b % m; }
ll powm(ll a, ll b, ll m) {
    ll r = 1; a %= m;
    for (; b; b >>= 1, a = mulm(a, a, m))
        if (b & 1) r = mulm(r, a, m);
    return r;
}
bool isPrime(ll n) {
    if (n < 2) return false;
    for (ll p : {2,3,5,7,11,13,17,19,23,29,31,37})
        if (n % p == 0) return n == p;
    ll d = n - 1; int s = 0;
    while (!(d & 1)) d >>= 1, s++;
    for (ll a : {2,3,5,7,11,13,17,19,23,29,31,37}) {
        ll x = powm(a, d, n);
        if (x == 1 || x == n - 1) continue;
        bool ok = false;
        rep(i, 1, s) {
            x = mulm(x, x, n);
            if (x == n - 1) { ok = true; break; }
        }
        if (!ok) return false;
    }
    return true;
}
ll pollard(ll n) {                 // a nontrivial factor
    if (!(n & 1)) return 2;
    static mt19937_64 rr(88172645463325252ULL);
    while (true) {
        ll x = rr() % n, c = rr() % n + 1, y = x, d = 1;
        auto f = [&](ll v){ return (mulm(v,v,n)+c) % n; };
        while (d == 1) {
            x = f(x); y = f(f(y));
            if (x == y) break;
            d = __gcd(x > y ? x - y : y - x, n);
        }
        if (d != 1 && d != n) return d;
    }
}
void factor(ll n, vi& out) {
    if (n == 1) return;
    if (isPrime(n)) { out.push_back(n); return; }
    ll d = pollard(n);
    factor(d, out); factor(n / d, out);
}
\end{lstlisting}

\entry{Discrete log (baby-step giant-step)}{$O(\sqrt m)$}{%
Smallest $x > 0$ with $a^x \equiv b \pmod m$, or \code{-1}. Works even when
$\gcd(a, m) \ne 1$. Check $x = 0$ (i.e. $b = 1$) yourself if it is allowed.
\code{modLog(a, 1, m)} is the \emph{order} of $a$: the length of the cycle of powers.\\[0.3ex]
Idea: write $x = n i - j$ with $n \approx \sqrt m$. Store $b a^j$ for all small $j$ in a
hash map (baby steps), then try $a^{n i}$ for $i = 1, 2, \dots$ (giant steps) until one
is in the map. $m$ up to $\sim 3\cdot10^9$ (products must fit in \code{ll}; above that
use \code{i128}). From KACTL.}

%LABEL modLog
\begin{lstlisting}
ll modLog(ll a, ll b, ll m) {
    a %= m; b %= m;
    ll n = (ll)sqrtl(m) + 1, e = 1, f = 1, j = 1;
    unordered_map<ll, ll> A;
    while (j <= n && (e = f = e * a % m) != b)
        A[e * b % m] = j++;
    if (e == b) return j;
    if (__gcd(m, e) == __gcd(m, b))
        rep(i, 2, n + 2) if (A.count(e = e * f % m))
            return n * i - A[e];
    return -1;
}
\end{lstlisting}

\entry{Modular square root (Tonelli-Shanks)}{$O(\log^2 p)$}{%
Needs \code{power}. For a \textbf{prime} $p$, finds $x$ with $x^2 \equiv a \pmod p$; the
other root is $p - x$. A root exists iff $a \equiv 0$ or
\code{power(a, (p-1)/2, p) == 1} (Euler's criterion) --- test that first, the code
asserts it. Used for solving $x^2 + bx + c \equiv 0$ (complete the square) and in some
NTT tricks. From KACTL.}

%DEP modbasics
%LABEL modSqrt
\begin{lstlisting}
ll modSqrt(ll a, ll p) {
    a %= p; if (a < 0) a += p;
    if (a == 0) return 0;
    assert(power(a, (p - 1) / 2, p) == 1);   // no root
    if (p % 4 == 3) return power(a, (p + 1) / 4, p);
    ll s = p - 1, n = 2;
    int r = 0, m;
    while (s % 2 == 0) ++r, s /= 2;
    while (power(n, (p - 1) / 2, p) != p - 1) ++n;
    ll x = power(a, (s + 1) / 2, p);
    ll b = power(a, s, p), g = power(n, s, p);
    for (;; r = m) {
        ll t = b;
        for (m = 0; m < r && t != 1; ++m) t = t * t % p;
        if (m == 0) return x;
        ll gs = power(g, 1LL << (r - m - 1), p);
        g = gs * gs % p;
        x = x * gs % p;
        b = b * g % p;
    }
}
\end{lstlisting}

\entry{Floor sums}{$O(\log m)$}{%
\code{divsum(to, c, k, m)} $= \sum_{i=0}^{to-1} \lfloor (k i + c) / m \rfloor$ and
\code{modsum(to, c, k, m)} $= \sum_{i=0}^{to-1} (k i + c) \bmod m$, for huge
\code{to} (like $10^{18}$). It works like Euclid's gcd: swap the roles of $k$ and $m$ and
recurse.\\[0.3ex]
Where it shows up: counting lattice points under a line segment (that is exactly
\code{divsum}), ``how many $i \le n$ have $(ai+b) \bmod m < c$'' (difference of two
floor sums), and sums of \code{x \% m} over an arithmetic progression. All arguments
non-negative for \code{divsum}; \code{modsum} accepts negative \code{c}, \code{k}.
Overflow is handled with \code{ull} as long as the true answer fits. From KACTL.}

%LABEL floorsum
\begin{lstlisting}
ull sumsq(ull to) { return to / 2 * ((to - 1) | 1); }
ull divsum(ull to, ull c, ull k, ull m) {
    ull res = k / m * sumsq(to) + c / m * to;
    k %= m; c %= m;
    if (!k) return res;
    ull to2 = (to * k + c) / m;
    return res + (to - 1) * to2 - divsum(to2, m-1-c, m, k);
}
ll modsum(ull to, ll c, ll k, ll m) {
    c = ((c % m) + m) % m;
    k = ((k % m) + m) % m;
    return to * c + k * sumsq(to) - m * divsum(to, c, k, m);
}
\end{lstlisting}

\entry{Binomials mod p}{precompute $O(n)$, query $O(1)$ \quad mem $n$}{%
Call \code{precompute()} once, then \code{C(n, k)}. \code{MOD} must be \textbf{prime}
and larger than \code{MAXN}. Returns 0 outside the triangle, so no guards at the call
site.\\[0.3ex]
Non-prime modulus, or $n$ small: build Pascal's triangle instead
(\code{c[i][j] = c[i-1][j-1] + c[i-1][j]}), which needs no inverses. Huge $n$ with small
prime $p$: Lucas' theorem, multiply \code{C} of the base-$p$ digits.\\[0.3ex]
Stars and bars: ways to write $n$ as an ordered sum of $k$ non-negative parts is
\code{C(n+k-1, k-1)}. Catalan $C_n = $ \code{C(2n,n) * inv(n+1)}.}

%DEP modbasics
%LABEL binom
\begin{lstlisting}
const ll MOD = 1000000007, MAXN = 500001;
ll fact[MAXN], inv_fact[MAXN];

void precompute() {
    fact[0] = 1;
    rep(i, 1, MAXN) fact[i] = fact[i-1] * i % MOD;
    inv_fact[MAXN-1] = power(fact[MAXN-1], MOD - 2, MOD);
    down(i, MAXN - 1, 0)
        inv_fact[i] = inv_fact[i+1] * (i+1) % MOD;
}
ll C(ll n, ll k) {
    if (n < 0 || k < 0 || k > n || n >= MAXN) return 0;
    return fact[n] * inv_fact[k] % MOD
                   * inv_fact[n-k] % MOD;
}
\end{lstlisting}

\entry{Matrix power}{$O(k^3 \log n)$ \quad mem $k^2$}{%
For linear recurrences: put the recurrence in a $k\times k$ matrix, raise it, read off
an entry. Fibonacci is \code{\{\{1,1\},\{1,0\}\}} to the $n$, entry \code{[0][1]}.\\[0.3ex]
Same code with \code{min}/\code{+} instead of \code{+}/\code{*} gives ``cheapest walk
using exactly $n$ edges''; with \code{or}/\code{and} it gives reachability in exactly
$n$ steps. \code{MOD} is only used in \code{mul}, so drop it for those.}

%DEP modbasics binom
\begin{lstlisting}
matrix mul(const matrix& a, const matrix& b) {   // % MOD
    int n = (int)a.size(), m = (int)b[0].size(),
        k = (int)b.size();
    matrix c(n, vi(m, 0));
    rep(i, 0, n) rep(j, 0, k) if (a[i][j])
        rep(l, 0, m)
            c[i][l] = (c[i][l] + a[i][j] * b[j][l]) % MOD;
    return c;
}
matrix mpow(matrix a, ll e) {
    int n = (int)a.size();
    matrix r(n, vi(n, 0));
    rep(i, 0, n) r[i][i] = 1;
    for (; e; e >>= 1, a = mul(a, a))
        if (e & 1) r = mul(r, a);
    return r;
}
\end{lstlisting}

\entry{Gaussian elimination}{$O(n^3)$ \quad mem $n^2$}{%
Solves $Ax=b$ over the reals. Returns the rank; \code{x} holds \emph{a} solution with
free variables set to 0. Check consistency yourself afterwards: if a row of \code{a} is
all zero but its \code{b} is not, there is no solution; if \code{rank < n} and it is
consistent, there are infinitely many.\\[0.3ex]
Partial pivoting (largest absolute pivot) is what keeps it stable --- do not remove it.
\code{eps} at $10^{-9}$ suits inputs of moderate size; raise it for large
coefficients.\\[0.3ex]
Over $\mathbb{F}_2$ use \code{bitset<N>} rows and \code{\textasciicircum=}, which is
64x faster and exact --- that is the version for XOR-basis and lights-out problems.
Mod a prime: replace \code{/} with \code{inv}.}

%LABEL gauss
\begin{lstlisting}
// solves a*x = b in place; returns the rank
int gauss(vector<vector<ld>> a, vector<ld> b,
          vector<ld>& x) {
    int n = (int)a.size(), m = (int)a[0].size(), rank = 0;
    const ld eps = 1e-9;
    vector<int> where(m, -1);
    for (int col = 0; col < m && rank < n; col++) {
        int sel = rank;
        for (int i = rank; i < n; i++)
            if (fabsl(a[i][col]) > fabsl(a[sel][col]))
                sel = i;
        if (fabsl(a[sel][col]) < eps) continue;
        swap(a[sel], a[rank]); swap(b[sel], b[rank]);
        where[col] = rank;
        for (int i = 0; i < n; i++) if (i != rank) {
            ld c = a[i][col] / a[rank][col];
            for (int j = col; j < m; j++)
                a[i][j] -= c * a[rank][j];
            b[i] -= c * b[rank];
        }
        rank++;
    }
    x.assign(m, 0);
    for (int j = 0; j < m; j++)
        if (where[j] != -1)
            x[j] = b[where[j]] / a[where[j]][j];
    return rank;
}
\end{lstlisting}

\entry{FFT / NTT (convolution)}{$O(n \log n)$ \quad mem $n$}{%
\code{conv(a, b)} multiplies two polynomials: \code{c[i] = sum a[j]*b[i-j]}, length
\code{a.size()+b.size()-1}. This is also ``for every pair of elements, bucket their
sum'' --- count pairs summing to $s$, subset-sum counts, string matching with
wildcards.\\[0.3ex]
\textbf{NTT}, exact, mod \code{P = 998244353} (a prime with $2^{23} \mid P-1$, root 3).
So the result must be wanted mod that prime, and the transform length must be a power
of two $\le 2^{23}$. Other friendly primes: $1004535809$ (root 3),
$167772161$ (root 3).\\[0.3ex]
If the true coefficients are needed exactly and exceed \code{P}, run it under two
different primes and \code{crt} the results. For real/complex input use a
\code{complex<double>} FFT instead --- but it loses precision above $\sim10^{14}$ in the
output, so prefer NTT whenever the answer is modular.}

%LABEL ntt
\begin{lstlisting}
const ll P = 998244353, G = 3;
ll pw(ll a, ll b, ll m = P) {
    ll r = 1; a %= m;
    for (; b; b >>= 1, a = a * a % m)
        if (b & 1) r = r * a % m;
    return r;
}
void ntt(vi& a, bool inv) {
    int n = (int)a.size();
    for (int i = 1, j = 0; i < n; i++) {
        int b = n >> 1;
        for (; j & b; b >>= 1) j ^= b;
        j ^= b;
        if (i < j) swap(a[i], a[j]);
    }
    for (int len = 1; len < n; len <<= 1) {
        ll w = pw(G, (P - 1) / (2 * len));
        if (inv) w = pw(w, P - 2);
        for (int i = 0; i < n; i += 2 * len) {
            ll cur = 1;
            rep(j, 0, len) {
                ll u = a[i+j], v = a[i+j+len] * cur % P;
                a[i+j] = (u + v) % P;
                a[i+j+len] = (u - v + P) % P;
                cur = cur * w % P;
            }
        }
    }
    if (inv) {
        ll ni = pw(n, P - 2);
        for (ll& x : a) x = x * ni % P;
    }
}
vi conv(vi a, vi b) {
    if (a.empty() || b.empty()) return {};
    int need = (int)(a.size() + b.size() - 1), n = 1;
    while (n < need) n <<= 1;
    a.resize(n); b.resize(n);
    ntt(a, false); ntt(b, false);
    rep(i, 0, n) a[i] = a[i] * b[i] % P;
    ntt(a, true);
    a.resize(need);
    return a;
}
\end{lstlisting}

\section{Strings}

\entry{KMP}{$O(n + m)$ \quad mem $n$}{%
\code{failure(p)[i]} is the length of the longest proper prefix of \code{p[0..i]} that is
also a suffix --- so \code{m - failure.back()} is the shortest period of \code{p}, and
\code{p} is a repetition of that period iff the period divides \code{m}.
\code{kmp(txt, pat)} returns every 0-indexed start where \code{pat} occurs, overlaps
included.\\[0.3ex]
Borders of \code{p} are \code{m}, \code{f[m-1]}, \code{f[f[m-1]-1]}, \dots\ --- walk the
chain. Counting occurrences of every prefix: add 1 at each \code{f} position and push
the counts down the chain in reverse order.}

%LABEL kmp
\begin{lstlisting}
vector<int> failure(const string& p) {
    int m = (int)p.size(); vector<int> f(m, 0);
    for (int i = 1, k = 0; i < m; i++) {
        while (k && p[i] != p[k]) k = f[k-1];
        if (p[i] == p[k]) k++;
        f[i] = k;
    }
    return f;
}
vector<int> kmp(const string& txt, const string& pat) {
    vector<int> r;
    if (pat.empty()) return r;
    vector<int> f = failure(pat);
    int n = (int)txt.size(), m = (int)pat.size();
    for (int i = 0, k = 0; i < n; i++) {
        while (k && txt[i] != pat[k]) k = f[k-1];
        if (txt[i] == pat[k]) k++;
        if (k == m) { r.push_back(i - m + 1); k = f[k-1]; }
    }
    return r;
}
\end{lstlisting}

\entry{Z function}{$O(n)$ \quad mem $n$}{%
\code{z[i]} is the length of the longest common prefix of \code{s} and \code{s[i..]}
(\code{z[0]} is left 0). Pattern matching: build \code{pat + '\textbackslash1' + txt} and
look for \code{z[i] == pat.size()}; the separator must not occur in either
string.\\[0.3ex]
Often replaces KMP, and is the easier one to get right. Distinct substrings, smallest
rotation, and ``is this a repetition'' all fall out of it.}

%LABEL zfunc
\begin{lstlisting}
vector<int> zfunc(const string& s) {
    int n = (int)s.size();
    vector<int> z(n, 0);
    for (int i = 1, l = 0, r = 0; i < n; i++) {
        if (i < r) z[i] = min(r - i, z[i - l]);
        while (i + z[i] < n && s[z[i]] == s[i + z[i]]) z[i]++;
        if (i + z[i] > r) { l = i; r = i + z[i]; }
    }
    return z;
}
\end{lstlisting}

\entry{Manacher}{$O(n)$ \quad mem $n$}{%
Every palindromic substring at once. \code{d1[i]} is the radius of the longest
\emph{odd} palindrome centred at \code{i} (always $\ge 1$), \code{d2[i]} the radius of
the longest \emph{even} one centred between \code{i-1} and \code{i} (may be 0). The
number of palindromic substrings is $\sum d1 + \sum d2$.\\[0.3ex]
\code{isPal(l, r)} answers ``is \code{s[l..r)} a palindrome'' in $O(1)$, which is what
you usually actually want --- e.g. for ``split the string into the fewest
palindromes'' DP. Longest palindromic substring:
$\max(2 d1[i] - 1, 2 d2[i])$.}

%LABEL manacher
\begin{lstlisting}
struct Manacher {
    int n; vector<int> d1, d2;
    Manacher(const string& s) : n((int)s.size()),
                                d1(n, 0), d2(n, 0) {
        for (int i = 0, l = 0, r = -1; i < n; i++) {
            int k = i > r ? 1 : min(d1[l+r-i], r-i+1);
            while (i-k >= 0 && i+k < n && s[i-k] == s[i+k])
                k++;
            d1[i] = k--;
            if (i + k > r) { l = i - k; r = i + k; }
        }
        for (int i = 0, l = 0, r = -1; i < n; i++) {
            int k = i > r ? 0 : min(d2[l+r-i+1], r-i+1);
            while (i - k - 1 >= 0 && i + k < n
                   && s[i-k-1] == s[i+k]) k++;
            d2[i] = k--;
            if (i + k > r) { l = i - k - 1; r = i + k; }
        }
    }
    bool isPal(int l, int r) {                // s[l..r)
        int len = r - l, c = l + len / 2;
        return len % 2 ? 2*d1[c]-1 >= len : 2*d2[c] >= len;
    }
};
\end{lstlisting}

\entry{String hashing}{build $O(n)$, compare $O(1)$ \quad mem $n$}{%
\code{get(l, r)} is the hash of \code{s[l..r)}; two substrings are equal iff their
hashes match. Two independent moduli, so a collision needs both to fail at once ---
with a single 64-bit overflow hash an adversarial test \emph{will} break you.\\[0.3ex]
This is the workhorse: compare any two substrings, binary-search the longest common
prefix of two suffixes, dedupe substrings by putting \code{get} results in a
\code{set}, Rabin-Karp matching. Often replaces a suffix array
outright.\\[0.3ex]
For randomised bases (safer against hash-killers) pick \code{B1}, \code{B2} at run time
from \code{mt19937} in $[256, M)$.}

%LABEL hash
\begin{lstlisting}
struct Hash {
    static constexpr ll M1 = 1000000007, M2 = 998244353,
                        B1 = 131, B2 = 137;
    int n; vi h1, h2, p1, p2;
    Hash(const string& s) : n((int)s.size()), h1(n+1, 0),
            h2(n+1, 0), p1(n+1, 1), p2(n+1, 1) {
        rep(i, 0, n) {
            h1[i+1] = (h1[i] * B1 + s[i]) % M1;
            h2[i+1] = (h2[i] * B2 + s[i]) % M2;
            p1[i+1] = p1[i] * B1 % M1;
            p2[i+1] = p2[i] * B2 % M2;
        }
    }
    pii get(int l, int r) {                   // s[l..r)
        ll a = (h1[r] - h1[l] * p1[r-l]) % M1;
        ll b = (h2[r] - h2[l] * p2[r-l]) % M2;
        return {(a + M1) % M1, (b + M2) % M2};
    }
};
\end{lstlisting}

\entry{Suffix array + LCP}{$O(n \log n)$ \quad mem $n$}{%
\code{sa[i]} is the start of the \code{i}-th smallest suffix; \code{rnk} is its inverse;
\code{lcp[i]} is the longest common prefix of \code{sa[i]} and \code{sa[i+1]}.
The input must not contain the byte \code{1}, which is used as a
sentinel.\\[0.3ex]
Number of distinct substrings is $\frac{n(n+1)}{2} - \sum \mathrm{lcp}$. Longest
repeated substring is $\max \mathrm{lcp}$. Longest common substring of two strings: join
with a separator and take the best \code{lcp} between suffixes from different sides.
LCP of two arbitrary suffixes = range-min of \code{lcp} between their ranks --- feed
\code{lcp} to \code{SpTab}.\\[0.3ex]
Consider \code{Hash} plus binary search first: shorter to write, usually fast enough.}

%LABEL sufarr
\begin{lstlisting}
struct SuffixArray {
    int n; vector<int> sa, rnk, lcp;
    SuffixArray(const string& str) {
        string s = str; s += char(1);       // sentinel
        n = (int)s.size();
        vector<int> p(n), c(n), pn(n), cn(n);
        vector<int> cnt(max(256, n), 0);
        for (int i = 0; i < n; i++)
            cnt[(unsigned char)s[i]]++;
        for (int i = 1; i < 256; i++) cnt[i] += cnt[i-1];
        for (int i = 0; i < n; i++)
            p[--cnt[(unsigned char)s[i]]] = i;
        int cls = 1; c[p[0]] = 0;
        for (int i = 1; i < n; i++) {
            if (s[p[i]] != s[p[i-1]]) cls++;
            c[p[i]] = cls - 1;
        }
        for (int h = 0; (1 << h) < n; h++) {
            for (int i = 0; i < n; i++) {
                pn[i] = p[i] - (1 << h);
                if (pn[i] < 0) pn[i] += n;
            }
            fill(cnt.begin(), cnt.begin() + cls, 0);
            for (int i = 0; i < n; i++) cnt[c[pn[i]]]++;
            for (int i = 1; i < cls; i++) cnt[i] += cnt[i-1];
            for (int i = n - 1; i >= 0; i--)
                p[--cnt[c[pn[i]]]] = pn[i];
            cn[p[0]] = 0; cls = 1;
            for (int i = 1; i < n; i++) {
                int a = p[i], b = p[i-1], d = 1 << h;
                if (c[a] != c[b]
                    || c[(a+d) % n] != c[(b+d) % n]) cls++;
                cn[a] = cls - 1;
            }
            c = cn;
        }
        sa.assign(p.begin() + 1, p.end());   // drop sentinel
        n--;
        rnk.assign(n, 0);
        for (int i = 0; i < n; i++) rnk[sa[i]] = i;
        lcp.assign(max(n - 1, 0), 0);
        for (int i = 0, k = 0; i < n; i++) {
            if (rnk[i] == n - 1) { k = 0; continue; }
            int j = sa[rnk[i] + 1];
            while (i + k < n && j + k < n
                   && str[i+k] == str[j+k]) k++;
            lcp[rnk[i]] = k;
            if (k) k--;
        }
    }
};
\end{lstlisting}

\entry{Aho-Corasick}{build $O(\Sigma \cdot \mathrm{tot})$, scan $O(n + \mathrm{hits})$ \quad mem $\Sigma \cdot \mathrm{tot}$}{%
Many patterns at once. \code{add(s, id)} before \code{build()}, then \code{match(txt, f)}
calls \code{f(i, id)} for every pattern \code{id} ending at text position \code{i}.
Lowercase \code{a-z}; change the 26 and the \code{- 'a'} for another
alphabet.\\[0.3ex]
\code{out} is the \emph{output link} --- the nearest fail-ancestor that is itself a
pattern end --- so reporting costs $O(1)$ per hit instead of walking the whole fail
chain, which is the difference between $O(n + \mathrm{hits})$ and $O(n \cdot
\mathrm{depth})$. Empty patterns are not allowed (the root is never a match).\\[0.3ex]
To \emph{count} hits per pattern without listing them, mark \code{cnt[u]++} per text
position and push counts along \code{fail} in reverse BFS order --- $O(n +
\mathrm{tot})$ regardless of how many matches there are.}

%LABEL aho
\begin{lstlisting}
struct Aho {
    struct Nd { int go[26], fail = 0, out = 0, head = -1; };
    vector<Nd> t;
    vector<int> idv, nxt;        // id lists on terminal nodes
    Aho() { newNode(); }
    int newNode() {
        t.push_back(Nd{});
        fill(t.back().go, t.back().go + 26, -1);
        return (int)t.size() - 1;
    }
    void add(const string& s, int id) {  // nonempty s
        int u = 0;
        for (char ch : s) {
            int c = ch - 'a';
            if (t[u].go[c] < 0) {
                int v = newNode(); t[u].go[c] = v;
            }
            u = t[u].go[c];
        }
        idv.push_back(id); nxt.push_back(t[u].head);
        t[u].head = (int)idv.size() - 1;
    }
    void build() {
        vector<int> q;
        for (int c = 0; c < 26; c++) {
            int v = t[0].go[c];
            if (v < 0) t[0].go[c] = 0;
            else { t[v].fail = 0; q.push_back(v); }
        }
        for (int i = 0; i < (int)q.size(); i++) {
            int u = q[i], f = t[u].fail;
            t[u].out = t[f].head >= 0 ? f : t[f].out;
            for (int c = 0; c < 26; c++) {
                int v = t[u].go[c], w = t[f].go[c];
                if (v < 0) t[u].go[c] = w;
                else { t[v].fail = w; q.push_back(v); }
            }
        }
    }
    template<class F>       // f(end index, pattern id)
    void match(const string& s, F f) {
        int u = 0;
        for (int i = 0; i < (int)s.size(); i++) {
            u = t[u].go[s[i] - 'a'];
            for (int v = t[u].head >= 0 ? u : t[u].out; v;
                 v = t[v].out)
                for (int e = t[v].head; e >= 0; e = nxt[e])
                    f(i, idv[e]);
        }
    }
};
\end{lstlisting}

\entry{XOR trie}{insert / query $O(\mathrm{BITS})$ \quad mem $n \cdot \mathrm{BITS}$}{%
\code{maxXor(x)} is the largest \code{x \textasciicircum{} y} over everything inserted:
walk from the top bit and always take the opposite branch if it exists. Raise
\code{BITS} to 62 for \code{ll} values.\\[0.3ex]
Keep a counter per node and you can also answer ``how many \code{y} with
\code{x \textasciicircum{} y < k}'', and delete by decrementing. Maximum XOR
\emph{subarray}: insert prefix xors as you sweep.\\[0.3ex]
A plain prefix trie over letters is the same shape with 26 children --- see
\code{Aho} without the fail links.}

\begin{lstlisting}
struct XorTrie {
    static const int BITS = 30;
    vector<array<int,2>> ch{{-1, -1}};
    void insert(ll x) {
        int u = 0;
        down(b, BITS + 1, 0) {
            int c = x >> b & 1;
            if (ch[u][c] < 0) {
                ch.push_back({-1, -1});
                ch[u][c] = (int)ch.size() - 1;
            }
            u = ch[u][c];
        }
    }
    ll maxXor(ll x) {         // needs >= 1 inserted value
        int u = 0; ll r = 0;
        down(b, BITS + 1, 0) {
            int c = (x >> b & 1) ^ 1;
            if (ch[u][c] >= 0) { r |= 1LL<<b; u = ch[u][c]; }
            else u = ch[u][c ^ 1];
        }
        return r;
    }
};
\end{lstlisting}

\entry{Minimum rotation}{$O(n)$}{%
Index where the lexicographically smallest rotation of \code{s} starts:
\code{rotate(s.begin(), s.begin() + minRotation(s), s.end())}. Two strings (or
necklaces, or cyclic sequences of numbers) are rotations of each other iff their
minimum rotations are equal --- the standard way to put a cycle into one canonical form
before hashing or comparing. Works on \code{vector<int>} too if you change the type.
From KACTL.}

%LABEL minrot
\begin{lstlisting}
int minRotation(string s) {
    int a = 0, N = (int)s.size(); s += s;
    rep(b, 0, N) rep(k, 0, N) {
        if (a + k == b || s[a+k] < s[b+k]) {
            b += max(0, (int)k - 1); break;
        }
        if (s[a+k] > s[b+k]) { a = (int)b; break; }
    }
    return a;
}
\end{lstlisting}

\section{Geometry}

\entry{Point / vector}{}{%
One template instead of a float-only \code{Point}. \code{Pl} is \code{P<ll>} and is
\textbf{exact} --- use it for everything decided by a sign (orientation, hulls,
point-in-polygon, segment crossing) and you have no \code{eps} to tune and no wrong
answers. \code{Pd} is \code{P<ld>}, for the metric answers where you must divide or take
a root.\\[0.3ex]
\begin{ctab}{0.22}{0.73}
\code{cross(p)} & $x p_y - y p_x$. Sign = which side; $|\cdot|$ = area of the
parallelogram.\\
\code{a.cross(b,c)} & \code{(b-a).cross(c-a)}, the turn at \code{a}.\\
\code{dot(p)} & $<0$ means the angle is obtuse.\\
\code{dist2()} & squared length, exact for \code{Pl}. Compare these, never
\code{dist()}.\\
\code{perp()} & rotate $90^\circ$ counter-clockwise.\\
\code{angle()} & \code{atan2}, in $(-\pi, \pi]$.\\
\end{ctab}\\[0.3ex]
\code{unit()} and \code{rotate()} only make sense for \code{Pd} (they divide / call
\code{cos}); they simply will not compile for \code{Pl}, which is the point.
Overflow: \code{cross} on coordinates up to $10^9$ reaches $2\cdot10^{18}$, just inside
\code{ll} --- above that use \code{i128}.\\[0.3ex]
\code{ccw(a,b,c)} is \code{+1} if \code{a->b->c} turns left, \code{-1} right, \code{0}
collinear. Nearly every predicate below is built from it.}

%LABEL point
\begin{lstlisting}
template<class T> struct P {
    T x, y;
    P(T x = 0, T y = 0) : x(x), y(y) {}
    P operator+(P p) const { return {x + p.x, y + p.y}; }
    P operator-(P p) const { return {x - p.x, y - p.y}; }
    P operator*(T k) const { return {x * k, y * k}; }
    P operator/(T k) const { return {x / k, y / k}; }
    T dot(P p)   const { return x * p.x + y * p.y; }
    T cross(P p) const { return x * p.y - y * p.x; }
    T cross(P a, P b) const {         // turn at *this
        return (a - *this).cross(b - *this);
    }
    T dist2()  const { return x * x + y * y; }
    ld dist()  const { return sqrtl((ld)dist2()); }
    ld angle() const { return atan2l((ld)y, (ld)x); }
    P perp()   const { return {-y, x}; }      // 90 deg ccw
    P unit() const { return *this / (T)dist(); } // Pd only
    P rotate(ld a) const {             // Pd only
        return {x * cosl(a) - y * sinl(a),
                x * sinl(a) + y * cosl(a)};
    }
    bool operator<(P p) const {
        return tie(x, y) < tie(p.x, p.y);
    }
    bool operator==(P p) const {
        return tie(x, y) == tie(p.x, p.y);
    }
};
using Pl = P<ll>;          // exact: signs, hulls, areas
using Pd = P<ld>;          // metric: distances, intersections

template<class T> int sgn(T v) { return (v > 0) - (v < 0); }
template<class T> int ccw(P<T> a, P<T> b, P<T> c) {
    return sgn(a.cross(b, c));
}
\end{lstlisting}

\entry{Segments and lines}{exact where it matters}{%
\code{onSeg} and \code{segInter} are \textbf{exact} on \code{Pl} and handle every
degenerate case --- collinear overlap, shared endpoints, zero-length segments. That is
the usual source of wrong answers, so prefer these over an orientation-only
test.\\[0.3ex]
\code{lineInter} works on infinite lines and returns \code{false} when they are parallel
(including identical). \code{segDist} is point-to-\emph{segment}; for
segment-to-segment take the min of the four endpoint-to-segment distances, or 0 if they
cross.}

%DEP point
%LABEL seg
\begin{lstlisting}
bool onSeg(Pl p, Pl a, Pl b) {          // p on segment ab
    return (b - a).cross(p - a) == 0
        && (a - p).dot(b - p) <= 0;
}
bool segInter(Pl a, Pl b, Pl c, Pl d) {
    int d1 = ccw(c, d, a), d2 = ccw(c, d, b);
    int d3 = ccw(a, b, c), d4 = ccw(a, b, d);
    if (d1 * d2 < 0 && d3 * d4 < 0) return true;
    return onSeg(a, c, d) || onSeg(b, c, d)
        || onSeg(c, a, b) || onSeg(d, a, b);
}
// lines ab and cd; false if parallel
bool lineInter(Pd a, Pd b, Pd c, Pd d, Pd& out) {
    ld den = (b - a).cross(d - c);
    if (fabsl(den) < 1e-12) return false;
    out = a + (b - a) * ((c - a).cross(d - c) / den);
    return true;
}
ld segDist(Pd p, Pd a, Pd b) {          // point to segment
    if (a == b) return (p - a).dist();
    ld t = (p - a).dot(b - a) / (ld)(b - a).dist2();
    t = max((ld)0, min((ld)1, t));
    return (p - (a + (b - a) * t)).dist();
}
\end{lstlisting}

\entry{Polygons}{$O(n)$, \code{inConvex} $O(\log n)$}{%
\code{area2} is \emph{twice} the signed area (shoelace), exact: positive means the
vertices are counter-clockwise. The real area is \code{abs(area2)/2.0} --- keep the
doubled integer as long as you can.\\[0.3ex]
\code{inPoly} works for any simple polygon, convex or not, and returns \code{2} on the
boundary so you can decide what counts. \code{inConvex} is the $O(\log n)$ version and
needs a strictly convex counter-clockwise polygon, i.e. exactly what \code{hull}
returns.\\[0.3ex]
\textbf{Pick's theorem} (lattice polygon): $A = I + B/2 - 1$, so with \code{area2} and
the boundary count $B = \sum \gcd(|dx|, |dy|)$ over the edges, the interior count is
$I = (\texttt{area2} - B + 2)/2$.}

%DEP point seg
%LABEL poly
\begin{lstlisting}
ll area2(const vector<Pl>& p) {        // signed, 2x area
    ll a = 0; int n = (int)p.size();
    rep(i, 0, n) a += p[i].cross(p[(i+1) % n]);
    return a;
}
// 0 = outside, 1 = inside, 2 = on the boundary
int inPoly(const vector<Pl>& p, Pl q) {
    int n = (int)p.size(), c = 0;
    rep(i, 0, n) {
        Pl a = p[i], b = p[(i+1) % n];
        if (onSeg(q, a, b)) return 2;
        if (a.y > b.y) swap(a, b);
        if (a.y <= q.y && q.y < b.y
            && (b - a).cross(q - a) > 0) c ^= 1;
    }
    return c;
}
// p = strictly convex, ccw (what hull() gives)
int inConvex(const vector<Pl>& p, Pl q) {
    int n = (int)p.size();
    if (n == 1) return p[0] == q ? 2 : 0;
    if (n == 2) return onSeg(q, p[0], p[1]) ? 2 : 0;
    if (ccw(p[0], p[1], q) < 0) return 0;
    if (ccw(p[0], p[n-1], q) > 0) return 0;
    int lo = 1, hi = n - 1;
    while (hi - lo > 1) {
        int m = (lo + hi) / 2;
        if (ccw(p[0], p[m], q) >= 0) lo = m; else hi = m;
    }
    int c = ccw(p[lo], p[lo+1], q);
    if (c < 0) return 0;
    if (c == 0) return 2;
    if (onSeg(q, p[0], p[1]) || onSeg(q, p[0], p[n-1]))
        return 2;
    return 1;
}
\end{lstlisting}

\entry{Convex hull}{$O(n \log n)$ \quad mem $n$}{%
Andrew's monotone chain. Counter-clockwise, \textbf{no collinear points}, first vertex
not repeated at the end. Duplicates in the input are removed. Fewer than 3 distinct
points come back as-is.\\[0.3ex]
To \emph{keep} collinear boundary points change \code{<= 0} to \code{< 0} --- but then
the ``is every input point on the hull'' test becomes
\code{hull(p).size() == p.size()}, which is what Ruiming's original was checking.\\[0.3ex]
Diameter of a point set = farthest pair on the hull: rotating calipers, or just
$O(h^2)$ over hull vertices when $h$ is small. Minimum-width / minimum-area enclosing
rectangle also come from calipers on the hull.}

%DEP point
%LABEL hull
\begin{lstlisting}
vector<Pl> hull(vector<Pl> p) {
    sort(all(p)); p.erase(unique(all(p)), p.end());
    int n = (int)p.size();
    if (n < 3) return p;
    vector<Pl> h;
    rep(phase, 0, 2) {              // lower, then upper
        size_t start = h.size();
        for (Pl q : p) {
            while (h.size() >= start + 2
                   && ccw(h[h.size()-2], h.back(), q) <= 0)
                h.pop_back();
            h.push_back(q);
        }
        h.pop_back();
        reverse(all(p));
    }
    return h;
}
\end{lstlisting}

\entry{Closest pair}{$O(n \log n)$ \quad mem $n$}{%
Returns the smallest \textbf{squared} distance and the two points. Exact on \code{Pl};
the only floating point is a deliberately generous \code{sqrtl} bound on the sweep
window, so rounding cannot lose the answer. Needs $n \ge 2$; duplicate points give 0,
which is the correct answer.\\[0.3ex]
Sweep in \code{x}, keep a \code{set} ordered by \code{y} holding everything within the
current best distance, and only compare against the \code{y}-band. Coordinates up to
$10^9$ keep \code{dist2} inside \code{ll}.}

%DEP point
%LABEL closest
\begin{lstlisting}
pair<ll, pair<Pl,Pl>> closest(vector<Pl> p) {
    sort(all(p));                        // by x, then y
    ll best = LLONG_MAX;
    pair<Pl,Pl> bp{p[0], p[1]};
    set<pii> s;                          // (y, x)
    int j = 0;
    rep(i, 0, (ll)p.size()) {
        while (j < i) {
            ll dx = p[i].x - p[j].x;
            if (dx * dx <= best) break;
            s.erase({p[j].y, p[j].x}); j++;
        }
        ll d = (ll)sqrtl((ld)best) + 1;
        auto lo = s.lower_bound({p[i].y - d, LLONG_MIN});
        auto hi = s.upper_bound({p[i].y + d, LLONG_MAX});
        for (auto it = lo; it != hi; ++it) {
            Pl q(it->second, it->first);
            ll dd = (p[i] - q).dist2();
            if (dd < best) { best = dd; bp = {q, p[i]}; }
        }
        s.insert({p[i].y, p[i].x});
    }
    return {best, bp};
}
\end{lstlisting}

\entry{Circles}{}{%
\code{circleLine} intersects a circle with the \emph{infinite} line through \code{a,b};
\code{circleCircle} with another circle. Both return 0, 1 (tangent) or 2 points. For a
\emph{segment} instead, intersect with the line and then test each result with
\code{segDist} or a parameter check.\\[0.3ex]
Two identical circles return \code{\{\}} --- infinitely many points is not something the
caller can use, so check that case yourself.\\[0.3ex]
Tangent length from an outside point \code{p} to circle \code{(c,r)} is
$\sqrt{|pc|^2 - r^2}$. Circle through three points: the circumcentre is the intersection
of two perpendicular bisectors, i.e. \code{lineInter} of
\code{(m1, m1+(b-a).perp())} and \code{(m2, m2+(c-a).perp())} with \code{m} the
midpoints.}

%DEP point
%LABEL circle
\begin{lstlisting}
vector<Pd> circleLine(Pd c, ld r, Pd a, Pd b) {
    Pd ab = b - a;
    Pd m = a + ab * ((c - a).dot(ab) / ab.dist2());
    ld h2 = r * r - (m - c).dist2();
    if (h2 < -1e-12) return {};
    if (h2 <= 1e-12) return {m};
    Pd h = ab.unit() * sqrtl(h2);
    return {m - h, m + h};
}
vector<Pd> circleCircle(Pd c1, ld r1, Pd c2, ld r2) {
    Pd d = c2 - c1;
    ld dd = d.dist2(), dist = sqrtl(dd);
    if (dist < 1e-12) return {};          // concentric
    if (dist > r1 + r2 + 1e-12) return {};
    if (dist < fabsl(r1 - r2) - 1e-12) return {};
    ld a = (dd + r1 * r1 - r2 * r2) / (2 * dist);
    ld h2 = r1 * r1 - a * a;
    Pd m = c1 + d.unit() * a;
    if (h2 <= 1e-12) return {m};
    Pd per = d.unit().perp() * sqrtl(h2);
    return {m - per, m + per};
}
\end{lstlisting}

\entry{Projection, reflection, angle sort}{}{%
\code{lineProj(a, b, p)} is the point on line \code{ab} closest to \code{p} (the foot
of the perpendicular); with \code{refl = true} it is \code{p} mirrored across the line.
Use \code{Pd}; with \code{Pl} the division truncates.\\[0.3ex]
\code{angleLess} sorts vectors by angle around the origin, counter-clockwise starting
from the positive $x$-axis, \textbf{exactly} on \code{Pl} (no \code{atan2}). Ruiming's
\code{cmp} with \code{atan2l} does the same in floating point; fine for small
coordinates, but two nearly equal angles can swap. To sort around a point \code{c},
sort \code{p - c}. The origin itself must not be in the list. Vectors at the same angle
compare equal; break ties by \code{dist2()} if needed.}

%DEP point
%LABEL lineproj
\begin{lstlisting}
template<class T>
P<T> lineProj(P<T> a, P<T> b, P<T> p, bool refl = false) {
    P<T> v = b - a;
    return p - v.perp() * (1 + refl) * v.cross(p - a)
             / v.dist2();
}
bool upperHalf(Pl p) {
    return p.y > 0 || (p.y == 0 && p.x > 0);
}
bool angleLess(Pl a, Pl b) {         // ccw from +x axis
    bool ha = upperHalf(a), hb = upperHalf(b);
    if (ha != hb) return ha;
    return a.cross(b) > 0;
}
\end{lstlisting}

\entry{Circumcircle, minimum enclosing circle}{\code{mec} expected $O(n)$}{%
\code{ccCenter(A, B, C)} is the centre of the circle through three points (they must
not be collinear) and \code{ccRadius} its radius.\\[0.3ex]
\code{mec(ps)} is the smallest circle containing every point: returns
\code{\{centre, radius\}}. Randomised (Welzl): the points are shuffled, and whenever a
point lies outside the current circle it must be on the boundary of the new one, which
is rebuilt from 1, 2 or 3 boundary points. Expected linear time --- $10^6$ points is
fine. Also answers ``place one tower to cover all houses with the smallest range''.
Ruiming's \code{min\_center\_circle} and KACTL's \code{mec} are the same algorithm; this
is it on \code{Pd}.}

%DEP point
%LABEL mec
\begin{lstlisting}
ld ccRadius(Pd A, Pd B, Pd C) {
    return (B - A).dist() * (C - B).dist() * (A - C).dist()
         / fabsl((B - A).cross(C - A)) / 2;
}
Pd ccCenter(Pd A, Pd B, Pd C) {
    Pd b = C - A, c = B - A;
    return A + (b * c.dist2() - c * b.dist2()).perp()
                 / b.cross(c) / 2;
}
pair<Pd, ld> mec(vector<Pd> ps) {
    shuffle(all(ps), mt19937(1234));
    Pd o = ps[0]; ld r = 0, EPS = 1 + 1e-10;
    rep(i, 0, (ll)ps.size()) if ((o - ps[i]).dist() > r*EPS) {
        o = ps[i], r = 0;
        rep(j, 0, i) if ((o - ps[j]).dist() > r * EPS) {
            o = (ps[i] + ps[j]) / 2;
            r = (o - ps[i]).dist();
            rep(k, 0, j) if ((o - ps[k]).dist() > r * EPS) {
                o = ccCenter(ps[i], ps[j], ps[k]);
                r = (o - ps[i]).dist();
            }
        }
    }
    return {o, r};
}
\end{lstlisting}

\entry{Circle tangents}{}{%
\code{tangents(c1, r1, c2, r2)} gives the lines that touch both circles, each as a pair
of touching points (on circle 1, on circle 2). These are the \emph{outer} tangents; pass
\code{-r2} for the \emph{inner} ones (the ones that cross between the circles). 0, 1 or
2 results: 0 if one circle is inside the other (or they are equal), 1 if they touch.\\[0.3ex]
Tangents from a point \code{p} to a circle: \code{tangents(c, r, p, 0)}; the touching
points are the \code{.first}s. Shortest path around a round obstacle goes along a
tangent, then the arc, then a tangent --- this is the building block. From KACTL.}

%DEP point
%LABEL tangents
\begin{lstlisting}
vector<pair<Pd,Pd>> tangents(Pd c1, ld r1, Pd c2, ld r2) {
    Pd d = c2 - c1;
    ld dr = r1 - r2, d2 = d.dist2(), h2 = d2 - dr * dr;
    if (d2 == 0 || h2 < 0) return {};
    vector<pair<Pd,Pd>> out;
    for (ld sg : {-1, 1}) {
        Pd v = (d * dr + d.perp() * sqrtl(h2) * sg) / d2;
        out.push_back({c1 + v * r1, c2 + v * r2});
    }
    if (h2 == 0) out.pop_back();
    return out;
}
\end{lstlisting}

\entry{Circle-polygon intersection area}{$O(n)$}{%
Area of the part of a polygon (counter-clockwise, need not be convex) that lies inside
the circle \code{(c, r)}. Works by adding up signed pieces for each edge as seen from
the centre, so it handles every case without special code. A clockwise polygon gives the
negative answer. From KACTL.}

%DEP point
%LABEL circlePoly
\begin{lstlisting}
ld circlePoly(Pd c, ld r, vector<Pd> ps) {
    auto arg = [](Pd p, Pd q) {
        return atan2l(p.cross(q), p.dot(q)); };
    auto tri = [&](Pd p, Pd q) {
        ld r2 = r * r / 2;
        Pd d = q - p;
        ld a = d.dot(p) / d.dist2();
        ld b = (p.dist2() - r * r) / d.dist2();
        ld det = a * a - b;
        if (det <= 0) return arg(p, q) * r2;
        ld s = max((ld)0, -a - sqrtl(det));
        ld t = min((ld)1, -a + sqrtl(det));
        if (t < 0 || 1 <= s) return arg(p, q) * r2;
        Pd u = p + d * s, v = q + d * (t - 1);
        return arg(p, u) * r2 + u.cross(v) / 2
             + arg(v, q) * r2;
    };
    ld sum = 0;
    rep(i, 0, (ll)ps.size())
        sum += tri(ps[i] - c, ps[(i + 1) % ps.size()] - c);
    return sum;
}
\end{lstlisting}

\entry{Polygon cut, centroid}{$O(n)$}{%
\code{polygonCut(poly, s, e)} keeps the part of the polygon to the \textbf{right} of the
directed line \code{s -> e} (everything on the left is cut away). The polygon should be
convex; for a non-convex one the result can contain zero-width bridges.\\[0.3ex]
\textbf{Half-plane intersection for small inputs}: start from a huge square and cut
once per half-plane --- $O(n^2)$, simple, and enough for $n \le 5000$. Keep the
\emph{left} side of \code{s -> e} by cutting with \code{(e, s)}.\\[0.3ex]
\code{polygonCenter} is the centre of mass of the filled polygon (not the average of
the vertices). Both from KACTL.}

%DEP point
%LABEL polyCut
\begin{lstlisting}
vector<Pd> polygonCut(const vector<Pd>& poly, Pd s, Pd e) {
    vector<Pd> res;
    rep(i, 0, (ll)poly.size()) {
        Pd cur = poly[i], prev = i ? poly[i-1] : poly.back();
        ld a = s.cross(e, cur), b = s.cross(e, prev);
        if ((a < 0) != (b < 0))
            res.push_back(cur + (prev - cur) * (a / (a - b)));
        if (a < 0) res.push_back(cur);
    }
    return res;
}
Pd polygonCenter(const vector<Pd>& v) {
    Pd res(0, 0); ld A = 0;
    for (int i = 0, j = (int)v.size() - 1; i < (int)v.size();
         j = i++) {
        res = res + (v[i] + v[j]) * v[j].cross(v[i]);
        A += v[j].cross(v[i]);
    }
    return res / A / 3;
}
\end{lstlisting}

\entry{Hull diameter (rotating calipers)}{$O(n)$}{%
The two points of a convex polygon that are farthest apart. Input is what \code{hull}
returns (counter-clockwise, no collinear points); for a point set call \code{hull}
first. Returns the pair; \code{(a - b).dist2()} is the squared diameter, exact.\\[0.3ex]
How: walk \code{i} around the hull and keep a second pointer \code{j} at the vertex
farthest from edge \code{i}; \code{j} only ever moves forward, so the total is
linear. The same two-pointer walk gives minimum width (smallest edge-to-farthest-vertex
distance) and the smallest enclosing rectangle. From KACTL.}

%DEP point
%LABEL hullDiam
\begin{lstlisting}
array<Pl,2> hullDiameter(const vector<Pl>& S) {
    int n = (int)S.size(), j = n < 2 ? 0 : 1;
    pair<ll, array<Pl,2>> res({0, {S[0], S[0]}});
    rep(i, 0, j)
        for (;; j = (j + 1) % n) {
            res = max(res, {(S[i] - S[j]).dist2(),
                            {S[i], S[j]}});
            if ((S[(j+1)%n] - S[j]).cross(S[i+1] - S[i]) >= 0)
                break;
        }
    return res.second;
}
\end{lstlisting}

\entry{Manhattan MST}{$O(n \log n)$}{%
Minimum spanning tree where the cost of joining two points is
$|x_1 - x_2| + |y_1 - y_2|$. The complete graph has $n^2$ edges, too many; this returns
at most $4n$ candidate edges \code{\{w, a, b\}} that are guaranteed to contain an MST.
Feed them straight to \code{kruskal}. (For each point it only keeps the nearest point
in each of 8 directions, found by a sweep.)\\[0.3ex]
Chebyshev distance $\max(|dx|, |dy|)$ instead: map every point to $(x + y, x - y)$.
Manhattan distance between the new points is exactly \emph{twice} the Chebyshev
distance between the old ones, so the same tree works; halve the total. From KACTL.}

%DEP point
%LABEL manhattan
\begin{lstlisting}
vector<array<ll,3>> manhattanMST(vector<Pl> ps) {
    vector<int> id(ps.size()); iota(all(id), 0);
    vector<array<ll,3>> edges;
    rep(k, 0, 4) {
        sort(all(id), [&](int i, int j) {
            return (ps[i] - ps[j]).x < (ps[j] - ps[i]).y; });
        map<ll, int> sweep;
        for (int i : id) {
            for (auto it = sweep.lower_bound(-ps[i].y);
                 it != sweep.end(); sweep.erase(it++)) {
                int j = it->second;
                Pl d = ps[i] - ps[j];
                if (d.y > d.x) break;
                edges.push_back({d.y + d.x, i, j});
            }
            sweep[-ps[i].y] = i;
        }
        for (Pl& p : ps)
            if (k & 1) p.x = -p.x; else swap(p.x, p.y);
    }
    return edges;
}
\end{lstlisting}

\section{Searching}

\entry{Binary search on a predicate}{$O(\log)$}{%
\code{pred} must be monotone false$\dots$false true$\dots$true. Returns the first
\code{x} in \code{[lo, hi]} with \code{pred(x)}, or \code{hi+1} if none --- so test the
return value before using it.\\[0.3ex]
``Largest \code{x} that still works'' is the same call on the negated predicate minus
one, or just search for the first failing \code{x} and subtract 1. Almost every
``minimise the maximum'' / ``can we do it within $k$'' problem is this function plus a
greedy check.\\[0.3ex]
Reals: fix the iteration count (100 halvings covers any range to $10^{-30}$) rather than
looping on a tolerance, which can spin forever.}

\begin{lstlisting}
template<class F>
ll bsearch(ll lo, ll hi, F pred) {   // first true in [lo,hi]
    hi++;                            // hi = "none found"
    while (lo < hi) {
        ll m = lo + (hi - lo) / 2;   // no overflow
        if (pred(m)) hi = m; else lo = m + 1;
    }
    return lo;
}
template<class F>
ld bsearchReal(ld lo, ld hi, F pred) {
    rep(it, 0, 100) {
        ld m = (lo + hi) / 2;
        if (pred(m)) hi = m; else lo = m;
    }
    return lo;
}
\end{lstlisting}

\entry{Ternary search}{$O(\log)$}{%
For a \textbf{strictly unimodal} function --- one valley, no flat stretches. A flat
region breaks it silently, so if ties are possible make the comparison
\code{<=} vs \code{<} deliberately, or binary-search on
\code{f(m+1) - f(m)} instead, which is the more robust trick.\\[0.3ex]
The integer version finishes by scanning the last three candidates, which is where
off-by-ones normally live. For a maximum, negate \code{f}.}

\begin{lstlisting}
template<class F>              // argmin of a unimodal f
ll tsearch(ll lo, ll hi, F f) {
    while (hi - lo > 2) {
        ll m1 = lo + (hi - lo) / 3, m2 = hi - (hi - lo) / 3;
        if (f(m1) < f(m2)) hi = m2; else lo = m1;
    }
    ll best = lo;
    for (ll x = lo + 1; x <= hi; x++)
        if (f(x) < f(best)) best = x;
    return best;
}
template<class F>
ld tsearchReal(ld lo, ld hi, F f) {
    rep(it, 0, 200) {
        ld m1 = lo + (hi - lo) / 3, m2 = hi - (hi - lo) / 3;
        if (f(m1) < f(m2)) hi = m2; else lo = m1;
    }
    return (lo + hi) / 2;
}
\end{lstlisting}

\entry{Coordinate compression}{$O(n \log n)$}{%
Maps arbitrary values onto \code{0 .. k-1} keeping order, so a \code{BIT} or \code{Seg}
can index them. This is what lets you avoid \code{SparseSeg} whenever the whole input is
readable up front --- which is almost always.\\[0.3ex]
Remember to compress \emph{every} value you will ever look up, queries included, not
just the ones in the array. For half-open interval problems add both \code{l} and
\code{r} (and sometimes \code{r+1}).}

\begin{lstlisting}
// returns the sorted distinct values; id(v) is v's index
vi compress(vi vals) {
    sort(all(vals));
    vals.erase(unique(all(vals)), vals.end());
    return vals;
}
int id(const vi& c, ll v) {
    return (int)(lower_bound(all(c), v) - c.begin());
}
\end{lstlisting}

\entry{Interval cover}{$O(n \log n)$}{%
Fewest intervals from \code{I} that together cover the target \code{G}; returns their
indices, or empty if impossible. Intervals are \textbf{half-open} \code{[a, b)}; for
closed ones on integers use \code{[a, b+1)}. Greedy: from the current point, among the
intervals starting at or before it, take the one that reaches furthest. Same idea for
``fewest jumps'' and ``fewest refuelling stops'' problems. From KACTL.}

%LABEL cover
\begin{lstlisting}
template<class T>
vector<int> cover(pair<T,T> G, vector<pair<T,T>> I) {
    vector<int> S(I.size()), R;
    iota(all(S), 0);
    sort(all(S), [&](int a, int b) { return I[a] < I[b]; });
    T cur = G.first;
    int at = 0;
    while (cur < G.second) {
        pair<T, int> mx = make_pair(cur, -1);
        while (at < (int)I.size() && I[S[at]].first <= cur) {
            mx = max(mx, make_pair(I[S[at]].second, S[at]));
            at++;
        }
        if (mx.second == -1) return {};
        cur = mx.first;
        R.push_back(mx.second);
    }
    return R;
}
\end{lstlisting}

\entry{Fast input}{$\sim$5x faster than \code{cin}}{%
For inputs of $10^7$ numbers, where even \code{sync\_with\_stdio(false)} is too slow.
Reads from \code{stdin} in blocks; do not mix with \code{cin}. \code{readInt} skips
spaces and newlines and handles a minus sign; for \code{ll} change the types. From
KACTL.}

%LABEL fastin
\begin{lstlisting}
inline char gc() {                 // like getchar()
    static char buf[1 << 16];
    static size_t bc, be;
    if (bc >= be) {
        buf[0] = 0, bc = 0;
        be = fread(buf, 1, sizeof(buf), stdin);
    }
    return buf[bc++];              // 0 on end of file
}
int readInt() {
    int a, c;
    while ((a = gc()) < 40);
    if (a == '-') return -readInt();
    while ((c = gc()) >= 48) a = a * 10 + c - 480;
    return a - 48;
}
\end{lstlisting}

\section{Reference}

\entry{What fits in the time limit}{$\sim10^8$ simple ops/s at \code{-O2}}{%
{\ttfamily\footnotesize
\begin{ctab}{0.20}{0.75}
n <= 10 & O(n!) \hfill brute force permutations\\
n <= 20 & O(2\^{}n n) \hfill bitmask DP, PathDP\\
n <= 40 & O(2\^{}(n/2)) \hfill meet in the middle\\
n <= 100 & O(n\^{}4) \hfill 4 nested loops, Hungarian\\
n <= 500 & O(n\^{}3) \hfill Floyd, matrix, interval DP\\
n <= 5000 & O(n\^{}2) \hfill pairwise DP\\
n <= 1e5 & O(n log\^{}2 n) \hfill HLD + Seg\\
n <= 1e6 & O(n log n) \hfill sort, Seg, Dijkstra\\
n <= 1e7 & O(n) \hfill sieve, two pointers\\
n <= 1e18 & O(log n) \hfill binary search, matrix power\\
\end{ctab}}\\[0.3ex]
A \code{map} / \code{set} operation costs $\sim$10x a \code{vector} one; sorting
$10^6$ \code{ll} is $\sim$0.1 s; $10^8$ additions is $\sim$0.1 s but $10^8$ random
memory accesses is $\sim$1 s. If $n \le 18$ and nothing else fits, it is a bitmask.}

\entry{Overflow and precision}{read this when a big test fails}{%
\code{ll} holds $9.2\cdot10^{18}$, \code{int} only $2.1\cdot10^9$ --- and
\code{n*(n-1)/2} already overflows \code{int} at \code{n = 1e5}.\\[0.3ex]
Use \code{INF = 1e18}, never \code{LLONG\_MAX}: \code{INF + INF} still fits, so a
forgotten guard gives a wrong answer instead of undefined behaviour. Guard anyway.\\[0.3ex]
\code{a * b \% m} overflows once $m > 3\cdot10^9$ --- cast to \code{i128}.
\code{cross} on coordinates up to $10^9$ reaches $2\cdot10^{18}$: fine, barely.\\[0.3ex]
\code{-7 \% 3 == -1} in C++. Normalise with \code{((x \% m) + m) \% m}.\\[0.3ex]
\code{pow(10, 2)} can return \code{99.999...}; never use \code{pow} for integers.
\code{sqrt(x)} likewise --- take \code{(ll)sqrtl(x)} then fix up by $\pm1$ in a loop.\\[0.3ex]
Compare squared distances, never \code{sqrt}ed ones. \code{double} has 15--16 significant
digits, \code{long double} 18--19 on x86.}

\entry{Pitfalls}{the usual wrong answers}{%
\textbf{Multiple test cases:} clear every global, resize every vector, reset
\code{cnt}/\code{timer}. The most common repeat-offender bug there is.\\[0.3ex]
\textbf{Recursion depth:} a DFS on a path of $2\cdot10^5$ vertices can blow the stack.
The iterative \code{euler}, \code{SCC}, \code{Bridges}, \code{centroids},
\code{eulerPath}, \code{Dinic}, \code{MCMF}, \code{Matching}, \code{TwoSat} and
\code{centroidTree} here are safe; \code{HLD} is not.\\[0.3ex]
\textbf{Output:} \code{'\textbackslash n'} not \code{endl} (\code{endl} flushes).
Keep \code{sync\_with\_stdio(false)}, and then never mix \code{printf} with
\code{cout}.\\[0.3ex]
\textbf{\code{unordered\_map}} is $O(n)$ per operation against a crafted test; use
\code{map}, or add a random offset to every key.\\[0.3ex]
\textbf{Comparators} must be a strict weak ordering --- \code{<=} in a \code{sort}
comparator is a crash, not a wrong answer.\\[0.3ex]
\textbf{\code{vector<bool>}} is a proxy type; \code{auto\& x = v[i]} does not do what you
expect. \textbf{Iterators} invalidate on \code{push\_back}.\\[0.3ex]
\textbf{Reading:} after \code{cin >> n}, a following \code{getline} reads the rest of that
line. Interactive problems need an explicit flush.\\[0.3ex]
\textbf{Edge cases:} $n = 1$, $n = 0$, an empty string, all elements equal, a
disconnected graph, self-loops, multi-edges, duplicate points, collinear points.}

\entry{Combinatorics}{}{%
$\binom{n}{k} = \binom{n-1}{k-1} + \binom{n-1}{k}$; hockey stick
$\sum_{i=r}^{n}\binom{i}{r} = \binom{n+1}{r+1}$; Vandermonde
$\sum_k \binom{m}{k}\binom{n}{p-k} = \binom{m+n}{p}$.\\[0.3ex]
\textbf{Stars and bars.} $n$ identical items into $k$ ordered boxes:
$\binom{n+k-1}{k-1}$, and $\binom{n-1}{k-1}$ if every box must be non-empty.\\[0.3ex]
\textbf{Catalan} $C_n = \frac{1}{n+1}\binom{2n}{n}$: balanced bracket strings,
binary trees on $n$ nodes, triangulations of an $(n{+}2)$-gon, paths under the
diagonal. $1, 1, 2, 5, 14, 42, 132, 429$.\\[0.3ex]
\textbf{Derangements} $D_n = (n-1)(D_{n-1} + D_{n-2})$, $D_0 = 1$, $D_1 = 0$.\\[0.3ex]
\textbf{Inclusion-exclusion}
$|\bigcup A_i| = \sum|A_i| - \sum|A_i \cap A_j| + \dots$ --- over $\le 20$ sets this is
a mask loop with $(-1)^{\mathrm{popcount}}$.\\[0.3ex]
\textbf{Burnside.} Orbits = average number of fixed points over the group. Necklaces of
$n$ beads in $k$ colours under rotation: $\frac{1}{n}\sum_{i=0}^{n-1} k^{\gcd(i,n)}$.\\[0.3ex]
\textbf{Lucas.} $\binom{n}{k} \bmod p$ for prime $p$ = product of
$\binom{n_i}{k_i}$ over the base-$p$ digits (0 if any $k_i > n_i$).\\[0.3ex]
$\sum_{i\le n} i = \frac{n(n+1)}{2}$, $\sum i^2 = \frac{n(n+1)(2n+1)}{6}$,
$\sum i^3 = \left(\frac{n(n+1)}{2}\right)^2$.}

\entry{Number theory}{}{%
$\gcd(a,b)\cdot\mathrm{lcm}(a,b) = ab$ --- compute \code{a / gcd * b} to avoid
overflow.\\[0.3ex]
$\varphi(n) = n\prod_{p\mid n}(1-\frac1p)$; $\sum_{d\mid n}\varphi(d) = n$.
Fermat: $a^{p-1} \equiv 1$ for prime $p \nmid a$. Euler:
$a^{\varphi(m)} \equiv 1$ when $\gcd(a,m)=1$.\\[0.3ex]
$d(n) = \prod(e_i+1)$ divisors, $\sigma(n) = \prod\frac{p_i^{e_i+1}-1}{p_i-1}$.
The number of divisors stays small: at most 240 below $10^9$, 1344 below
$10^{18}$.\\[0.3ex]
\textbf{Legendre.} The exponent of $p$ in $n!$ is
$\sum_{i\ge1}\lfloor n/p^i \rfloor$ --- that is how you count trailing zeros of a
factorial.\\[0.3ex]
\textbf{M\"obius} $\mu(n)$ is $0$ with a squared factor, else $(-1)^{\#\text{primes}}$;
$\sum_{d\mid n}\mu(d) = [n=1]$, which turns ``count coprime pairs'' into a divisor
sum.\\[0.3ex]
Primes below $10^k$: $4$, $25$, $168$, $1229$, $9592$, $78498$, $664579$, $5761455$.
The gap between consecutive primes near $10^{18}$ stays under 1500.}

\entry{Graph facts}{}{%
\textbf{Max flow = min cut.} \textbf{K\"onig} (bipartite): max matching = min vertex
cover, and max independent set $= n - $ max matching. \textbf{Hall:} a perfect matching
on the left exists iff every subset $S$ of the left has
$|N(S)| \ge |S|$.\\[0.3ex]
\textbf{Dilworth:} the minimum number of chains covering a poset equals the largest
antichain --- in practice a bipartite matching or an LIS.\\[0.3ex]
Bipartite iff no odd cycle. A tree has $n-1$ edges and a unique path between any two
vertices. A connected graph has an Euler circuit iff every degree is even (directed: in
= out everywhere).\\[0.3ex]
\textbf{Planar:} $V - E + F = 2$, so $E \le 3V - 6$ and some vertex has degree
$\le 5$. Four colours always suffice, five are easy to find.\\[0.3ex]
\textbf{Spanning trees:} Cayley gives $n^{n-2}$ labelled trees on $n$ vertices; for a
general graph it is any cofactor of the Laplacian (Matrix-Tree), i.e. one
\code{gauss} determinant.\\[0.3ex]
\textbf{Pick:} a lattice polygon has $A = I + B/2 - 1$, with
$B = \sum \gcd(|dx|,|dy|)$ over the edges.}

\entry{Games (Sprague-Grundy)}{}{%
Two players alternate, the one who cannot move loses, both play perfectly, nothing is
hidden and there is no luck. Every position gets a number, its \textbf{Grundy value}:
$g(x) = \mathrm{mex}\{g(y) : y \text{ reachable from } x\}$, where mex is the smallest
non-negative integer \emph{not} in the set. A position with no moves has $g = 0$. The
player to move \textbf{loses iff $g = 0$}.\\[0.3ex]
\textbf{Several independent games at once} (move in exactly one of them per turn): the
Grundy value of the whole thing is the \textbf{XOR} of the parts. Nim: piles of sizes
$a_i$, take any amount from one pile; $g = a_1 \oplus a_2 \oplus \dots$. A move that
splits one game into two gives the XOR of the two halves as one option.\\[0.3ex]
In practice: compute $g$ for small sizes by brute force with memo, print them, look for
the pattern (periodic, $n \bmod k$, \dots), then prove or trust it. Mis\`ere Nim (last
to move loses): same as Nim unless every pile is $\le 1$; then the player to move wins
iff the number of piles is even.
Staircase Nim: only piles on odd steps matter. ``Take $1..k$ from one pile'':
$g = a \bmod (k+1)$.}

\entry{More counting}{}{%
\textbf{Stirling 2nd kind} $S(n,k)$: ways to split $n$ labelled items into $k$
non-empty unlabelled groups. $S(n,k) = k\,S(n{-}1,k) + S(n{-}1,k{-}1)$. Onto functions
$n \to k$: $k!\,S(n,k)$.\\[0.3ex]
\textbf{Bell} $B_n = \sum_k S(n,k)$: all ways to split $n$ items into groups.
$1, 1, 2, 5, 15, 52, 203, 877, 4140$.\\[0.3ex]
\textbf{Stirling 1st kind} $c(n,k)$: permutations of $n$ with exactly $k$ cycles.
$c(n,k) = (n{-}1)\,c(n{-}1,k) + c(n{-}1,k{-}1)$.\\[0.3ex]
\textbf{Partitions} $p(n)$ of $n$ into a sum of positive integers (order does not
matter): dp \code{p[0] = 1; for each part k: for i >= k: p[i] += p[i-k]}. $p(100) \approx
1.9\cdot10^8$.\\[0.3ex]
\textbf{Multinomial} $\binom{n}{k_1, \dots, k_m} = \frac{n!}{k_1! \cdots k_m!}$:
arrangements of a word with repeated letters.\\[0.3ex]
\textbf{Expected value} is linear even for dependent events: expected number of X =
sum over all candidates of P(candidate is X). Expected tries until a success of
probability $p$: $1/p$. Collect all $n$ coupons: $n H_n \approx n \ln n$.\\[0.3ex]
\textbf{Grid paths} from $(0,0)$ to $(a,b)$: $\binom{a+b}{a}$. Those that never touch
the line $y = x + 1$ ($a \ge b$, reflection trick): $\binom{a+b}{a} - \binom{a+b}{a+1}$.}

\entry{Geometry formulas}{}{%
Triangle with sides $a, b, c$, $s = (a+b+c)/2$: area
$\sqrt{s(s-a)(s-b)(s-c)}$ (Heron); $c^2 = a^2 + b^2 - 2ab\cos\gamma$;
$\frac{a}{\sin\alpha} = 2R$; circumradius $R = \frac{abc}{4A}$; inradius $r = A / s$.\\[0.3ex]
Angle between vectors $u, v$: \code{atan2(cross, dot)} --- signed and stable, better than
\code{acos}. Distance from point $p$ to line $ab$: $|(b-a)\times(p-a)| / |b-a|$.\\[0.3ex]
Circle sector area $\frac12 r^2\theta$, segment area $\frac12 r^2(\theta - \sin\theta)$,
chord length $2r\sin(\theta/2)$.\\[0.3ex]
Sphere: $V = \frac43\pi r^3$, $A = 4\pi r^2$; cap of height $h$: $V = \frac{\pi h^2
(3r - h)}{3}$, $A = 2\pi r h$. Cone / pyramid: $V = \frac13 \cdot$ base $\cdot$ height.\\[0.3ex]
Great-circle distance, latitude $\varphi$, longitude $\lambda$:
$d = r \arccos(\sin\varphi_1\sin\varphi_2 + \cos\varphi_1\cos\varphi_2\cos\Delta\lambda)$.\\[0.3ex]
Polygon on lattice points: Pick, see \code{Polygons}. Volume of a polyhedron with
outward faces: $\frac16\sum_{\text{faces}} (a \times b)\cdot c$ over triangulated
faces.}

\entry{Technique by symptom}{}{%
\begin{ctab}{0.43}{0.52}
\textbf{You see} & \textbf{Try}\\
answer is ``smallest $x$ such that \dots works'' & binary search on $x$\\
min of a max / max of a min & binary search, or sort + DSU\\
queries offline, any order & sort queries; Mo; DFS + rollback\\
``for every pair'' and $n \le 10^5$ & sort, two pointers, sweep, FFT\\
sum over all subarrays & contribution of each element\\
subsets and $n \le 20$ & bitmask DP, SOS\\
$n \le 40$ & meet in the middle\\
shortest path, few weights & 0-1 BFS (deque), Dijkstra\\
``number of ways'' mod $p$ & DP, combinatorics, Berlekamp-Massey\\
$k$ up to $10^{18}$ steps & matrix power, cycle finding, doubling\\
grid, ``each cell one of two'' & bipartite / flow / 2-SAT\\
intervals & sort by end (greedy), sweep with events\\
tree, all paths & centroid decomposition, rerooting\\
tree, path queries & HLD, LCA + prefix sums\\
DAG / game & DP in topological order, Grundy\\
strings, many patterns & Aho-Corasick, suffix array, hashing\\
XOR of subsets & XOR basis, trie\\
points, ``nearest / enclosing'' & hull, closest pair, \code{mec}\\
exact formula unknown & brute force small cases, find the pattern\\
\end{ctab}}

\entry{When stuck}{}{%
Re-read the statement and the constraints. The bound on $n$ usually names the intended
complexity outright --- check the table above before designing anything.\\[0.3ex]
Write the brute force first, then stress test against it. A wrong answer on $n=3$ is
worth more than an hour of staring.\\[0.3ex]
Reframe: is it a graph? a flow? a matching? Can you binary-search the answer? Sort the
input? Process queries offline, sorted? Think backwards from the end? Count the
complement? Fix one variable and see what is left?\\[0.3ex]
Check the samples \emph{by hand}. If the sample passes but the judge does not, the bug is
almost always an edge case from the Pitfalls list --- most often multiple test cases, or
an overflow.\\[0.3ex]
\footnotesize
\emph{Team Wery Accepted: Ruben, Ruiming, Teo. Sources: templates by Ruben, Ruiming,
b5xu and Martin (earlier teammate), consolidated. Entries marked ``From KACTL'' and
\code{LineContainer} are adapted from KACTL (KTH, CC0). Provenance and the per-entry
diff against the original files are in \code{notebook/README.md}.}}


\end{multicols*}
\end{document}

